\documentclass[11pt,a4paper]{article}
\usepackage[margin=27mm,headheight=15pt]{geometry}
\usepackage{fontspec}
\setmainfont{Latin Modern Roman}
\setsansfont{Latin Modern Sans}
\setmonofont{Latin Modern Mono}
\usepackage{amsmath,amssymb,amsthm,mathtools}
\usepackage{booktabs,longtable,array,tabularx,enumitem,microtype,xcolor}
\usepackage[authoryear,round]{natbib}
\usepackage[hidelinks,unicode]{hyperref}
\hypersetup{pdftitle={CATHOT: What categories modulo natural transformations can represent},pdfauthor={AI-produced by GPT-6 Astra (OpenAI, ChatGPT Pro mode)}}
\usepackage{fancyhdr}
\pagestyle{fancy}\fancyhf{}\fancyhead[L]{\small CATHOT}\fancyhead[R]{\small Categories, transformations, and homotopy types}\fancyfoot[C]{\thepage}
\setlength{\parindent}{1.3em}\setlength{\parskip}{3pt}
\setlength{\emergencystretch}{3em}
\setcounter{secnumdepth}{0}
\newcolumntype{L}[1]{>{\raggedright\arraybackslash}p{#1}}
\newcommand{\Cat}{\mathrm{Cat}}
\newcommand{\Cathot}{{\mathrm{Cat}_{\mathrm{hot}}}}
\newcommand{\Hot}{\mathrm{Hot}}
\newcommand{\Ch}{\mathrm{Ch}}
\newcommand{\Set}{\mathrm{Set}}
\newcommand{\Spc}{\mathcal S}
\newcommand{\U}{\mathcal U}\newcommand{\V}{\mathcal V}
\newcommand{\W}{\mathcal W}
\newcommand{\Z}{\mathbb Z}\newcommand{\Q}{\mathbb Q}\newcommand{\F}{\mathbb F}\newcommand{\N}{\mathbb N}
\newcommand{\op}{\mathrm{op}}\newcommand{\id}{\mathrm{id}}
\DeclareMathOperator{\Hom}{Hom}\DeclareMathOperator{\Fun}{Fun}
\DeclareMathOperator{\Ob}{Ob}\DeclareMathOperator{\Mor}{Mor}
\DeclareMathOperator{\Map}{Map}\DeclareMathOperator{\colim}{colim}
\DeclareMathOperator{\Ind}{Ind}\DeclareMathOperator{\Aut}{Aut}
\DeclareMathOperator{\End}{End}
\DeclareMathOperator{\hocolim}{hocolim}\DeclareMathOperator{\holim}{holim}
\DeclareMathOperator{\Ho}{Ho}\DeclareMathOperator{\sk}{sk}
\newcommand{\msp}[2]{\textup{#1 ms.\ #2}}
\makeatletter
\newenvironment{claim}[3]{\par\medskip\noindent\def\@currentlabel{#1}\phantomsection\label{stmt:#1}\textbf{#1. #3}\quad{\sffamily\small [#2].}\enspace\ignorespaces}{\par\medskip}
\makeatother
\newcommand{\partsection}[2]{\section*{#1. #2}\addcontentsline{toc}{section}{#1. #2}}
\title{CATHOT:\\What categories modulo natural transformations can represent}
\author{AI-produced by GPT-6 Astra\\OpenAI, ChatGPT Pro mode}
\date{10 November 2026}
\begin{document}
\begin{center}
{\Large\bfseries CATHOT: What categories modulo\\natural transformations can represent\par}
\vspace{6pt}
{\normalsize AI-produced by GPT-6 Astra (OpenAI, ChatGPT Pro mode)\par}
{\small 10 November 2026\par}
\end{center}
\noindent\small\textit{Research audit of Grothendieck's \emph{Les D\'erivateurs}, Chapters V, VII, IX, and XII appendix \citep{Grothendieck1991}. No priority is claimed. Numbered assertions carry their proof status.}
\begingroup\small
\section*{Conventions and scope}
We fix Grothendieck universes $\U\in\V$ with $\N\in\U$; a category called small has its sets of objects and arrows in $\U$. Categories of all $\U$-small objects and their presheaves are formed in $\V$. For a small category $X$, $NX$ is its nerve and $BX=|NX|$ its classifying space. All homotopy classes are unpointed unless a subscript $*$ occurs. We put $\Hom_{\Cathot}(X,K)=\pi_0\Fun(X,K)$, using finite zigzags of natural transformations. We write $\W_\infty$ for functors whose nerves are weak homotopy equivalences and $\Hot=\Cat[\W_\infty^{-1}]$. Equality, categorical isomorphism, categorical equivalence, isomorphism in $\Cathot$, and isomorphism in $\Hot$ are distinct notions.

For $u:A\to B$ and $b\in B$, $A/b$ has objects $(a,u(a)\to b)$ and $b\backslash A$ has objects $(a,b\to u(a))$. A functor is final in the covariant convention when each $b\backslash A$ is nonempty and connected. Presheaf coefficients have variance $X^{\op}\to D(e)$. A sieve is full and closed under sources of arrows, and a cosieve is full and closed under targets. We translate the manuscript's terminology according to this convention at each use.

A fundamental localizer $\W$ is used with the following precise version of Loc(1)--(3). Loc(1): $\W$ contains identities, has two-out-of-three, and if $i:A\to B$, $r:B\to A$ satisfy $ri=\id_A$ and $ir\in\W$, then $i,r\in\W$. Loc(2): if $A$ has a terminal object then $A\to e$ belongs to $\W$. Loc(3), relative Theorem A: for a strictly commutative triangle $A\xrightarrow{u}B\xrightarrow{v}C$, if $u/c:A/c\to B/c$ belongs to $\W$ for every $c\in C$, then $u\in\W$. Here the slice structures are those induced by $vu$ and $v$. We call $A$ $\W$-aspherical when $A\to e$ belongs to $\W$. Our homotopy-theoretic conclusions use $\W_\infty$; formal localization statements explicitly indicate when another class is allowed.

A finite category has finitely many arrows, including identities. A finite homotopy type, when this expression is used in the compact-generation argument, means a homotopy type represented by a finite CW complex; it does not mean a space with finitely many homotopy groups or a nerve of a finite category. Empty categories and empty spaces are permitted.

\endgroup\normalsize
\clearpage
\tableofcontents
\clearpage
\partsection{A}{The localization and the formal arguments}

\begin{claim}{A.1}{Proved}{The category and its localization}
There is a category $\Cathot_\U$ with the same objects as $\Cat_\U$ and
\[
 \Hom_{\Cathot}(X,K)=\pi _0\Fun(X,K).
\]
Here two functors represent the same morphism precisely when they are joined by a finite zigzag of natural transformations, whose orientations may vary. The class $\overline\W_\infty$ consists of the classes of functors in $\W_\infty$. The functor
$\varphi:\Cathot\to\Hot$, $[f]\mapsto[Bf]$, is a localization at this class, and
$\varphi[f]$ is invertible precisely when $f\in\W_\infty$. The identification
$\Hot\simeq\Ho(\Spc_\U)$ uses weak homotopy types of spaces, including the empty space.
\end{claim}
\begin{proof}
Whiskering a natural transformation by functors preserves its source, target, and orientation. It therefore preserves finite zigzags and defines associative composition on the displayed sets. A natural transformation $f\Rightarrow g$ is a functor $X\times[1]\to K$ restricting to $f$ and $g$. Its nerve gives a homotopy $Bf\simeq Bg$. Thus $\W_\infty$ is constant on these classes.

More intrinsically, the projection $p:X\times[1]\to X$ is in $\W_\infty$ and its two sections $i_0,i_1:X\to X\times[1]$ become the same inverse of $p$ in every category receiving a functor that inverts $\W_\infty$. Such a functor therefore identifies $f=Hi_0$ and $g=Hi_1$. The universal property of localization now factors through $\Cathot$ and gives the claimed localization. The nerve induces the equivalence with the homotopy category of simplicial sets, for example by the category-of-simplices resolution; alternatively this follows from the Quillen equivalence in A.5. Weak equivalences of simplicial sets are saturated, so $Bf$ is an isomorphism of homotopy types if and only if it is a weak equivalence. These arguments also cover $X=\varnothing$; its functor category into any $K$ is terminal.
\end{proof}

\begin{claim}{A.2}{Proved}{The specified fundamental-localizer convention}
The class $\W_\infty$ satisfies Loc(1)--(3) as stated on the first page. For any fundamental localizer in that convention, localization identifies naturally transformed functors. In particular, the categorical quotient used in A.1 is compatible with the variance of those axioms.
\end{claim}
\begin{proof}
Two-out-of-three and the stated retraction axiom hold for weak homotopy equivalences. A category with a terminal object has a natural transformation from its identity to the constant functor at that object, so its nerve is contractible. For the relative assertion, the covariant diagrams $c\mapsto A/c$ and $c\mapsto B/c$ on $C$ have Grothendieck constructions whose nerves are weakly equivalent to $NA$ and $NB$, respectively. To see the first comparison explicitly, the total category has objects $(a,c,vu(a)\to c)$. Its projection to $A$ has the section $a\mapsto(a,vu(a),\id)$, and the defining arrow supplies a natural transformation from section--projection to the identity. The construction for $B$ is identical. Objectwise weak equivalences $A/c\to B/c$ give a weak equivalence of these total categories by the homotopy-colimit theorem \citep[Theorem~26.8 and Corollary~26.9]{ChacholskiScherer2001}. This proves the relative Theorem A axiom; compare also \citet[\S1, Theorem~A]{Quillen1973}.

For an arbitrary $\W$ satisfying the stated axioms, $X/x$ and $(X\times[1])/x=(X/x)\times[1]$ have terminal objects. Their projections to the point belong to $\W$, and two-out-of-three shows that the projection of these slices belongs to $\W$. Loc(3) gives $X\times[1]\to X$ in $\W$. The section argument in A.1 then identifies naturally transformed functors in $\Cat[\W^{-1}]$. We have used the slices $A/c$, not the coslices $c\backslash A$; replacing one by the other without reversing the convention would change the hypothesis.
\end{proof}

\begin{claim}{A.3}{Proved}{Proposition 1 is a conditional formal theorem}
Let $L:C\to D$ and $R:D\to C$ be functors between locally small categories, let $L\dashv R$, and suppose $R$ is fully faithful. Write $\eta:\id_C\to RL$ and $\epsilon:LR\to\id_D$ for the unit and counit, and $W_L=\{u\mid Lu\text{ is invertible}\}$. For $K\in C$, the following conditions are equivalent: $\Hom_C(-,K)$ inverts $W_L$; it inverts every $\eta_X$; precomposition with $\eta_K$ is bijective; $\eta_K$ has a retraction; $\eta_K$ is invertible; $K$ is in the essential image of $R$; and the canonical map
\[
 \Hom_C(X,K)\longrightarrow\Hom_D(LX,LK)
\]
is bijective for every $X\in C$. If $H':D^{\op}\to\Set$ and $H=H' L^{\op}$, then $H$ is representable if and only if $H'$ is representable. In that event their representing objects satisfy $K\cong RK'$ and $K'\cong LK$, compatibly with the specified representations.
\end{claim}
\begin{proof}
Full faithfulness of $R$ makes $\epsilon$ invertible, and the triangle identity makes $L\eta_X$ invertible. The first three conditions successively imply the next, since lifting $\id_K$ under precomposition gives $p:RLK\to K$ with $p\eta_K=\id_K$. Naturality of $\eta$ gives
\[
 \eta_Kp=RL(p)\eta_{RLK}.
\]
Applying $L$ to $p\eta_K=\id$ shows that $Lp=(L\eta_K)^{-1}$. The other triangle identity and invertibility of $\epsilon$ identify $\eta_{RLK}$ with $RL\eta_K$. The right side is therefore the identity. Thus $p$ is the inverse of $\eta_K$.

If $K=RD$ up to isomorphism, adjunction and the counit identify the canonical comparison with a bijection. Conversely, that comparison, followed by adjunction, is precisely $\Hom_C(X,\eta_K)$; ordinary Yoneda makes $\eta_K$ invertible. Finally the comparison makes $\Hom_C(-,K)$ invert $W_L$. This closes the equivalences, including the apparently weaker retraction condition. If $H'$ is represented by $K'$, adjunction represents $H$ by $RK'$. If $H$ is represented by $K$, it inverts $W_L$, so the equivalence just proved identifies it with $\Hom_D(L-,LK)$. Restriction along a localization is fully faithful on presheaves; in this reflective case $L$ is such a localization. Hence $H'$ is represented by $LK$. Every uniqueness asserted here is the uniqueness supplied by Yoneda relative to the given natural isomorphism.
\end{proof}

\begin{claim}{A.4}{Proved}{Ind-categories and Proposition 2}
Use one fixed auxiliary universe for the filtered indexing categories. For $I,J$ small filtered, $K:I\to C$ and $T:J\to C$, put
\[
 h_K(X)=\colim_{i\in I}\Hom_C(X,K_i),\qquad
 \Hom_{\Ind C}(K,T)=\lim_{i\in I}\colim_{j\in J}\Hom_C(K_i,T_j).
\]
The association $K\mapsto h_K$ is fully faithful. Under the hypotheses of A.3, $\Ind L\dashv\Ind R$ and $\Ind R$ is fully faithful. For an ind-object $K$, the following are equivalent: $h_K$ inverts $W_L$; it inverts all $\eta_X$ for $X\in C$; the unit $K\to(\Ind R)(\Ind L)K$ is invertible; $K$ belongs to the essential image of $\Ind R$; and the canonical comparison
\[
 h_K(X)\longrightarrow
 \Hom_{\Ind D}(LX,(\Ind L)K)
\]
is bijective for every $X\in C$. These are precisely the valid, typed contents of V ms.~40--42.
\end{claim}
\begin{proof}
Yoneda and the universal property of the indicated colimit give
\[
 \operatorname{Nat}(h_K,h_T)
 =\lim_i\operatorname{Nat}(h_{K_i},h_T)
 =\lim_i\colim_j\Hom_C(K_i,T_j).
\]
No interchange of a limit and a colimit has occurred. Applying the adjunction inside this formula gives the ind-adjunction; applying full faithfulness inside it proves full faithfulness of $\Ind R$.

We verify the step on which the manuscript comments. Extend $h_K$ to ind-objects by $\overline h_K(X)=\Hom_{\Ind C}(X,K)=\lim_jh_K(X_j)$. If $h_K(\eta_A)$ is bijective for every $A\in C$, then $\overline h_K$ inverts the levelwise unit of every ind-object, since a limit of componentwise bijections is a bijection. Apply A.3 in $\Ind C$ to obtain the invertible unit of $K$. All remaining implications follow from A.3. Conversely, the comparison on constant source objects is exactly a natural isomorphism of the associated presheaves, and full faithfulness just proved detects the unit isomorphism. This argument uses $R$ where stated; it supplies no $R$ for $\varphi$.
\end{proof}

\begin{claim}{A.5}{Known}{The precise model-category mapping formula}
Thomason's model category has weak equivalences $\W_\infty$ and fibrations detected by $\operatorname{Ex}^2N$. Its adjunction
\[
 c\operatorname{Sd}^2:\mathbf{sSet}\rightleftarrows\Cat:
 \operatorname{Ex}^2N
\]
is a Quillen equivalence. Choose a fibrant replacement $K\to RK$ in this model category. Put $QX=c\operatorname{Sd}^2NX$. The composite last-vertex map gives a natural weak equivalence $q_X:QX\to X$, and $QX$ is cofibrant. There are natural bijections
\[
 [BX,BK]\cong\pi_0\Map_{\mathbf{sSet}}(NX,\operatorname{Ex}^2NRK)
 \cong\Fun(QX,RK)_0/\sim_T .
\]
Here $\Fun(QX,RK)_0$ is the set of functors, and $\sim_T$ is generated by endpoint restrictions of functors
$c\operatorname{Sd}^2(NX\times\Delta^1)\to RK$.
For any cofibrant replacement $Q'X\to X$, one may equivalently use $\Hom_{\Cat}(Q'X,RK)$ modulo model-category homotopy. The displayed explicit relation specifies that homotopy; it does not silently replace it by a different cylinder relation.
\end{claim}
\begin{proof}
The model structure and properness are \citet[Theorem~4.9 and Corollary~5.5]{Thomason1980}; the adjoint comparison is given in the introduction and \S5, and is recalled with the corrected cofibration argument by \citet[\S2]{Raptis2010}. Both subdivision maps and the adjunction comparison are weak equivalences, which proves the assertion about $q_X$. Since every simplicial set is cofibrant, the left Quillen functor takes $NX$ to a cofibrant category. The precise cosimplicial resolution $c\operatorname{Sd}^2(NX\times\Delta^\bullet)$ is also recorded in \citet[Remark~4.9]{Raptis2010}. The precise cosimplicial resolution $c\operatorname{Sd}^2(NX\times\Delta^\bullet)$ is also recorded in \citet[Remark~4.9]{Raptis2010}. The right-hand simplicial target is Kan by the definition of Thomason fibrancy. Its function complex therefore computes the derived mapping space out of $NX$. By adjunction, its vertices and edges are exactly the functors specified in the assertion. Connected components are the quotient by those edges. This proves the formula with its exact comparison and exact equivalence relation. The general cofibrant--fibrant formula is the defining computation of the homotopy category of a model category. In particular, calling $RK$ fibrant is not a proof of a formula with the unreplaced source $X$.
\end{proof}

\begin{claim}{A.6}{Refuted}{A fibrant target does not suffice}
The assertion that $\pi_0\Fun(X,P)\to[BX,BP]$ is always a bijection for a Thomason-fibrant $P$ is false, even for a finite $X$. The assertion that $\Fun(X,P)$ is always Thomason-fibrant is also false, even for a finite $X$ and a countable fibrant $P$.
\end{claim}
\begin{proof}
For the first assertion, let $X=B C_2$ be the one-object, two-arrow category of the cyclic group of order two. By \citet[Theorem~2.6]{Raptis2010}, choose a fibrant replacement within posets of a poset having classifying space $BC_2$. Denote it by $P$. A functor $BC_2\to P$ has to send every group element to the identity of its chosen object. Natural transformations are precisely the order relations in $P$. Thus $\Fun(BC_2,P)\cong P$. Its component set is a singleton, whereas $[BC_2,BP]$ contains a constant map and a homotopy equivalence. They induce respectively the trivial map and an isomorphism on fundamental groups, so cannot be unpointedly homotopic. The source is not cofibrant: every cofibrant category is a poset \citep[Proposition~5.7]{Thomason1980}.

For the second assertion we verify the separate construction in the companion. Let $D$ be the four-element crown, with minimal elements $a_0,a_1$, maximal elements $b_0,b_1$, and exactly the four strict relations $a_i<b_j$. Its nerve is a circle. It is not fibrant: $\operatorname{Ex}^2ND$ has finitely many edges, but a Kan complex represents each element of its based fundamental group by a single loop edge. Its fundamental group would consequently be finite, contradicting $\pi_1BD=\Z$.

Let $G=C_2$ act trivially on $D$. Starting with $Q_0=D$, at stage $r$ attach, for every Thomason generating horn $A\hookrightarrow B$ and every $F:A\to Q_r$, the free $G$-cell $G\times B$ along $(g,a)\mapsto gF(a)$. These are pushouts of Dwyer inclusions of finite posets. The old poset remains full and every new object is in a free $G$-orbit; this is the explicit pushout description in \citet[Lemmas~4.7--4.8]{MayStephanZakharevich2017}. Hence $Q=\colim_rQ_r$ is countable, fibrant, and $Q^G=D$. Fibrancy follows because each horn map has finite domain and receives a filler at a subsequent stage.

Define $E=Q\rtimes G$: an arrow $x\to y$ labelled $g$ exists when $gx\leq y$, and the composite of labels $g,h$ is $hg$. For a horn diagram $A\to E$, its labels define $A\to BG$. The horn's classifying space is contractible, so choose labels $h_a$ with $h_bg_{ab}=h_a$. Replacing the object at $a$ by $h_a x_a$ gives a diagram $A\to Q$. Extend it to $B\to Q$, extend the choices $h_a$ to the new objects, and undo the change. This produces a strict extension to $E$, proving that $E$ is fibrant.

A functor $BG\to E$ has label homomorphism either $0$ or $\id_G$. In the latter case its object satisfies $gx\leq x$, hence $gx=x$, and so lies in $D$. Naturality of a transformation labelled $h$ requires $h\phi(g)=\psi(g)h$; because $G$ is abelian it separates the two cases. Consequently
\[
 \Fun(BG,E)\cong E\amalg(D\times BG).
\]
The category $D$ is a retract of this category: project the second summand to $D$, send the first constantly to $a_0$, and use identity labels for the section. Fibrant objects are closed under retracts. Since $D$ is not fibrant, this functor category is not fibrant. The finite crown and its nerve are checked in \texttt{code/a\_finite\_checks.py}; the transfinite construction is proved above, not certified by a finite calculation.
\end{proof}

\begin{claim}{A.7}{Proved}{Descent of ind-representability needs no right adjoint}
Let $L:C\to D$ be any localization, and let $H=\colim_{i\in I}h_{K_i}$ be a small filtered colimit of representable presheaves. If $H\cong H' L^{\op}$, then the canonical comparison identifies
\[
 H'\cong\colim_{i\in I}\Hom_D(-,LK_i).
\]
For a single representing object, it follows that $h_K$ inverts the localized class if and only if its canonical comparison with $\Hom_D(L-,LK)$ is bijective. This assertion, and not a reflection, is sufficient for the descent of the manuscript's external Hom ind-objects.
\end{claim}
\begin{proof}
Restriction $L^*$ on presheaves is fully faithful, by the universal property of a localization including natural transformations. For any presheaf $F:D^{\op}\to\Set$, put $T=\colim_i h_{LK_i}$ and compute
\[
 \operatorname{Nat}(T,F)=\lim_iF(LK_i)
 =\operatorname{Nat}(H,L^*F)
 =\operatorname{Nat}(H',F).
\]
Yoneda in the presheaf category identifies $T$ with $H'$. The comparison is the one induced by $L$: this follows by tracing the identity transformation under the displayed bijections. The single-object case gives the last assertion. Thus the reverse implication, from representability on $D$ to representability on $C$, is exactly the step that is absent without $R$; the forward descent is formal.
\end{proof}

\partsection{B}{Cardinal obstructions to representability}

We write $\mathsf B G$ for the category with one object and endomorphism
group $G$, with categorical composition given by group multiplication, and
$BG=B(\mathsf B G)$ for its classifying space.  All groups constructed below
are $\U$-small.  Our cardinal parameters are infinite cardinals in $\U$;
there is no passage to a group outside the universe of the source category.
For every $\U$-small set $S$, an infinite cardinal $\lambda$ in $\U$ with
$\lambda>|S|$ exists: one may take the cardinal of a sufficiently large
power set of $S\amalg\N$.  Free products, group presentations on such sets,
quotients, finite products, and countable unions remain in $\U$.

\begin{claim}{B.1}{Proved}{What a sufficiently large group sees in a small category}
Let $K$ be a $\U$-small category and let $G$ be a $\U$-small group such that
every group homomorphism $G\longrightarrow H$ is trivial whenever
$|H|\leq |\Mor K|$.  The constant-functor map and evaluation map
\[
 \delta_K:K\longrightarrow\Fun(\mathsf B G,K),\qquad
 \operatorname{ev}_*: \Fun(\mathsf B G,K)\longrightarrow K
\]
are mutually inverse isomorphisms of categories.  Consequently the map
\[
 p^*:\Hom_{\Cathot}(e,K)\longrightarrow
       \Hom_{\Cathot}(\mathsf B G,K),\qquad p:\mathsf B G\longrightarrow e,
\]
is a bijection.  Both assertions are natural in $K$ among categories
whose automorphism groups receive only trivial homomorphisms from $G$.
\end{claim}
\begin{proof}
A functor $F:\mathsf B G\to K$ chooses an object $k=F(*)$ and a
homomorphism $G\to\operatorname{Aut}_K(k)$.  Invertibility is preserved by
every functor, even when $K$ has noninvertible endomorphisms or nontrivial
idempotents.  Since
$|\operatorname{Aut}_K(k)|\leq |\Mor K|$, this homomorphism is trivial.
Thus $F=\delta_K(k)$ as an actual functor.  A natural transformation
$\delta_K(k)\Rightarrow\delta_K(l)$ consists of one arrow $k\to l$ of
$K$; its naturality equations are identities.  Evaluation therefore gives
a bijection on objects and on every hom-set, and preserves composition.
If $K$ is empty, both categories in the assertion are empty and the same
conclusion holds.  Passing to components proves the last assertion.
\end{proof}

\begin{claim}{B.2}{Known}{The finite Schur-multiplier input}
For a finite set $F$ with $|F|\geq8$,
\[
 H_1(\operatorname{Alt}(F);\Z)=0,
 \qquad H_2(\operatorname{Alt}(F);\Z)\cong\Z/2.
\]
If $a,b,c,d$ are distinct elements of $F$, the commuting permutations
$x=(ab)(cd)$ and $y=(ac)(bd)$ define a torus class
\[
 t_F=[x\mid y]-[y\mid x]\in H_2(\operatorname{Alt}(F);\Z)
\]
which generates this group.  In particular, inclusions between finite
alternating groups on sets containing these four elements and having at
least eight elements carry generator to generator.
\end{claim}
\begin{proof}[Source and the required transition-map check]
The multiplier calculation is Schur's classical finite calculation.  We
use the explicitly recorded formula and generator in
\citet[Eq.~(35), Proposition~20, and Corollary~21]{BC_DominguezSegovia2021};
the cited numbering is that of arXiv version~2, rather than the later
journal version.  This is the one imported finite group calculation in
the proof of Lead~1.  We do not infer a direct-limit calculation merely
from the abstract isomorphism $H_2\cong\Z/2$: the cycle displayed above
uses the same four permutations before and after any inclusion, and its
image is the displayed nonzero generator.  This verifies the required
transition maps.

There is also a direct check of its nontriviality.  In the real Clifford
algebra on orthonormal $e_a,e_b,e_c,e_d$, put
\[
 \widetilde x=\tfrac12(e_a-e_b)(e_c-e_d),\qquad
 \widetilde y=\tfrac12(e_a-e_c)(e_b-e_d).
\]
These are spin lifts of $x$ and $y$.  They satisfy
$\widetilde x^2=\widetilde y^2=-1$ and
$\widetilde x\widetilde y=-\widetilde y\widetilde x$.
Thus the central double cover evaluates nontrivially on the torus cycle.
The exact rational Clifford-algebra verification is included in
\texttt{code/bc\_checks.py}.  This check proves nontriviality; the assertion
that the whole finite multiplier has order two is the cited classical
input.
\end{proof}

\begin{claim}{B.3}{Proved}{Finitary alternating groups at every infinite cardinal}
Let $\Omega$ be an infinite $\U$-small set and let
$A_\Omega=\operatorname{Alt}_{\mathrm{fin}}(\Omega)$ be the group of even
permutations with finite support.  Then
\[
 |A_\Omega|=|\Omega|,\qquad
 A_\Omega\text{ is nonabelian simple},\qquad
 H_1(A_\Omega;\Z)=0,\quad H_2(A_\Omega;\Z)=\Z/2.
\]
Every homomorphism from $A_\Omega$ to a group of cardinality strictly
less than $|\Omega|$ is trivial.
\end{claim}
\begin{proof}
There are $|\Omega|$ finite subsets of $\Omega$, and each supports finitely
many permutations.  This gives the upper bound on the cardinality.  The
3-cycles using two fixed elements and a varying third element give the
opposite bound.

Let $N$ be a nontrivial normal subgroup and choose $g\in N\setminus\{1\}$.
Choose $a$ moved by $g$, put $b=g(a)$, and choose $c,d$ outside the support
of $g$ and distinct from $a,b$.  With $\tau=(a\,c\,d)$, the commutator
\[
 g\tau g^{-1}\tau^{-1}=(b\,c\,d)(a\,d\,c)=(a\,b\,c)
\]
belongs to $N$.  Every 3-cycle is conjugate to this one by an even
finitary permutation: a finite permutation carrying the ordered triple
to the desired triple can have its parity adjusted on two fresh points.
Every even finitary permutation is a product of 3-cycles, by decomposing
it into an even number of transpositions and replacing pairs of
transpositions by 3-cycles.  Hence $N=A_\Omega$.  Noncommuting 3-cycles
show that the group is nonabelian.  Its commutator subgroup is therefore
a nontrivial normal subgroup, so the group is perfect and $H_1=0$.

Fix four elements $a,b,c,d$ of $\Omega$.  The finite subsets $F$ containing
these elements and with $|F|\geq8$ form a filtered cofinal subsystem of
the finite subsets of $\Omega$.  The normalized bar-chain complex of
$A_\Omega$ is the filtered colimit of those of $\operatorname{Alt}(F)$:
every finite tuple of permutations, and every finite chain, has support
in one finite $F$.  Filtered colimits of abelian groups are exact, so
\[
 H_2(A_\Omega;\Z)
   =\colim_F H_2(\operatorname{Alt}(F);\Z).
\]
By B.2 this is the constant system $\Z/2$ when its generators are chosen
to be $t_F$.  It therefore has colimit $\Z/2$.  Finally, a nontrivial
homomorphism from a simple group is injective.  Such an injection into a
strictly smaller group is impossible.
\end{proof}

\begin{claim}{B.4}{Refuted}{The asserted right adjoint and the first proposed envelope}
There is no $\U$-small category $K$ representing
\[
 X\longmapsto [BX,K(\Z/2,2)]
 \quad\text{on }\Cathot^{\op}.
\]
Consequently $\varphi:\Cathot\to\Hot$ has no right adjoint, fully
faithful or otherwise.  The assertion at V~ms.~37 and the unrestricted
Thomason--Kan-envelope construction at V~ms.~42--43 are false.
\end{claim}
\begin{proof}
Suppose there were natural bijections
$\theta_X:\pi_0\Fun(X,K)\to[BX,K(\Z/2,2)]$ for all $\U$-small $X$.
Taking $X=e$ shows that $\pi_0K$ is a singleton.  Choose an infinite
$\U$-cardinal $\lambda>|\Mor K|$, a set $\Omega$ of cardinality
$\lambda$, and $G=A_\Omega$.  By B.1 and B.3,
$\pi_0\Fun(\mathsf B G,K)=\pi_0K$ is a singleton.  On the other hand,
ordinary unpointed cohomological representability and the universal
coefficient theorem give
\[
\begin{split}
 [BG,K(\Z/2,2)]
 &=H^2(BG;\Z/2)\\
 &\cong\Hom_{\Z}(H_2(BG;\Z),\Z/2)
 \cong\Z/2.
\end{split}
\]
Here the omitted universal-coefficient term is
$\operatorname{Ext}^1_{\Z}(H_1(BG;\Z),\Z/2)=0$, because $G$ is perfect.
The two cohomology classes remain distinct as unpointed homotopy classes:
the target is simply connected.  We have a contradiction.  If a right
adjoint existed, its value at this homotopy type would be such a
representing category, so no right adjoint exists.
\end{proof}

\begin{claim}{B.5}{Proved}{A simple overgroup with a prescribed infinite cardinality}
Let $H$ be a $\U$-small group and let $\lambda$ be an infinite
$\U$-cardinal with $|H|\leq\lambda$.  There is an injective homomorphism
$H\hookrightarrow S$ into a nonabelian simple $\U$-small group $S$ with
$|S|=\lambda$.
\end{claim}
\begin{proof}
We give the construction to make its cardinal control explicit.  Start
with $H_0=H*F_\lambda$, so $H\hookrightarrow H_0$ and $|H_0|=\lambda$.
For an arbitrary group $L$ of cardinality $\lambda$, form
\[
 L'=L*F\{t_{g,h}:g\in L\setminus\{1\},\ h\in L\}.
\]
For each pair $(g,h)$ abbreviate $t=t_{g,h}$ and put
\[
 w_{g,h}=(gtg^{-1}t^{-1})(gt^2g^{-1}t^{-2}).
\]
Both $w_{g,h}$ and $w_{g,h}h$ have infinite order in the free factor
$L*\langle t\rangle$ of $L'$.  The first word is cyclically reduced of
syllable length eight.  For the second, cyclically moving its initial
$g$ to the end gives a cyclically reduced word of length eight when
$hg\ne1$.  When $hg=1$, a further cyclic reduction gives
\[
 g^{-1}t^{-1}gt^2g^{-1}t^{-1},
\]
which has syllable length six.  A cyclically reduced free-product word
of length greater than one has infinite order, since its positive powers
are reduced concatenations.

We may therefore adjoin stable letters $s_{g,h}$ identifying the two
infinite cyclic subgroups by
\[
 s_{g,h}^{-1}w_{g,h}s_{g,h}=w_{g,h}h.
\]
The simultaneous HNN embedding theorem ensures that $L'$ embeds in the
resulting group $E(L)$; its precise form, allowing an arbitrary family
of subgroup isomorphisms, is
\citet[Theorem~II]{BC_HNN1949}.  Equivalently, the HNN normal form says
that a nontrivial word in the base group remains nontrivial after these
stable letters are adjoined.  This theorem is applicable because each
specified map $w_{g,h}^r\mapsto(w_{g,h}h)^r$ is an isomorphism between
infinite cyclic subgroups of the same base group $L'$.  There are
$\lambda$ generators and relations, so $|E(L)|=\lambda$.

In $E(L)$, the element $h=w_{g,h}^{-1}s_{g,h}^{-1}w_{g,h}s_{g,h}$ lies
in the normal closure of $g$: each factor of $w_{g,h}$ already lies in
that normal closure.  Iterate the injective construction
$H_{r+1}=E(H_r)$, and set $S=\bigcup_{r<\omega}H_r$.  This union has
cardinality $\lambda$.  For any $g\ne1$ and $h$ in $S$, one stage
contains both, and the next stage puts $h$ in the normal closure of $g$.
Every nontrivial normal subgroup of $S$ is consequently $S$.  The
embedded free group $F_\lambda$ makes $S$ nonabelian.  This also proves
that none of the inclusions has collapsed the original group $H$.
\end{proof}

\begin{claim}{B.6}{Proved}{An acyclic overgroup without a cardinal increase}
For every infinite $\U$-small group $H$ there is an injective
homomorphism $H\hookrightarrow M(H)$ with
\[
 |M(H)|=|H|,\qquad \widetilde H_i(BM(H);\Z)=0
 \quad\text{for every }i\geq0.
\]
For a finite $H$, the same construction gives a countable overgroup.
\end{claim}
\begin{proof}
We use the mitosis construction of
\citet[Proposition~4.1, Theorem~4.2, and Lemma~5.4]{BC_BDH1980}.
Its algebraic form also supplies a short verification of the size bound.
Embed $H$ as the first factor of $H\times H$, and adjoin two HNN stable
letters $d,s$ satisfying
\[
 d^{-1}(h,1)d=(h,h),\qquad s^{-1}(h,1)s=(1,h)
 \quad(h\in H).
\]
Both identifications are isomorphisms between subgroups of $H\times H$.
The base group embeds, and the resulting group $m(H)$ is generated by
$H,d,s$, because the second factor is conjugate to the first.  In the
notation $h^s=s^{-1}hs$ it satisfies
\[
 h^d=hh^s,\qquad [h',h^s]=1\quad(h,h'\in H).
\]
Put $M(H)=\bigcup_{r<\omega}m^r(H)$ using these injections.  Its
cardinality is at most $\max(|H|,\aleph_0)$, and it contains $H$.

For completeness, the homology-annihilation argument is as follows.
Suppose $f:A\to B$ induces zero in field homology in positive degrees
less than $n$, and $j:B\hookrightarrow m(B)$ is the mitosis inclusion.
The homomorphism
\[
 k:B\times B\longrightarrow m(B),\qquad k(b,c)=b\,c^s
\]
is defined because the two displayed subgroups commute.  The two factor
inclusions followed by $k$ are conjugate to $j$, and the diagonal
followed by $k$ is also conjugate to $j$.  For $z\in H_n(A;k)$ the
K\"unneth formula and the vanishing of the lower-degree maps give
\[
 (f\times f)_*\Delta_*(z)=f_*(z)\otimes1+1\otimes f_*(z).
\]
Applying $k_*$ and using invariance under inner automorphisms yields
$(jf)_*(z)=2(jf)_*(z)$, hence $(jf)_*(z)=0$ over every field.  Induction
shows that a composite of $n$ mitosis inclusions kills field homology
through degree $n$.  Every bar cycle in $M(H)$ occurs at a finite
stage, so all its positive field homology vanishes.  Rational vanishing
makes each positive integral homology group torsion.  The universal
coefficient sequence in the next degree, with coefficients $\F_p$,
shows that it has no $p$-torsion for any prime $p$.  It is therefore zero.
The group is nonempty, so its classifying space is connected as required.
\end{proof}

\begin{claim}{B.7}{Proved}{Acyclic simple overgroups of arbitrary size}
For every group $H$ and every infinite $\U$-cardinal
$\lambda\geq\max(|H|,\aleph_0)$ there is a monomorphism of groups
$i:H\hookrightarrow A$ such that $A$ is nonabelian simple,
$|A|=\lambda$, and $BA$ is integrally acyclic.
\end{claim}
\begin{proof}
Start with $H_0=H*F_\lambda$.  Alternate B.5 and B.6 to obtain
injective homomorphisms
\[
 H_0\hookrightarrow S_0\hookrightarrow M_0\hookrightarrow
 S_1\hookrightarrow M_1\hookrightarrow\cdots,
\]
where every $S_r$ is simple, every $M_r$ is acyclic, and every term has
cardinality $\lambda$.  Let $A$ be the union.  If a normal subgroup
$N\lhd A$ contains $g\ne1$, then for any $a\in A$ a sufficiently late
$S_r$ contains both $g$ and $a$.  The subgroup $N\cap S_r$ is normal in
$S_r$ and nontrivial, hence contains $a$.  This proves simplicity.
Every finite bar cycle lies in some $M_r$ and bounds there, so $A$ is
acyclic.  Its cardinality is $\lambda$, and the inclusions give the
asserted monomorphism from $H$.  The alternating-union observation is a
formal combination of the two preceding constructions; no priority is
claimed for it.
\end{proof}

\begin{claim}{B.8}{Proved}{Groups invisible to smaller groups and mapping to spheres}
For every integer $n\geq2$ and every infinite $\U$-cardinal $\lambda$
there is a $\U$-small group $G_{n,\lambda}$ and a pointed continuous map
of CW models
\[
 q_{n,\lambda}:BG_{n,\lambda}\longrightarrow S^n
\]
such that $|G_{n,\lambda}|=\lambda$, the map $q_{n,\lambda}$ induces
isomorphisms on integral homology in all degrees, and every homomorphism
from $G_{n,\lambda}$ to a group of cardinality less than $\lambda$ is
trivial.  The group $G_{n,\lambda}$ is perfect.  It is not asserted to be
simple.
\end{claim}
\begin{proof}
For a group $H$ with $|H|\leq\lambda$, choose the embedding
$i:H\hookrightarrow A$ of B.7 and form the amalgamated free product
\[
 D_\lambda(H)=A*_{H}A.
\]
The maps from both factors $A$ to this amalgam are injective by the
amalgamated-product normal form.  Thus $|D_\lambda(H)|=\lambda$.
A homomorphism $D_\lambda(H)\to L$ with $|L|<\lambda$ restricts
trivially to each factor, because each factor is simple of cardinality
$\lambda$.  Since the two factors generate the amalgam, the homomorphism
itself is trivial.

The double mapping cylinder of
$BA\leftarrow BH\rightarrow BA$, whose structure maps are $Bi$, is a
classifying CW space for $D_\lambda(H)$.  To see the asphericity required
here, the universal cover is the tree of spaces over the Bass--Serre
tree of the amalgam, with contractible vertex and edge spaces.  A finite
subtree gives a contractible union, by induction over its terminal
vertices; every sphere and homotopy meets a finite subtree after cellular
approximation.  The total cover is therefore weakly contractible.  The
ordinary van Kampen theorem identifies the fundamental group with the
specified amalgam.  This establishes the comparison with a classifying
space without substituting a homotopy pushout for an ordinary pushout
without justification.

Collapse the two copies of $BA$ in this double mapping cylinder to the
two suspension vertices.  We obtain a pointed map
\[
 r_H:B D_\lambda(H)\longrightarrow \Sigma BH.
\]
Here $\Sigma$ denotes the unreduced suspension of the connected CW
space $BH$, with one suspension vertex chosen as basepoint; it agrees
up to pointed homotopy with its reduced suspension after choosing the
standard vertex of $BH$.  The Mayer--Vietoris sequence, compared with
that for the suspension, proves that $r_H$ is an integral homology
isomorphism, because both copies of $BA$ are acyclic.  In particular
$H_1(D_\lambda(H);\Z)=0$.

Start with $G_{1,\lambda}=\Z$ and a chosen pointed homotopy equivalence
$q_{1,\lambda}:B\Z\to S^1$.  Inductively set
\[
 G_{n,\lambda}=D_\lambda(G_{n-1,\lambda}),\qquad
 q_{n,\lambda}=(\Sigma q_{n-1,\lambda})r_{G_{n-1,\lambda}},
\]
using the usual pointed identification $\Sigma S^{n-1}\simeq S^n$.
Suspension preserves integral homology isomorphisms, so induction proves
the homology assertion.  Every group with $n\geq2$ is the specified
double and therefore has all the remaining properties already checked.
\end{proof}

\begin{claim}{B.9}{Proved}{The target condition needed to use these homology equivalences}
Let $u:A\to B$ be a pointed integral homology equivalence between
connected CW spaces, and let $Z$ be a pointed simply connected CW space.
Precomposition induces a weak equivalence of derived pointed mapping
spaces
\[
 u^*:\Map_*(B,Z)\longrightarrow\Map_*(A,Z).
\]
In particular it induces a bijection $[B,Z]_*\to[A,Z]_*$.
\end{claim}
\begin{proof}
For an abelian group $M$, the universal coefficient theorem gives an
isomorphism $\widetilde H^j(B;M)\to\widetilde H^j(A;M)$ for every
$j\geq0$.  The homotopy groups of
$\Map_*(A,K(M,r))$, in every component, are
$\widetilde H^{r-i}(A;M)$ for $0\leq i\leq r$, and vanish for $i>r$;
the mapping space is an abelian group object.  Hence the assertion holds
for every $K(M,r)$.

For a simply connected $Z$, choose its Postnikov tower with
$P_1Z=*$ and with $P_rZ$ the homotopy fibre of its $k$-invariant
\[
 P_{r-1}Z\longrightarrow K(\pi_rZ,r+1).
\]
All coefficient systems in this assertion are constant, precisely
because $Z$ is simply connected.  Derived mapping spaces preserve
homotopy limits in the target.  Induction on $r$, using the homotopy
fibre descriptions, proves the assertion for every $P_rZ$.  Finally
$Z\to\operatorname{holim}_rP_rZ$ is a weak equivalence: for each
homotopy group the tower is eventually constant with the value of that
homotopy group of $Z$.  Taking the homotopy limit of the already proved
mapping-space equivalences gives the assertion for $Z$.  This last step
is essential; we have proved equivalences of mapping spaces at all
stages, rather than merely bijections of their component sets, so no
uncontrolled inverse-limit or phantom-map inference occurs.  The result
is the integral-homology locality of simply connected spaces; see
\citet[Example~9.2]{BC_Lawson2019}.
\end{proof}

\begin{claim}{B.10}{Proved}{Essential maps from arbitrarily large groups for every higher type}
Let $Y$ be a $\U$-small homotopy type with a point $y$ and a nonzero
element $\alpha\in\pi_n(Y,y)$ for some $n\geq2$.  For every infinite
$\U$-cardinal $\lambda$, the group $G_{n,\lambda}$ of B.8 has no
nontrivial homomorphism to any group of cardinality less than $\lambda$
and admits a map $BG_{n,\lambda}\to Y$ which is not freely homotopic to
a constant map.  Thus this assertion holds for every homotopy type that
is not a $1$-type, including every simply connected noncontractible
homotopy type.
\end{claim}
\begin{proof}
Choose a CW model for the component $Y_y$ and its pointed universal
cover $p:\widetilde Y_y\to Y_y$.  The sphere map representing $\alpha$
lifts uniquely as a pointed map
$\widetilde\alpha:S^n\to\widetilde Y_y$, and its class is nonzero
because a covering induces isomorphisms in homotopy degrees at least
two.  By B.9 and B.8,
$\widetilde\alpha q_{n,\lambda}:BG_{n,\lambda}\to\widetilde Y_y$
is not pointed nullhomotopic.

The composite $p\widetilde\alpha q_{n,\lambda}$ cannot be freely
nullhomotopic in $Y$.  If a free nullhomotopy existed, its image would
lie in $Y_y$ and its lift beginning at
$\widetilde\alpha q_{n,\lambda}$ would end at a constant map in the
cover.  The lifted homotopy exists by the covering homotopy property.
Its terminal map is constant because its image lies in the discrete
fibre over the endpoint and $BG_{n,\lambda}$ is connected.  In a
simply connected target, a free nullhomotopy of a pointed map can be
made pointed: the obstruction is the action of the target fundamental
group, which is trivial.  We would obtain a pointed nullhomotopy,
contrary to B.9.  This proves essentiality in the unpointed sense
actually required.  The last sentence follows by choosing a nonzero
higher homotopy group; a connected simply connected CW type with all
homotopy groups zero is contractible by Whitehead's theorem.
\end{proof}

\begin{claim}{B.11}{Proved}{The full negative representability statement}
For every $\U$-small homotopy type $Y$ which is not a $1$-type, the
presheaf $X\mapsto[BX,Y]$ on $\Cathot$ is not representable by a
$\U$-small category.  No connectedness assumption on $Y$ is needed.
\end{claim}
\begin{proof}
Suppose that $K$ represents the presheaf, by natural bijections
$\theta_X:\pi_0\Fun(X,K)\to[BX,Y]$.  Choose
$\lambda>|\Mor K|$ in $\U$, and choose $n$ and the corresponding
$G=G_{n,\lambda}$ from B.10.  Naturality for the unique functor
$p:\mathsf B G\to e$ gives a commutative square
\[
\begin{array}{ccc}
\pi_0K&\xrightarrow{\ p^*\ }&\pi_0\Fun(\mathsf B G,K)\\
\theta_e\downarrow\phantom{\theta_e}&&\phantom{\theta_{\mathsf B G}}\downarrow\theta_{\mathsf B G}\\
\pi_0Y&\xrightarrow{\ (Bp)^*\ }&[BG,Y].
\end{array}
\]
The upper arrow is a bijection by B.1.  The lower arrow sends a component
of $Y$ to its constant-map class.  It is not surjective by B.10.  This
contradicts the vertical bijections.  This argument uses the specified
constant-map function, rather than the cardinalities of the two sets;
it therefore also applies when $Y$ has infinitely many components.
\end{proof}

\partsection{C}{Local categories and the groupoid reflection}

For a category $X$, let $\Pi_1X=X[\Mor(X)^{-1}]$ denote the groupoid
obtained by inverting all its arrows, and let
$\eta_X:X\to\Pi_1X$ be the canonical functor.  Its objects are those
of $X$; its arrows are finite composable words of arrows and formal
inverses, modulo the identity, inverse, and composition relations.
This construction is $\U$-small.  It is equivalent to the fundamental
groupoid of $BX$ on the vertices of its nerve: the nerve's edges give
the generating paths, its $2$-simplices impose the composition
relations, and higher cells impose no further relations on the
fundamental groupoid.  We allow the empty groupoid.

\begin{claim}{C.1}{Proved}{Groupoids are local, with the canonical comparison}
For every $\U$-small category $X$ and $\U$-small groupoid $G$, the
canonical map
\[
 \pi_0\Fun(X,G)\longrightarrow[BX,BG]
\]
is a bijection, natural in both variables.  More precisely,
\[
 N\Fun(X,G)\cong\operatorname{Map}_{\mathrm{sSet}}(NX,NG)
\]
is an isomorphism of simplicial sets, and the right side is a derived
mapping space.  Restriction along $\eta_X$ is an isomorphism of
categories $\Fun(\Pi_1X,G)\to\Fun(X,G)$.
\end{claim}
\begin{proof}
In degree $r$, the asserted simplicial isomorphism is the chain of
canonical bijections
\[
\begin{split}
N\Fun(X,G)_r
 &=\Fun([r],\Fun(X,G))\\
 &=\Fun(X\times[r],G)\\
 &=\Hom_{\mathrm{sSet}}(NX\times\Delta[r],NG).
\end{split}
\]
The last equality uses full faithfulness of the nerve and preservation
of products.  These bijections commute with faces and degeneracies.
The nerve of a groupoid is a Kan complex, and every simplicial set is
cofibrant in the Kan--Quillen model structure.  Thus the displayed
mapping space already computes derived maps of spaces, giving precisely
the canonical map on components.  This proves localness, including for
$X$ or $G$ empty.

Every functor $X\to G$ extends uniquely over $\eta_X$, since its values
on arrows are invertible.  A natural transformation between two such
functors has invertible components.  Its naturality square for an arrow
of $X$ also gives the naturality square for its formal inverse, by
multiplying the square by the inverses.  Therefore transformations extend
uniquely as well, and restriction is the asserted isomorphism of
categories.  This last assertion concerns ordinary categories, before
passing to components.
\end{proof}

\begin{claim}{C.2}{Proved}{Classification of all representable homotopy types}
A $\U$-small homotopy type $Y$ has its presheaf
$X\mapsto[BX,Y]$ represented on $\Cathot$ by a $\U$-small category if
and only if $Y$ is a $1$-type.  If $Y$ is a $1$-type, a fundamental
groupoid $G_Y$ represents it.  Every other representing category is
isomorphic to $G_Y$ in $\Cathot$, and its classifying space is
isomorphic to $Y$ in $\Hot$.  This is a complete affirmative answer to
Lead~2's representability conjecture.
\end{claim}
\begin{proof}
The negative direction is B.11.  If $Y$ is a $1$-type, choose its
fundamental groupoid, with one representative object for each component
and automorphism group $\pi_1(Y,y)$ there.  This is $\U$-small, and
$BG_Y\simeq Y$: the comparison identifies components and fundamental
groups, and both sides have zero higher homotopy groups.  C.1 gives
the required natural bijections.

If a category $K$ represents the same presheaf, the resulting natural
isomorphism
$\Hom_{\Cathot}(-,K)\cong\Hom_{\Cathot}(-,G_Y)$ is represented by an
isomorphism $K\cong G_Y$ in $\Cathot$, by the ordinary Yoneda lemma.
Applying $B$ gives an isomorphism in $\Hot$.  More precisely, if
$\theta$ is the specified representing isomorphism, its universal
element $\theta_K(\id_K):BK\to Y$ is this weak equivalence in
$\Hot$, composed with the chosen identification $BG_Y\simeq Y$.
Thus an unspecified abstract bijection has not been substituted for the
canonical comparison.  When $Y$ is empty the representing category is
the empty groupoid, and the same Yoneda argument applies.
\end{proof}

\begin{claim}{C.3}{Proved}{Classification of local categories}
For a $\U$-small category $K$, the following conditions are equivalent:
the canonical function $\pi_0\Fun(X,K)\to[BX,BK]$ is bijective for
every $\U$-small $X$; the category $K$ is isomorphic in $\Cathot$ to a
groupoid; and the canonical functor $\eta_K:K\to\Pi_1K$ is an
isomorphism in $\Cathot$.  In particular every local $K$ has $BK$ a
$1$-type.  Conversely, for a category with $BK$ a $1$-type, localness
is equivalent to being isomorphic in $\Cathot$ to its groupoidification.
The hypothesis that $BK$ is a $1$-type alone is insufficient.
\end{claim}
\begin{proof}
If $K$ is local, it represents $X\mapsto[BX,BK]$.  C.2 implies that
$BK$ is a $1$-type.  Hence $B\eta_K:BK\to B\Pi_1K$ is a weak
equivalence: it always induces isomorphisms on components and
fundamental groups, and in this case both source and target have no
higher homotopy groups.

Localness at the source $X=\Pi_1K$ produces a morphism
$[j]:\Pi_1K\to K$ in $\Cathot$ whose realization represents
$(B\eta_K)^{-1}$ in $\Hot$.  Localness of $K$ at $X=K$ then implies
$[j][\eta_K]=\id_K$.  C.1, applied to the groupoid $\Pi_1K$, gives
$[\eta_K][j]=\id_{\Pi_1K}$.  These are equalities in $\Cathot$;
they do not assert actual equalities of functors or an equivalence of
the underlying categories.

The third condition implies the second.  To prove the converse, we
observe that postcomposition by an isomorphism $K\to G$ in $\Cathot$
induces bijections on all representable hom-sets, and realization of
that isomorphism induces bijections $[BX,BK]\to[BX,BG]$.  The square
formed by the two canonical comparisons commutes, and the comparison
for $G$ is bijective by C.1.  Thus the comparison for $K$ is bijective.
The failure of the weaker converse is exhibited in C.5 below.
\end{proof}

\begin{claim}{C.4}{Proved}{The reflector is fundamental groupoidification}
Let $\mathcal L\subset\Cathot$ be the full subcategory of local
objects.  It is reflective.  After identifying it with groupoids modulo
natural isomorphism, its reflector is
\[
 \Pi_1:\Cathot\longrightarrow\operatorname{Ho}(\mathrm{Grpd}),
 \qquad \Pi_1\dashv I,
\]
where $I$ is the full inclusion and the unit is $[\eta_X]$.
This reflector inverts exactly those functors that induce a bijection
on components and isomorphisms on fundamental groups at every
basepoint.  Its unit is in $W_\infty$ precisely when $BX$ is a
$1$-type.  In particular this reflection does not supply the nonexistent
right adjoint to $\varphi:\Cathot\to\Hot$.
\end{claim}
\begin{proof}
Applying $\Pi_1$ to a natural transformation produces a natural
isomorphism.  It therefore descends to $\Cathot$.  C.1 gives natural
bijections
\[
 \Hom_{\operatorname{Ho}(\mathrm{Grpd})}(\Pi_1X,G)
   =\pi_0\Fun(\Pi_1X,G)
   \cong\pi_0\Fun(X,G)
   =\Hom_{\Cathot}(X,IG).
\]
This is the asserted adjunction, with the displayed unit.  C.3
identifies its replete essential image with $\mathcal L$.
A functor of groupoids is an equivalence exactly when it gives the
specified bijection of components and isomorphisms of isotropy groups;
thus the maps inverted are exactly the stated $1$-equivalences.
For every $X$, the map $B\eta_X$ preserves components and fundamental
groups and has a $1$-type as its target.  It is a weak equivalence if
and only if all the higher homotopy groups of $BX$ vanish.

Although every $W_\infty$-map is inverted by this reflector, many more
maps are inverted.  For example a category with classifying space $S^2$
maps to a contractible groupoidification.  The target $\mathcal L$
therefore models $1$-truncation; it does not model the localization
$\Hot$ at $W_\infty$.  This distinction is necessary when using the
word ``reflective'' here.
\end{proof}

\begin{claim}{C.5}{Refuted}{Coreflectivity, and the weaker converse from the classifying space}
The inclusion $I:\mathcal L\hookrightarrow\Cathot$ has no right
adjoint.  Moreover, a category with contractible classifying space need
not be local.
\end{claim}
\begin{proof}
Let $F$ be the poset with object set $\N$, identity arrows, and exactly
the following nonidentity arrows:
\[
 2r\longrightarrow2r+1\quad(r\geq0),\qquad
 2r\longrightarrow2r-1\quad(r\geq1).
\]
There are no composable pairs of nonidentity arrows, so composition
with identities specifies the entire category.  Its nerve has only
vertices and these nondegenerate edges.  Its realization is the
infinite ray, hence is contractible.

Let $D$ be the discrete category with object set $\N$.  Define functors
$f,g:D\to F$ by $f(r)=0$ and $g(r)=2r$.  A natural transformation
between two functors $D\to F$ is an objectwise family of arrows.  Thus
each one-step transformation, in either direction, changes each
coordinate by at most one in the undirected path metric on $F$.
A zigzag of length $m$ changes every coordinate by at most $m$.
Since the distance from $f(r)$ to $g(r)$ is $2r$, no finite zigzag
connects $f$ to $g$.  Hence
\[
 |\Hom_{\Cathot}(D,F)|>1,
 \qquad |\Hom_{\Cathot}(e,F)|=1.
\]
In particular the canonical comparison to $[BD,BF]$, which is a
singleton, is not injective, proving that $F$ is not local.  The
finite truncations and their exact distances are checked by
\texttt{code/bc\_checks.py}; the unbounded-coordinate argument proves
the infinite assertion.

Within $\mathcal L$, the object $D$ is the coproduct of countably many
copies of $e$.  Indeed groupoids have this coproduct, and a family of
componentwise natural isomorphisms can be assembled into a single
natural isomorphism on a discrete union.  If $I$ had a right adjoint,
then, being a left adjoint, it would preserve this coproduct.
Equivalently its adjunction would force the canonical restriction map
\[
 \Hom_{\Cathot}(D,F)\longrightarrow
     \prod_{r\in\N}\Hom_{\Cathot}(e,F)
\]
to be a bijection.  Its target is a singleton and its source has the
two distinct elements just constructed.  This contradiction proves
the failure of coreflectivity.  The use of an infinite source is part
of the counterexample, not a suppressed finiteness assumption.
\end{proof}

\begin{claim}{C.6}{Refuted}{Thomason fibrancy is not the manuscript's localness}
It is false that every Thomason-fibrant category is local in the sense
of C.3.  There is no change of target by a $W_\infty$-equivalence that
makes a homotopy type with nonzero higher homotopy groups into a local
category.
\end{claim}
\begin{proof}
Let $P$ be the face poset of the boundary of a tetrahedron, and choose
a Thomason fibrant replacement $P\to RP$.  By the definition of
fibrant replacement this is a specified $W_\infty$-map, and
$B(RP)\simeq BP\simeq S^2$.  If $RP$ were local, C.3 would make this
a $1$-type, a contradiction.  The same reasoning applies to any
category whose nerve has a nonzero higher homotopy group and to any
$W_\infty$-equivalent target.  Fibrancy can be used for derived
mapping objects with the appropriate source and cylinder; it cannot
produce the envelope asserted at V~ms.~37--43.
\end{proof}

\partsection{D}{Bounded tests and ind-representations}

\begin{claim}{D.1}{Proved}{Exact bounded representability}
Let $Y$ be a $\U$-small Kan complex and let $\lambda\in\U$ be an infinite cardinal with $|Y|\leq\lambda$, where size counts simplices in all degrees. There is a $\U$-small category $K_\lambda(Y)$ and a natural transformation
\[
 \theta_{X,\lambda}:\pi_0\Fun(X,K_\lambda(Y))\longrightarrow[NX,Y]
 \qquad (X\in\Cat_\U)
\]
which is bijective whenever $|\Mor X|\leq\lambda$. One can take $|\Mor K_\lambda(Y)|\leq2^\lambda$. Consequently, for every infinite $\kappa\in\U$, the functor in Part D(i), on categories with fewer than $\kappa$ arrows, is represented by a small category after allowing the representing category to exceed the test bound.
\end{claim}
\begin{proof}
Take all contractible simplicial sets $S$ having at most $\lambda$ simplices, encoded with each set $S_n$ a subset of a fixed set of size $\lambda$, together with a simplicial map $a:S\to Y$. Let $K_\lambda(Y)$ have these pairs as objects and all simplicial maps over $Y$ as arrows. The collection of structures, maps, and subsets has size at most $2^\lambda$. Relabeling a simplicial set produces an isomorphic object; thus this choice of codes loses no object of the required size. Here ``contractible'' means weakly equivalent to the point, so in particular nonempty. If $Y$ is empty the resulting category is empty, as it must be.

A functor $F:X\to K_\lambda(Y)$ is exactly an objectwise contractible diagram $A:X\to\mathbf{sSet}$ equipped with $a:A\to Y_X$, the constant diagram at $Y$. Its augmentation $A\to *_X$ is an objectwise weak equivalence. In the projective model category of $X$-diagrams it therefore defines the morphism
\[
 *_X\xleftarrow{\simeq}A\xrightarrow{a}Y_X
 \quad\text{in }\Ho(\mathbf{sSet}^X).
\]
Every natural transformation over $Y$ between such diagrams is objectwise a weak equivalence, by two-out-of-three. It gives the same morphism. This constructs the comparison $\theta$ after identifying this Hom-set with $[NX,Y]$. That identification follows directly from the projectively cofibrant resolution
\[
 Q_X(x)=N(X/x),\qquad Q_X\longrightarrow *_X.
\]
Each value is contractible, with terminal object $\id_x$. Filtering by nondegenerate chains exhibits $Q_X$ as built by attaching the projective cells $X(z,-)\times\partial\Delta^n\to X(z,-)\times\Delta^n$. Its ordinary colimit is $NX$, so the simplicial mapping object from $Q_X$ to $Y_X$ is $Y^{NX}$. The latter is Kan. This proves the identification, and also its naturality under restriction along $u:X'\to X$: restriction preserves objectwise weak equivalences, and the natural comparison $Q_{X'}\to u^*Q_X$ is compatible with $Nu$.

Suppose now $|\Mor X|\leq\lambda$. The resolution $Q_X$ has at most $\lambda$ simplices at each object, and the entire diagram has size at most $\lambda$. Every representative $g:NX\to Y$ yields $Q_X\to Y_X$, hence a functor into $K_\lambda(Y)$ after objectwise relabeling. This proves surjectivity.

For injectivity, start with $A\to Y_X$ as above and factor it projectively as a trivial cofibration $A\to A^f$ followed by a fibration $A^f\to Y_X$. The small-object construction uses maps from finite horns tensored with covariant representables. At each countable stage there are at most $\lambda$ lifting problems, and each attaching cell has at most $\lambda$ simplices in the whole diagram. Hence $A^f$ can be chosen with values of size at most $\lambda$. Since $Y$ is Kan, $A^f\to *_X$ is a projective trivial fibration. Cofibrancy of $Q_X$ gives a lift $Q_X\to A^f$ over the terminal diagram $*_X$. Its composite with $A^f\to Y_X$ defines $g:Q_X\to Y_X$; no lift over a preassigned $g$ is required. Thus the original functor is joined, inside $\Fun(X,K_\lambda(Y))$, to the functor supplied by $g$ by the zigzag $A\to A^f\leftarrow Q_X$.

If two resulting maps $g_0,g_1$ represent the same derived morphism, the Kan function complex gives a simplicial homotopy $Q_X\times\Delta^1\to Y_X$ between them (or a finite chain of edges, which gives the same conclusion). Its values are contractible and of size at most $\lambda$. The two endpoint inclusions give a zigzag of natural transformations over $Y$. Consequently the original functors lie in the same component. This proves injectivity with the actual finite-zigzag equivalence relation. All relabelings used above are objectwise isomorphisms and hence give natural isomorphisms of the relabeled diagrams. They do not require any unproved strictification.
\end{proof}

\begin{claim}{D.2}{Proved}{Finite targets are not enough in the test examples}
On the full category of finite categories, neither $X\mapsto[BX,S^2]$ nor $X\mapsto[BX,K(C_2,2)]$ has a finite representing category in $\Cathot$. Both have small, generally infinite, representing categories by D.1. For $K(C_2,2)$ and an uncountable limit cardinal $\kappa$, a representative for all tests with fewer than $\kappa$ arrows cannot itself have fewer than $\kappa$ arrows.
\end{claim}
\begin{proof}
If $K$ is finite, $\Fun(P,K)$ is finite for each finite $P$. Take a finite face poset $P$ of a triangulated two-sphere. The set $[BP,S^2]=\Z$ is infinite, giving the first obstruction. For the second, choose $n\geq8$ with $|A_n|>|\Mor K|$. Simplicity of $A_n$ makes every functor $BA_n\to K$ trivial on its arrows. Thus $\Fun(BA_n,K)\cong K$, while evaluation at the point would require $\pi_0K$ to be a singleton. The Schur multiplier computation used in Part B gives $H^2(A_n;C_2)=C_2$. This is a contradiction on a finite test category. For a limit $\kappa$ and an infinite candidate size $\mu<\kappa$, choose $\mu<\lambda<\kappa$. The finitary alternating group on $\lambda$ letters has $\lambda$ elements and gives the same contradiction by Part B. A finite candidate is already excluded by the preceding finite-group argument. No successor-cardinal lower bound beyond what this argument proves is being asserted.
\end{proof}

\begin{claim}{D.3}{Proved}{Finite categories always give a filtered ind-representation}
Let $\mathcal C_f$ be a skeleton of finite categories in $\Cathot$, with all functors modulo natural-transformation zigzags as morphisms. For every $\U$-small homotopy type $Y$, put $F_Y(C)=[BC,Y]$. The category of elements $\mathcal E_Y$ of $F_Y|_{\mathcal C_f}$ is a small filtered category, and
\[
 F_Y|_{\mathcal C_f}\cong
 \colim_{(C,c)\in\mathcal E_Y}\Hom_{\Cathot}(-,C)|_{\mathcal C_f}.
\]
Thus Part D(ii) has an affirmative answer for filtered ind-representability by finite categories. The assertion does not require that a finite category have homotopy-finite nerve.
\end{claim}
\begin{proof}
An object is a pair $(C,c)$ with $c:BC\to Y$ a morphism of $\Hot$; an arrow $(C,c)\to(D,d)$ is a class $u:C\to D$ in $\Cathot$ with $c=d\,Bu$. The object $(\varnothing,!)$ proves nonemptiness, including $Y=\varnothing$. Coproducts of finite categories give a common successor to any two objects because maps out of the disjoint union are pairs of maps.

Let $u,v:(A,a)\rightrightarrows(B,b)$ be parallel arrows. Choose representative functors. Form the finite Grothendieck construction $C$ of the diagram with two objects and two arrows $A\rightrightarrows B$. Explicitly its objects are the disjoint union of those of $A$ and $B$. Its internal arrows are those of $A$ and $B$; an arrow from an $A$-object $a_0$ to a $B$-object $b_0$ consists of a tag $t\in\{u,v\}$ and an arrow $t(a_0)\to b_0$ of $B$. There are no arrows in the other direction. Precomposition uses $t$ on arrows of $A$, and postcomposition is composition in $B$. The two tagged identity arrows give natural transformations from the inclusion of $A$ to the two composites $wu,wv$, where $w:B\to C$ is the inclusion. Hence $wu=wv$ in $\Cathot$.

By the homotopy-colimit theorem \citep[Theorem~26.8 and Corollary~26.9]{ChacholskiScherer2001}, $BC$ is naturally weakly equivalent to the double mapping cylinder of $Bu,Bv$, hence to their homotopy coequalizer, compatibly with the copy of $BB$. Since $bBu=bBv$ in $\Hot$, choose an actual homotopy between representatives of these composites into a fibrant space representing $Y$. It extends $b$ over that homotopy coequalizer and yields $c:BC\to Y$ with $cBw=b$. Thus $(B,b)\to(C,c)$ equalizes the pair in $\mathcal E_Y$. This proves filteredness. The definition is independent of the chosen representatives, since the equalization needed is an equality in $\Cathot$.

For completeness, the displayed colimit sends $f:X\to C$ at $(C,c)$ to $cBf$. It is surjective because an element $x\in F_Y(X)$ is the image of $\id_X$ at $(X,x)$. If two representatives have equal image $x$, the arrows from $(X,x)$ to their indexing objects identify both with this same identity representative in the colimit. This is exactly the presheaf category-of-elements proof of density. Smallness follows from the countable skeleton of finite composition tables and the small sets $[BC,Y]$.
\end{proof}

\begin{claim}{D.4}{Proved}{The exact sequential criterion}
For $Y\in\Hot_\U$, the restriction $F_Y|_{\mathcal C_f}$ is a colimit of a sequence of representables by finite categories if and only if $[BC,Y]$ is countable for every finite category $C$. Such a sequence may be represented by actual functors of finite categories. Its colimit $K$ in $\Cat$ is a countable locally finite category and represents $F_Y$ on finite tests. Conversely, every countable locally finite representing category supplies such a sequence. Here a category is locally finite when each finitely generated subcategory has finitely many arrows.
\end{claim}
\begin{proof}
Each $\Hom_{\Cathot}(C,K_n)$ is finite when both categories are finite. A sequential colimit of these sets is countable. Conversely, the hypothesis makes $\mathcal E_Y$ countable: it has countably many objects and each of its Hom-sets is finite. Every countable filtered category admits a cofinal sequence. One constructs it by enumerating its objects and arrows, successively choosing common successors for the next object and equalizing the next finite set of parallel arrows; a diagonal enumeration ensures that each comma category is nonempty and connected. Applying this to D.3 gives the required sequence in $\Cathot$. Choosing representative functors on adjacent arrows defines an actual sequence in $\Cat$, since the ordered indexing category has no relations other than those imposed by composition.

A finite category is finitely presentable as a category: its finitely many objects and arrows, together with its finite composition and identity equations, give a finite presentation. It follows, by the same finite-data argument also for components of natural transformations, that
\[
 \Fun(C,\colim_nK_n)\cong\colim_n\Fun(C,K_n),\qquad
 \pi_0\Fun(C,\colim_nK_n)=\colim_n\pi_0\Fun(C,K_n).
\]
The second equality needs only that each zigzag is finite. The colimit is countable. Any finite generating family of arrows lifts to one finite stage. Move to a further common stage where every equality among the finitely many endpoint objects that holds in the colimit already holds. Every composable word in the chosen arrows then lifts to a composable word in that finite category. Its image is consequently contained in the image of a finite arrow set, so the generated subcategory is finite. This endpoint-equality step is necessary when transition functors identify objects. Hence the colimit is locally finite. Conversely, enumerate a countable locally finite category's arrows and take the finite subcategories they generate successively. Their union is the given category, and the same argument proves the representation. Full inclusions of finite stages are not required for this assertion.
\end{proof}

\begin{claim}{D.5}{Proved}{The Eilenberg--Mac Lane sequential example}
For $Y=K(C_2,2)$, the restriction to finite categories is ind-representable by a sequence of finite categories and representable by a countable locally finite category, although it is not representable on all $\Cathot_\U$ by any $\U$-small category.
\end{claim}
\begin{proof}
For finite $C$, the normalized degree-two cochain group of $NC$ with $C_2$ coefficients is finite: there are finitely many strings of two nonidentity arrows. Its subquotient $H^2(BC;C_2)$ is therefore finite. D.4 applies. The final assertion is Part B. The two statements have different ranges of the test variable $X$, so there is no contradiction.
\end{proof}

\begin{claim}{D.6}{Open}{The sequential two-sphere test}
The remaining exact question in Part D(ii) for $Y=S^2$ is whether $[BC,S^2]$ is countable for every finite category $C$. By D.4 this is equivalent to sequential ind-representability by finite categories; it is also equivalent to representability on finite tests by a countable locally finite category. General filtered ind-representability and ordinary small representability on these tests are already proved in D.3 and D.1. No assertion of sequential representability for $S^2$ is certified here.
\end{claim}
The nerve of a finite category can have infinitely many nondegenerate simplices, so finite-source compactness does not prove the required countability. Finite group categories are not an obstruction: the Sullivan theorem \citep[Theorem~A]{Miller1984} makes maps from their classifying spaces to $S^2$ constant up to unpointed homotopy. The finite monoids with homology in infinitely many degrees discussed in the research report do not by that information alone provide uncountably many maps to $S^2$. These are the two attempted reductions, and neither decides the stated question. This is an unresolved item of this audit, not a claim that the literature contains no answer.

\begin{claim}{D.7}{Proved}{The unbounded cardinal-indexed representation}
There is a compatible choice of the categories $K_\lambda(Y)$ for $\lambda\geq\max(\aleph_0,|Y|)$ and inclusion functors for increasing $\lambda$, such that, in the larger universe $\V$,
\[
 [B-,Y]\cong\colim_{\lambda\in\operatorname{Card}(\U),\ \lambda\geq\max(\aleph_0,|Y|)}
 \Hom_{\Cathot_\U}(-,K_\lambda(Y)).
\]
The indexing poset is filtered, $\V$-small, and unbounded in the $\U$-cardinals. Alternatively $F_Y$ is a $\V$-small filtered colimit over its category of elements on all $\Cathot_\U$.
\end{claim}
\begin{proof}
Use the subset encodings of D.1 on nested cardinal ordinals. The categories are then full subcategories of one another. The maps $\theta$ arise from the same roofs of diagrams and are compatible with the inclusions. For each fixed $\U$-small $X$, they are bijective from every stage $\lambda\geq|\Mor X|$ onward. This proves the displayed isomorphism pointwise and naturally. The index is a set in $\V$, but has no $\U$-small cofinal subset. The proof of D.3 applies verbatim with arbitrary $\U$-small categories because the two constructions used there, finite coproduct and the displayed finite-shape Grothendieck construction, preserve $\U$-smallness. Its whole category of elements is $\V$-small.
\end{proof}

\begin{claim}{D.8}{Proved}{Why the larger index is necessary}
If $Y$ has a component with a nonzero homotopy group in degree at least two, $F_Y$ on all $\Cathot_\U$ is not a $\U$-small colimit of representables, even without requiring the colimit to be filtered. If $Y$ is a 1-type it is representable by its fundamental groupoid. Thus, for the global functor, enlarging the index universe in D.7 is essential exactly outside the 1-types.
\end{claim}
\begin{proof}
Suppose $F_Y=\colim_{i\in I}h_{K_i}$ with $I\in\U$. Choose a $\U$-cardinal $\mu$ bounding the arrow sets of all $K_i$ simultaneously. Part B constructs a $\U$-small group $G$ with no nontrivial homomorphism into a group of cardinal at most $\mu$ and with an essential map $BG\to Y$. Every functor $BG\to K_i$ is therefore trivial on arrows; evaluation at the unique object is an actual isomorphism $\Fun(BG,K_i)\cong K_i$, natural in $K_i$. It identifies the map induced by $BG\to e$ with the identity on $\pi_0K_i$. Taking any colimit gives a bijection $F_Y(e)\to F_Y(BG)$. Its image consists precisely of the homotopy classes of constant maps, contradicting the constructed essential map. The 1-type assertion is Part C. The argument includes disconnected $Y$ because the essential map is chosen in one specified component.
\end{proof}

\begin{claim}{D.9}{Known}{The bounded conclusion also follows from Raptis's model structures}
Let $\mu\in\U$ be a regular infinite cardinal. There is a model structure $\Cat_{I(\mu)}$ with weak equivalences $\W_\infty$ in which every $\mu$-presentable category is cofibrant. If $K\to R_\mu K$ is a fibrant replacement in this structure, then the canonical comparison
\[
 \pi_0\Fun(X,R_\mu K)\longrightarrow[BX,BK]
\]
is a bijection for every $\mu$-presentable $X$, in particular for every $X$ with fewer than $\mu$ arrows. Thus bounded representability is also a consequence of a known model structure; the construction in D.1 additionally gives its explicit size bound and compatible cardinal filtration.
\end{claim}
\begin{proof}
The model structure and cofibrancy assertion are \citet[Theorem~3.5]{Raptis2010}. The product discussion at the end of that paper's Section~3 proves that $X\times-$ is left Quillen in $\Cat_{I(\mu)}$ whenever $X$ is $\mu$-presentable. Let $Z_2=c\operatorname{Sd}^2\Delta^1$, the five-object fence with four edges. Its two endpoint inclusions jointly form a cofibration, and its projection to $e$ is a weak equivalence. Hence
\[
 X\amalg X\longrightarrow X\times Z_2\longrightarrow X
\]
is a cylinder object for $X$ in this model structure. This is also the explicit cylinder in \citet[Proposition~4.8 and its proof]{Raptis2010}.

The cofibrant--fibrant computation identifies $[BX,BK]$ with functors $X\to R_\mu K$ modulo left homotopy through this cylinder. Such a left homotopy is a four-step natural-transformation zigzag, so equality in the model homotopy category implies equality in $\pi_0\Fun$. Conversely every natural transformation becomes equality on classifying spaces, by A.1, so every finite zigzag identifies model-homotopy classes. The two quotients therefore agree, with the asserted canonical comparison. For a prescribed cardinal bound $\kappa$, choose regular $\mu>\kappa$ in $\U$, and choose a small categorical model of $Y$ before applying this replacement. This supplies the alternative existence proof without asserting that a single replacement is local for all small sources.
\end{proof}

\partsection{E}{Corrections at the manuscript's uses of envelopes}

We use ``local'' only in the sense of C.3, and ``fibrant'' only for a specified model structure. At IX ms.~9 the letter $\psi$ denotes the restricted Yoneda functor $\Hot\to\Hot_0^{\wedge}$, not the right adjoint asserted at V ms.~37. Its existence is not in dispute. Its proposed detection properties are false. This distinction is necessary for reading the same symbol across the chapters.

\begin{claim}{E.1}{Proved}{V ms.~37--44: the permissible replacements}
The formal propositions on V ms.~37--42 remain theorems with the adjunction hypothesis in A.3--A.4. Their specialization to a right adjoint of $\varphi$ is refuted by B.4. For an individual category $K$, a local category weakly equivalent to $K$ exists if and only if $BK$ is a 1-type; in that case $\eta_K:K\to\Pi_1K$ is such a weak equivalence. For arbitrary $K$, a Thomason fibrant replacement exists, with the mapping property of A.5. For a fixed bounded family of sources, D.1 gives a target that represents the required Hom-functor on that family. For a weak-equivalence-invariant ind-represented functor, A.7 gives its descended ind-object in $\Hot$ without any envelope.
\end{claim}
\begin{proof}
Necessity for the individual local replacement follows from C.3; sufficiency follows because $\Pi_1K$ is local and $B\eta_K$ is a weak equivalence exactly for 1-types. The other assertions have the stated proofs and distinct universal properties. In particular a local category and a fibrant category cannot be substituted for one another in the sentence ``\emph{l'enveloppe de Thomason-Kan}'' at V ms.~42. The unit of the 1-truncation reflection is not a $\W_\infty$-map in the general case.
\end{proof}

\begin{claim}{E.2}{Proved}{V ms.~45: the uniform-stage lemma}
Let $K_0\subset K_1\subset\cdots$ be full subcategories with union $K_\infty$. If $X$ has finitely many objects, with no restriction on its number of arrows, the canonical comparison is an isomorphism of categories
\[
 \colim_n\Fun(X,K_n)\longrightarrow\Fun(X,K_\infty),
\]
and therefore a bijection on components. For an arbitrary filtered diagram of categories the same conclusion holds if $X$ has finitely many arrows. Neither conclusion identifies $\pi_0\Fun(X,K_\infty)$ with $[BX,BK_\infty]$ without an additional argument.
\end{claim}
\begin{proof}
For full inclusions, the finitely many images of objects of $X$ lie at one stage. Fullness puts all the images of arrows there, even if $X$ has infinitely many arrows. A natural transformation has one component for each object of $X$, and its components also lie in a common full stage. Every naturality equation then holds at that stage because its inclusion is faithful. This proves both surjectivity and injectivity on all Hom-sets and objects of the displayed functor. A path in a category is finite, so connectedness of two functors in the union is witnessed at a common stage. For arbitrary filtered diagrams, use the finite presentation argument of D.4. The final caution is witnessed by A.6 and by E.4 below; the comparison to derived maps is a different map.
\end{proof}

\begin{claim}{E.3}{Proved}{V ms.~43 and 46: what fibrant replacement and filtered union supply}
An actual sequence $K_n\to K_{n+1}$ admits a strictly commuting map to a sequence of Thomason-fibrant categories $R_n$ with each $K_n\to R_n$ a trivial cofibration. Its induced map of filtered colimits is in $\W_\infty$. Filtered colimits of Thomason-fibrant categories are Thomason-fibrant. In contrast, filtered unions of local categories need not be local, even with finite full stages.
\end{claim}
\begin{proof}
Choose trivial cofibrations $\alpha_n:K_n\to R_n$ with fibrant targets. The lifting problem whose left side is $\alpha_n$, whose right side is $R_{n+1}\to e$, and whose top composite is $K_n\to K_{n+1}\xrightarrow{\alpha_{n+1}}R_{n+1}$ produces $R_n\to R_{n+1}$. Adjacent squares suffice because the indexing sequence is freely generated. Thomason cofibrations are full sieve inclusions: the corrected statement uses pseudo-Dwyer maps rather than the false assertion that Dwyer maps are closed under retracts; see \citet[Proposition~2.4 and the following discussion]{Raptis2010}.

Nerves commute with filtered colimits, since a simplex and its defining composition equations involve finite data. Weak equivalences of simplicial sets are preserved by filtered colimits; one may verify this by a functorial Kan replacement using finite horns, followed by the finite simplicial definitions of homotopy groups. Hence the colimit of the $\alpha_n$ is a weak equivalence. The generating Thomason horn domains and codomains $c\operatorname{Sd}^2\Lambda_i^r\hookrightarrow c\operatorname{Sd}^2\Delta^r$ are finite categories. A lifting problem into a filtered colimit factors through one stage and receives a filler there, proving the fibrancy assertion. The last assertion is the explicit example in E.4. Thus the sentence at V ms.~46 is true with ``fibrant'' but false with ``local''.
\end{proof}

\begin{claim}{E.4}{Refuted}{V ms.~47--49: an ind-isomorphism need not induce a weak equivalence of unions}
There are finite full filtrations $T_0\subset T_1\subset\cdots$, every $T_n$ local and Thomason-fibrant, with fibrant union $T$, and functors $s,\id_T:T\rightrightarrows T$ such that they induce the same morphism of $(T_n)$ in $\Ind(\Cathot)$, but $Bs$ is not a weak equivalence and is not homotopic to the identity. Thus even the additional fibrancy of every stage and both unions does not rescue either the union-isomorphism inference or the uniqueness question of V ms.~47--49.
\end{claim}
\begin{proof}
Let $G_n$ be the symmetric group on $\{0,\ldots,n\}$, with the inclusions extending permutations by fixed points. Define $T_n$ with objects $0,\ldots,n$ and
\[
 \Hom_{T_n}(i,j)=\begin{cases}G_j&i\leq j,\\\varnothing&i>j.\end{cases}
\]
If $a:i\to j$ has label $g\in G_j$ and $b:j\to k$ has label $h\in G_k$, put $b\circ a=h\,\iota_{jk}(g)$. Group associativity and compatibility of the inclusions prove associativity, and the identity label at each object is its identity arrow. The inclusions $T_n\to T_{n+1}$ are full. The functor $r_n:T_n\to\mathsf BG_n$ sends every object to the unique object and sends labels by inclusion into $G_n$. The inclusion $j_n:\mathsf BG_n\to T_n$ selects object $n$. We have $r_nj_n=\id$ and a natural transformation $\id_{T_n}\Rightarrow j_nr_n$ whose components have identity labels. Hence $T_n$ is isomorphic in $\Cathot$ to a groupoid and is local.

We check the stronger fibrancy assertion separately. For a generating Thomason horn $A\hookrightarrow B$ and $F:A\to T_n$, first extend $r_nF:A\to\mathsf BG_n$ to $B$. Such an extension exists because groupoids are fibrant. Assign every new object of $B$ the object $n$ of $T_n$; keep the old objects and arrows exactly as in $F$. Assign to each arrow involving a new object the label given by the extended groupoid functor. The horn inclusion is a full sieve, so there is no arrow from a new object to an old object. Every remaining arrow with an old source and new target is permitted, because $\Hom_{T_n}(i,n)=G_n$. The label composition rule verifies all mixed composites. This is a strict filler, so $T_n$ is fibrant. E.3 makes $T$ fibrant.

Let $G=\bigcup_nG_n$ and define $s:T\to T$ by $i\mapsto i+1$ and by the shift permutation $s(g)$ which fixes $0$ and carries $g$ on $\{0,\ldots,j\}$ to its action on $\{1,\ldots,j+1\}$. The multiplication rule makes this a functor. For fixed $n$, put $m=n+1$. Both the inclusion $T_n\to T_m$ and the shifted functor $T_n\to T_m$ have natural transformations to the corresponding functors constant on objects at $m$, using identity labels. These last two functors are naturally isomorphic by the label $c=(0\,1\,\cdots\,n+1)$, because
$cgc^{-1}=s(g)$ for $g\in G_n$. Thus the inclusion and shift are connected by three natural transformations already in $\Fun(T_n,T_{n+1})$. Their induced ind-morphisms are equal.

The maps $r_n$ commute with the group inclusions. Filtered-colimit preservation of nerves and homotopy groups, as in E.3, gives $BT\simeq BG$: it is connected, its fundamental group is $G$, and its higher homotopy groups vanish. Under this identification $Bs$ induces the group shift, whose image fixes $0$ and omits $(0\,1)$. It is not surjective on the fundamental group, hence is not a weak equivalence or homotopic to the identity. In particular $T$ is not local. If it were local, the reflection in C.3 would identify its $\Cathot$ endomorphisms with the corresponding $\Hot$ endomorphisms, and E.2 applied to the finite group categories would not itself suffice; directly, the functor $\mathsf BG\to T$ required to represent $\id_{BG}$ would send its object to some finite stage $i$ and factor its entire group map through $G_i$, making a fundamental-group isomorphism impossible. The identity class is therefore missing. The code \texttt{e\_full\_filtration.py} checks every composition and all three transformations through stage three; the uniform formulas above prove the assertion at every stage and in the union.
\end{proof}

\begin{claim}{E.5}{Refuted}{V ms.~48: a local target does not force a strict retract}
Even when $A$ is local, $i:A\hookrightarrow B$ is a Thomason cofibration, and a functor $p:B\to A$ satisfies $[p][i]=\id_A$ in $\Cathot$, there need not exist any functor $r:B\to A$ with $ri=\id_A$ as an equality in $\Cat$.
\end{claim}
\begin{proof}
Let $A$ be the nine-object fence on $0,\ldots,8$, with an arrow from each even vertex to each adjacent odd vertex and no other nonidentity arrows. Let $B$ be the twelve-object circular fence on $\Z/12$, with the same parity rule and adjacency modulo twelve. Include the first nine vertices. This is $c\operatorname{Sd}^2$ applied to the inclusion of a two-edge path in the boundary of a triangle, so it is a Thomason cofibration. The category $A$ is categorically contractible: the successive functors $h_k(i)=\max(i,k)$, $0\leq k\leq8$, are order-preserving and adjacent ones are comparable pointwise. They join the identity to the constant functor at $8$. The constant functor $p:B\to A$ therefore satisfies the required equality in $\Cathot$, and $A$ is local by C.3.

Any functor of these fence posets takes adjacent vertices to equal or adjacent vertices, and hence cannot increase the undirected graph distance. A strict retraction would fix $0$ and $8$. In $B$ they are connected by the four-edge path through $11,10,9$, whereas in $A$ their distance is eight. This is impossible. The exhaustive check of all $9^3=729$ possible choices for the three new vertex images is in \texttt{e\_retraction\_and\_pointed.py}. This example concerns the local meaning of ``cat\'egorie de Thomason''; $A$ is not Thomason-fibrant. Indeed its connected $\operatorname{Ex}^2$-nerve has no edge from $0$ to $8$, because such an edge would be a four-edge-fence map. A connected Kan complex has an edge between every two vertices.
\end{proof}

\begin{claim}{E.6}{Open}{The additional fibrancy version of strict retraction}
For a Thomason-fibrant category $A$, a Thomason cofibration $i:A\hookrightarrow B$, and a functor $p:B\to A$ with $[p][i]=\id_A$ in $\Cathot$, this audit does not decide whether some strict retraction $r$ exists in the same $\Cathot$ class as $p$. The finite example in E.5 does not satisfy this extra fibrancy hypothesis.
\end{claim}
There is always a strict retraction in the same $\Hot$ class as $p$ under these fibrancy hypotheses; no cofibrancy of $A$ is needed for this weaker conclusion. To prove it, encode the specified finite zigzag by $A\times I\to A$, where $I$ is the free category on its finite oriented interval. The constant functor $c:A\to\Fun(I,A)$ is a weak equivalence by the finite interval contraction of E.20. Factor endpoint evaluation as a trivial cofibration $j:\Fun(I,A)\to P$ followed by a fibration $(d_0,d_1):P\to A^2$. Then $jc:A\to P$ is a weak equivalence, $d_0jc=d_1jc=\id_A$, and $A$ fibrant implies that both $d_i:P\to A$ are trivial fibrations. The zigzag supplies $h:A\to P$ with $(d_0h,d_1h)=(pi,\id_A)$. Lift $h$ along the cofibration $i$ against $d_0$, over the map $p:B\to A$. The lift $\widetilde h:B\to P$ gives $r=d_1\widetilde h$ with $ri=\id_A$. In $\Hot$, $d_0=d_1=(jc)^{-1}$, so $[r]=[p]$ there. Nothing in this factorization identifies that model path with a finite natural-transformation zigzag on $B$, which is the remaining requirement in E.6. If $A$ is an actual groupoid the stronger statement is elementary: the zigzag $pi\simeq\id_A$ composes to one natural isomorphism, and conjugating $p$ objectwise by that isomorphism on the full image of $i$, with identity choices elsewhere, gives a strict retraction naturally isomorphic to $p$. These are the verified reductions and the precise unresolved step.

\begin{claim}{E.7}{Proved}{V ms.~49--52: categorical homotopy and a canonical mapping object}
Stagewise natural-transformation homotopies do not imply a global one, even with full fibrant filtrations, by E.4. Without fibrancy the fence of C.5 realizes the manuscript's proposed unbounded-length mechanism. Fibrancy also does not make every categorical homotopy a single transformation. A canonical replacement for arbitrary relative-category mapping types is a derived mapping space of the simplicial localization, whose components are the localized Hom-sets; this construction has no right adjoint $\Hot\to\Cathot$ as an input.
\end{claim}
\begin{proof}
The first two assertions have the indicated explicit counterexamples. For the third take the terminal fork $0\to1\leftarrow2$. Every category with a terminal object is Thomason-fibrant: a functor from a full sieve extends by assigning the terminal object to all new objects and using its unique incoming arrows. This applies to every generating horn. The two functors $e\to C$ selecting $0$ and $2$ are connected by two transformations but by none of length one in either direction. The finite category is checked in \texttt{a\_finite\_checks.py}.

For the mapping object, let $(\mathcal M,W)$ be a relative category in a universe in which its simplicial localization is defined. The hammock localization $L^H(\mathcal M,W)$ has simplicial mapping objects $L^H(\mathcal M,W)(a,b)$ and composition, with
\[
 \pi_0 L^H(\mathcal M,W)(a,b)=\Hom_{\mathcal M[W^{-1}]}(a,b).
\]
This is \citet[\S2 and Proposition~3.1]{DwyerKan1980Calculating}. In a model category these objects compute the derived function complexes \citep[Proposition~4.4 and Corollary~4.7]{DwyerKan1980Function}. These theorems specify a different object from an unqualified ordinary union of strict zigzag categories. Agreement on $\pi_0$ does not give their weak equivalence. A comparison with the manuscript's exact all-arrows path model is proved separately in E.20 and following; for a general $W$ it requires an additional calculus or rectification hypothesis.
\end{proof}

\begin{claim}{E.8}{Proved}{VII ms.~108--124: total categories and actual filtered diagrams}
For an actual filtered diagram $D:I\to\Cat$, the canonical map from its Grothendieck construction to $\colim_I D$ is a weak equivalence. Objectwise weak equivalences of actual filtered diagrams induce weak equivalences of colimits. If $P$ has finite CW type, then
\[
 \colim_i[P,BD_i]\longrightarrow[P,B\colim_iD_i]
\]
is a bijection. These conclusions do not follow from a mere isomorphism of the diagrams as ind-objects in $\Cathot$ or $\Hot$, and the last assertion is false for a general nerve of a finite category.
\end{claim}
\begin{proof}
The nerve of the Grothendieck construction is naturally weakly equivalent to $\hocolim_I ND_i$ by \citet[Theorem~26.8 and Corollary~26.9]{ChacholskiScherer2001}. For a filtered category, homotopy colimit and ordinary colimit of simplicial sets are weakly equivalent: finite simplicial spheres and the finite data specifying their homotopies give the same filtered homotopy groups after a Kan replacement commuting with filtered colimits. The nerve commutes with this colimit by finite presentation of each $[n]$. These observations prove the first two assertions. Compactness of finite CW types proves the third, as in F.7. E.4 refutes the proposed implication from ind-isomorphisms even for full fibrant filtrations. F.1 gives the failure of the last assertion using $BC_2$ and finite skeleta of $BC_2$; for an actual categorical diagram, take $D_r=c\operatorname{Sd}^2\sk_rN(\mathsf BC_2)$. These are finite cofibrant posets, their transition maps are actual functors, and their classifying spaces and colimit are weakly equivalent to the corresponding skeleta and $BC_2$. The distinction invalidates the general inference at VII ms.~121 and its summary on ms.~123. It does not by itself refute a particular geometric comparison for which an independent proof is available.
\end{proof}

\begin{claim}{E.9}{Proved}{External Hom in the rectified setting, with its size condition}
Suppose a relative category $(\mathcal M,W)$ presents an $\infty$-category $\mathcal C$, and suppose, for every test category $X$ under consideration, its strict diagram localization presents the coherent diagram category:
\[
 \mathcal M^X[W_X^{-1}]\simeq h\Fun(NX,\mathcal C),
\]
naturally and compatibly with constant diagrams. Put $Z_{a,b}=\Map_{\mathcal C}(a,b)$. Then
\[
 \Hom_{\mathcal M^X[W_X^{-1}]}(a_X,b_X)
 \cong[BX,Z_{a,b}].
\]
The derived mapping object $Z_{a,b}$, not a Thomason--Kan envelope, is the correct representing object on $\Hot$. If $\mathcal M$ is $\U$-small and the manuscript's countable path ind-representation holds on all $\Cat_\U$, then this additional rectification assertion on all those shapes forces each $Z_{a,b}$ to be a 1-type.
\end{claim}
\begin{proof}
A coherent natural transformation between the constant diagrams at $a,b$ is a coherent $NX$-indexed family of points in $Z_{a,b}$. The end formula for the mapping space, or equivalently the limit of the constant space-valued diagram, gives
\[
 \Map_{\Fun(NX,\mathcal C)}(a_X,b_X)
 \simeq\Map_{\Spc}(BX,Z_{a,b}).
\]
Taking components and applying the stated presentation equivalence proves the formula. This also follows directly, in simplicial sets, from the projective bar resolution $N(X/-)$ used in D.1. A standard sufficient framework for diagram rectification is a combinatorial simplicial model category with its projective diagram structure; the precise coherent-nerve comparison is \citet[Proposition~4.2.4.4]{Lurie2009}. Such a category of all $\U$-small models is generally $\U$-large; one forms it in $\V$. This size distinction is material here.

For a $\U$-small $\mathcal M$, its fixed-length path categories and their countable index form a $\U$-small ind-object. If the displayed Hom-functor equals $[B-,Z_{a,b}]$ on all $\U$-small shapes, D.8 forbids any higher homotopy group of $Z_{a,b}$. Thus the formula for large model presentations, or for bounded shapes, cannot be transported to a small presentation with all shapes without checking this condition. This is also why an arbitrary choice of $W$ does not imply diagram rectification.
\end{proof}

\begin{claim}{E.10}{Proved}{VII ms.~128--136: strict arrows, strict diagrams, and path types}
For every category $C$, the canonical map
\[
 B\Fun([1],C)\longrightarrow\Map(B[1],BC)
\]
is a weak equivalence. Its strict fibre at $(a,b)$ need not be a homotopy fibre. For an arbitrary source $Y$, replacing $C$ by a Thomason-fibrant category does not make the corresponding comparison a weak equivalence, and it need not make $\Fun(Y,C)$ fibrant. A corrected mapping-space formula is $\Map_{\mathbf{sSet}}(NY,RNC)$, or the cosimplicial formula of A.5.
\end{claim}
\begin{proof}
Let $e_0:\Fun([1],C)\to C$ evaluate at the source, and let $c:C\to\Fun([1],C)$ select identity arrows. Then $e_0c=\id_C$ and the square $(\id_x,f)$ at an arrow $f:x\to y$ gives $ce_0\Rightarrow\id$. Thus $Be_0$ is a homotopy equivalence. Evaluation $\Map(B[1],BC)\to BC$ is likewise a homotopy equivalence by contraction of the interval. The canonical comparison commutes with these evaluations, proving the assertion by two-out-of-three. Its strict endpoint fibre is the discrete set $\Hom_C(a,b)$, whereas its derived endpoint fibre is the full path space in $BC$. For the terminal fork, the fibre from $0$ to $2$ is empty although that derived path space is contractible and nonempty. This corrects the absolute assertion at VII ms.~135 while preserving the reason to replace strict endpoint fibres. The two general fibrancy assertions are independently refuted in A.6. No claim that products with arbitrary categories preserve Thomason cofibrations is available: multiplying $\varnothing\to e$ by every $Y$ would already make every category cofibrant.

Grothendieck's statement that the path categories ``\emph{n'ont aucune tendance \`a \^etre des cat\'egories de Thomason-Kan}'' (VII ms.~136) has a precise reading. In the local sense a path category modeling $\Omega Y$ cannot be local whenever $\Omega Y$ has a nonzero homotopy group above degree one. For example paths at a point of a category of type $S^3$ model $\Omega S^3$, whose second homotopy group is $\Z$. No weakly equivalent local category exists by C.3. Fibrancy is a separate property and can be imposed for derived computations; it cannot change that obstruction.
\end{proof}

\begin{claim}{E.11}{Proved}{VII ms.~127, 150--156 and IX ms.~13: what the set-valued arguments prove}
An equality or an inverse verified only in $\Hot_0^{\wedge}$ cannot in general be pulled back to an equality or an inverse in $\Hot$. Thus the finite-test arguments for associativity, $T^2=1$, and weak equivalence of the displayed fibre-transition maps do not prove those conclusions. They are valid statements at the presheaf or ind-object level at which they were established. Fully coherent space-valued finite tests recover the conclusions, and height-two tests detect a specified weak equivalence. Equality of specified maps requires the stronger conditions detailed in Part F. For the actual all-arrows path construction, E.20 and following give direct proofs independent of these invalid inferences.
\end{claim}
\begin{proof}
F.2 gives a specified noninvertible map whose finite-test image is the identity. F.6 gives two unequal maps whose images agree, even with an explicit failure of the proposed epimorphism from an ind-object to its telescope. F.7--F.9 give the positive detection assertions and F.13 types the equality applications. In particular, the argument at VII ms.~152 and 156 only produces an isomorphism of restricted presheaves; conservativity at that level is false. The assertion that $T^2$ acts trivially on an ind-object is likewise insufficient by E.4. None of these counterexamples alone says that the manuscript's particular geometric operation is nonassociative, or that its particular reversal fails to square to one.
\end{proof}

\begin{claim}{E.12}{Proved}{IX ms.~9--13: product-preserving passage and the second psi}
The functors $\varphi$, their ind-extensions, restricted Yoneda $R_1^{\mathrm{cat}}$, and filtered-colimit realization in the presheaf category preserve finite products. They preserve the component sets obtained by evaluation at $e$. Therefore composition already specified as a morphism in one of these categories passes forward and gives the same ordinary category on components. The functor $R_1^{\mathrm{cat}}$, called $\psi$ on this page, is neither faithful, full, nor conservative. Its ind-extension is not faithful or full, since these failures already occur on constant ind-objects.
\end{claim}
\begin{proof}
Finite products of functor categories commute with $\pi_0$: two finite zigzags can be padded by identities to the same alternating length and performed coordinatewise. Realization of nerves preserves finite products up to the natural weak equivalence. Products in an ind-category are represented by pairwise products over the product of the two filtered index categories. Restricted Yoneda preserves products pointwise. Filtered colimits of sets commute with finite products, giving the last preservation assertion for presheaves. Evaluation at $e$ is $\pi_0$ at every level, so ordinary composition, its identities, and its associativity survive these passages. The three failures are F.2--F.3. Since the embedding of a category into its ind-category is fully faithful, the same witnesses between constant ind-objects give the claimed failures for the ind-extension. These facts validate the forward diagram at IX ms.~9 and invalidate the hopes of recovering the higher composition data by reversing its arrows.
\end{proof}

\begin{claim}{E.13}{Refuted}{IX ms.~10: nonconstant ind-objects can have ordinary colimits in Hot}
The assertion that an ind-object of $\Hot$ can have a colimit there only when it is essentially constant is false. The finite discrete spaces $D_n=\{0,\ldots,n\}$ have the discrete countable space $\N$ as colimit in $\Hot$, although $(D_n)$ is not essentially constant in $\Ind(\Hot)$.
\end{claim}
\begin{proof}
For every space $Y$, a map from a discrete space is a choice of a component for each point. Thus
\[
 [\N,Y]=(\pi_0Y)^{\N}=\lim_n(\pi_0Y)^{D_n}=\lim_n[D_n,Y],
\]
naturally, which is the colimit universal property. If the ind-object were represented by a constant $Z$, evaluating its presheaf at $e$ would give $\pi_0Z\cong\N$. The inverse ind-isomorphism $Z\to(D_n)$ would be represented at some finite stage, and composing with the supposed isomorphism back to $Z$ would factor $\id_Z$ in $\Hot$ through $D_n$. This is impossible on components. The existence or nonexistence of a colimit in a homotopy category must therefore be checked for the particular diagram; a homotopy colimit is not automatically that ordinary colimit.
\end{proof}

\begin{claim}{E.14}{Refuted}{IX ms.~22--23: the written tests permit a finite nontrivial relative category}
There is a finite pointed relative category $(\mathcal E,W)$, whose localization is not equivalent to a point, satisfying the representability tests for $\Sigma,\Omega$, the weak-equivalence invariance of the constant-diagram external Hom, and the displayed higher loop-mapping formulas, when those tests are interpreted exactly as written before the finite-limit hypothesis of IX ms.~24. A finite pointed category with all finite products, by contrast, is necessarily equivalent to the terminal category.
\end{claim}
\begin{proof}
Take the full category of pointed sets on the singleton $0$ and the two-element pointed set $x$, and let $W$ contain only identities. The category has five arrows. In particular $\End(x)=\{\id_x,z\}$ with $z^2=z\ne\id_x$; all other Hom-sets are singletons. Thus $0$ is a zero object and the localization is the same nonterminal category. For any small $X$,
\[
 \Hom_{\mathcal E^X}(a_X,b_X)=\Hom_{\mathcal E}(a,b)^{\pi_0X},
\]
since naturality forces the component arrow to be constant on every zigzag component. This inverts every $\W_\infty$-map. The coherent mapping spaces for the ordinary category are discrete, and the loop space of every pointed discrete space is a point. Consequently $\Sigma a=0=\Omega a$ gives both stated representabilities and both mapping-space formulas, with their zero basepoints.

For the actual identity-backward path categories, compose the finite forward chain to define $q:\mathrm{Ch}^W_\infty(a,b)\to\Hom_{\mathcal E}(a,b)_{\mathrm{disc}}$. Insert the composite at a fixed forward edge to define $s$. Replace the part before that edge by $a$, using its accumulated composite into the unchanged tail, to define $R$. Accumulated initial and final composites give natural transformations $R\to\id$ and $R\to sq$, and $qs=\id$. The backward arrows are identities, so all duplicated vertices have equal components and every square commutes by associativity. Thus these exact path categories are isomorphic in $\Cathot$ to discrete categories. The multiplication table and all typed associativity equations of $\mathcal E$ are verified in \texttt{e\_retraction\_and\_pointed.py}.

For the final assertion, if a finite category has two distinct arrows $a\to b$ and all finite powers of $b$, then it has at least $2^r$ arrows $a\to b^r$ for each $r$. This contradicts finiteness of its total arrow set. It is therefore a preorder. In a pointed category every Hom-set is nonempty because of the zero map, so every Hom-set is a singleton and all objects are isomorphic. The extra finite-product hypothesis, rather than tests (a)--(c), supplies this obstruction.
\end{proof}

\begin{claim}{E.15}{Proved}{IX ms.~22--25: representability versus derived loops}
In a pointed $\infty$-category $\mathcal C$ having the required finite limits and colimits, define $\Omega b=0\times_b^h0$ and $\Sigma a=0\amalg_a^h0$. There are canonical pointed equivalences
\[
 \Map_{\mathcal C}(a,\Omega b)\simeq
 \Omega\Map_{\mathcal C}(a,b)\simeq
 \Map_{\mathcal C}(\Sigma a,b),
\]
compatible with loop multiplication. They give $\Sigma\dashv\Omega$ on the homotopy category. In a pointed simplicial model category they can be calculated by replacing $b$ by a fibrant object and factoring its diagonal as
\[
 b\xrightarrow{\simeq}P b\xrightarrow{(d_0,d_1)}b\times b,
\]
with the second arrow a fibration, then taking the fibre at the zero pair. The first arrow must be a weak equivalence; merely calling it a cofibration is insufficient.
\end{claim}
\begin{proof}
The universal property of a homotopy pullback identifies maps into $0\times_b^h0$ with the homotopy pullback $*\times_{\Map(a,b)}^h*$, which is the based loop space. The dual pushout gives the second equivalence. The definitions use the same loop concatenation, so the comparison preserves multiplication and all its coherences. Taking components gives the two representabilities and their adjunction by Yoneda; an arbitrary component-level representability would not conversely prove these space-level equivalences.

For the model calculation take a cofibrant source $a$. The simplicial mapping functor $\Map(a,-)$ preserves the fibration and the defining strict pullback, and carries the weak equivalence between fibrant objects to a weak equivalence. It therefore supplies the corresponding path-object calculation of the loop space. If the first arrow is not required to be a weak equivalence, take a nontrivial Kan $b=K(C_2,1)$, set $P b=b\times b$, and take the second arrow to be the identity. The diagonal is a cofibration of simplicial sets and all stated fibrancy conditions hold, but its strict fibre at the zero pair is a point, whereas $\Omega b$ has two components. This is the exact missing hypothesis in the proposed factorization at IX ms.~24--25.
\end{proof}

\begin{claim}{E.16}{Proved}{The scope of the page corrections}
Every use of the envelope or of the finite-test inference in the specified ranges is assigned a replacement in the following table. A reference to a proved conditional theorem retains its displayed hypotheses; it does not assert that the manuscript has already established them.
\end{claim}
\begin{center}\small
\begin{longtable}{L{0.17\linewidth}L{0.43\linewidth}L{0.31\linewidth}}
\toprule
Manuscript location & Assertion or use & Correct statement and proof\\\midrule
\endfirsthead
\toprule Manuscript location&Assertion or use&Correct statement and proof\\\midrule\endhead
V 37 & Right adjoint asserted; editorial request for a proof & False, B.4 and B.11; positive range C.2.\\
V 37--40 & Reflective-localization Proposition 1 and representability corollary & A.3 under its actual adjunction hypothesis; one-way descent A.7 without it.\\
V 40--42 & Ind-adjunction, Proposition 2, units on ordinary objects & A.4; the ind-unit is tested by the extended functor.\\
V 42--44 & Envelopes and compatible transitions & E.1 and E.3; bounded targets D.1.\\
V 44--46 & External Hom ind-object versus ordinary union & A.7, E.2, E.9; the union need not be local.\\
V 46 & Invariant small ind-objects have local-stage presentations & Exact remaining implication and positive criterion E.17.\\
V 46--47 & Local stages and weak equivalence of union maps & E.4 refutes the inference from local fibrant stages; genuinely local unions satisfy E.18.\\
V 48 & Rectification of a categorical retraction & Local version false, E.5; extra-fibrancy version E.6.\\
V 48--49 & Stagewise homotopies and uniqueness of union maps & Fibrancy alone does not suffice, E.4; genuinely local unions satisfy E.18.\\
V 49 & Shortening a homotopy to one arrow in a fibrant target & False, terminal fork in E.7.\\
V 50--52 & Canonical external homotopy type; proposed return through $\psi$ & E.7 and E.9; no global return through an adjoint.\\
V 52 & Finite tests faithful, full, or conservative & All false for set-valued tests, F.2--F.4.\\
VII 108--115 & Total category, telescope, and all-type projection & Ordinary-colimit leg E.8; exact general all-type projection retained in E.26.\\
VII 116--120 & Independence of arbitrary cofinal systems and transitive comparisons & Canonical collar comparisons E.23 and E.25; the fully quantified original comparisons remain E.26.\\
VII 120--124 & Finite categories compact; conservativity; canonical union & General compactness and detection steps false, F.1--F.2 and E.4; correct compact tests F.7.\\
VII 125--127 & Chosen concatenation, arbitrary blocks, and associativity & Canonical all-arrows operation E.24; relative units E.25; remaining general block and relative assertions E.27.\\
VII 128--130 & Pass to an envelope; reduce to direct one-step paths & E.1, E.9--E.10; strict endpoint arrows do not model general path fibres.\\
VII 131--136 & Strict-diagram and mapping-category arguments & Require diagram rectification E.9; A.6 refutes both fibrancy shortcuts.\\
VII 135 & Absolute arrow-category comparison & True for $[1]$, E.10; the fibre inference is false.\\
VII 136 & Path categories as Thomason--Kan categories & Localness means a 1-type, C.3 and E.10; fibrancy is different.\\
VII 150 & Infer $T^2=1$ by faithfulness & Inference false, E.11; the actual translation satisfies E.22 and its relative extension E.25.\\
VII 150--152 & Endpoint transport inferred from restricted presheaves & Inference false, E.11; direct all-arrows lift-category proof E.21.\\
VII 153--156 & Mapping-path transport and slice transition maps & In the all-arrows scope declared at ms.~137, use E.21 and Quillen B; set-valued conservativity is unavailable.\\
IX 9--10 & Product diagram and the new meaning of $\psi$ & E.12; forward passage and components valid, reverse recovery false.\\
IX 10 & Colimit possible only for essentially constant ind-objects & False, E.13.\\
IX 11--13 & Variance from composition; associativity at ind level and at union & Ind-level Yoneda A.4; E.11, the canonical proof E.24, and the general scope E.27.\\
IX 22--23 & Local path categories identified with invariance; loops and suspension & E.9 and E.15; neither interpretation of ``Thomason--Kan'' can be silently substituted.\\
IX 23 & Proposed finite-category obstruction & Tests alone insufficient, E.14; finite products give the actual obstruction.\\
IX 24--25 & Path-object fibre computes loops & Require the weak-equivalence and fibrancy hypotheses in E.15.\\
\bottomrule
\end{longtable}
\end{center}

\subsection{The remaining assertion about local ind-presentations}

\begin{claim}{E.17}{Open}{V ms.~44 and 46: local presentations of invariant ind-objects}
The following sequential assertion remains unresolved in this audit.
Let $K=(K_n)_{n\in\N}$ be a sequence in $\Cathot_\U$, and assume
that its associated presheaf
\[
 H_K(X)=\colim_n\Hom_{\Cathot}(X,K_n)
\]
sends every map in $\W_\infty$ between $\U$-small categories to a
bijection. Must $K$ be isomorphic in $\Ind_\U(\Cathot_\U)$ to a
sequence of local categories, equivalently to a sequence of groupoids?
The corresponding assertion with an arbitrary $\U$-small filtered
indexing category in place of $\N$ also remains unresolved here.
We do not infer either assertion from descent in A.7, from the
representable-object classification in Part C, or from the obstruction
in D.8 for a single represented homotopy type.
\end{claim}

There is, however, a proved exact criterion for both versions. Let
$\mathcal G\subset\Cathot_\U$ be the full subcategory on groupoids,
let $j:\mathcal G\hookrightarrow\Cathot_\U$ be its inclusion, and
let $\eta_X:X\to j\Pi_1X$ be the unit of the reflection in C.4.
For any $\U$-small filtered $I$ and $K:I\to\Cathot_\U$, the
following conditions are equivalent:
\[
\begin{gathered}
 H_K(\eta_X)\text{ is bijective for every }X\in\Cat_\U;\\
 K\longrightarrow(\Ind_\U j)(\Ind_\U\Pi_1)K
       \text{ is an isomorphism};\\
 K\text{ is isomorphic to a }\U\text{-small ind-object of groupoids}.
\end{gathered}
\tag{E.17a}
\]
These conditions are also equivalent to invariance under $\W_1$,
the functors inducing bijections on components and isomorphisms on
fundamental groups. When they hold, the groupoid presentation can
be taken to be the canonical diagram $(\Pi_1K_i)_{i\in I}$ on the
\emph{same} indexing category. In particular the sequential assertion
does not require an unproved change from a general filtered index
back to a sequential one.

\begin{proof}[Proof of the criterion and the verified special cases]
C.4 gives the fully faithful right adjoint $j$ to $\Pi_1$ and
identifies its inverted class with $\W_1$. Apply A.4 to this actual
adjunction. It gives (E.17a), its $\W_1$ formulation, and the
specified unit. A diagram of local categories is equivalent to a
diagram of groupoids by applying its objectwise invertible units,
using C.3. The canonical comparison on presheaves is explicitly
\[
 \colim_i\Hom_{\Cathot}(X,K_i)
 \longrightarrow
 \colim_i\Hom_{\mathcal G}(\Pi_1X,\Pi_1K_i).
 \tag{E.17b}
\]

If the initial hypothesis of E.17 holds and the $BK_i$ are
$1$-types on a cofinal subdiagram, A.7 identifies the left side
of (E.17b) with $\colim_i[BX,BK_i]$. Each $B\eta_{K_i}$ on
that subdiagram is a weak equivalence, and groupoid locality
identifies the right side with the same colimit. Thus (E.17b)
is bijective and the conclusion holds in this case. Essentially
constant ind-objects with invariant presheaves are also covered
directly by A.7 and C.3, after replacing the diagram by a constant
presentation. More generally, if the descended presheaf on $\Hot_\U$ is
representable by a space $Y$, D.8 forces $Y$ to be a $1$-type.
Then $H_K$ is represented by its fundamental groupoid, so the
ind-object is itself essentially constant with a local representative.
\end{proof}

The missing implication is therefore precisely
\[
 \boxed{\quad
 H_K\text{ inverts }\W_\infty
 \quad\Longrightarrow\quad
 H_K\text{ inverts every }\eta_X:X\to\Pi_1X.
 \quad}
 \tag{E.17c}
\]
The reverse implication follows from (E.17a) because
$\W_\infty\subset\W_1$. The maps $\eta_X$ in (E.17c) are
not weak homotopy equivalences when $BX$ has higher homotopy.
Consequently A.4 applied to the genuine groupoid reflection does
not prove (E.17c) from its stated premise. A.7 only proves
$H_K'\cong\colim_i[-,BK_i]$ after descent; the target varies
with $i$, so D.8 cannot be applied to one of the stages as if it
represented that whole descended presheaf.

The size restriction is material. If the index is allowed to be
$\V$-small and unbounded in the $\U$-cardinals, D.7 gives an
invariant ind-presentation of $X\mapsto[BX,K(\Z/2,2)]$.
It cannot be isomorphic to any ind-object of groupoids: every
such ind-object sends the projection from a finite poset of type
$S^2$ to $e$ to a bijection, whereas this functor takes that map
to the nonbijective constant-map function
$\{*\}\to H^2(S^2;\Z/2)=\Z/2$. This is a counterexample to
the version with the enlarged index universe, and does not answer
the $\U$-small or sequential assertions in E.17. Likewise E.4
refutes the later inference about unions even for a diagram already
made of local categories; it does not refute the existence question
posed in E.17.

\begin{claim}{E.18}{Proved}{V ms.~47--49 with genuinely local unions}
Let $(K_n)$ and $(K'_n)$ be sequences of full inclusions of $\U$-small categories, with unions $K,K'$. Suppose that the unions themselves are local. Let $u:K\to K'$ be an actual functor carrying each $K_n$ into some $K'_{m(n)}$. If the induced morphism in $\Ind_\U(\Cathot_\U)$ is invertible, then $[u]$ is invertible in $\Cathot$. If two such functors $u,v:K\rightrightarrows K'$ induce the same ind-morphism, then $[u]=[v]$ in $\Cathot$. Localness of the unions is a separate hypothesis here; it does not follow from localness of the stages.
\end{claim}
\begin{proof}
For every $\U$-small group $H$, its one-object group category $\mathsf BH$ has finitely many objects, regardless of the size of $H$. E.2 therefore identifies
\[
 \Hom_{\Cathot}(\mathsf BH,K)
       =\colim_n\Hom_{\Cathot}(\mathsf BH,K_n),
\]
and gives the analogous identity for $K'$. These are the actual maps induced by the inclusions. Localness identifies their left sides with $[BH,BK]$ and $[BH,BK']$. An ind-isomorphism induced by $u$ thus gives bijections on these sets for every $H$, and the test $e$ gives a bijection on components.

Both $BK$ and $BK'$ are 1-types by C.3. The preceding bijections respect their decompositions by target component, because restriction along $e\to\mathsf BH$ is evaluation in that component. Every connected 1-type is equivalent to some $BH$. On each corresponding pair of components, ordinary Yoneda on the category of connected 1-types therefore identifies $Bu$ as an isomorphism. These componentwise inverses assemble over the small set of components to an inverse in $\Hot$. By C.3, the full subcategory of local categories maps fully faithfully to the category of 1-types, so that inverse lifts to an inverse of $[u]$ in $\Cathot$.

For the equality assertion, choose an isomorphism in $\Cathot$ from a groupoid of the form $\coprod_{c\in\pi_0K}\mathsf BG_c$ to $K$. Testing each of its component functors with E.2 shows that the two composites with $u,v$ are equal in $\Cathot$ on that component. Compose into the actual groupoid $\Pi_1K'$. A zigzag of transformations between groupoid-valued functors composes, after inverting arrows as needed, to one natural isomorphism. Choose such an isomorphism on every component and assemble them into a single natural isomorphism on the disjoint union. Cancelling the two local groupoidification isomorphisms proves $[u]=[v]$. This explicit assembly deals with infinitely many components and uses no uniform bound on lengths of the original zigzags. The counterexample in E.4 violates the genuine-local-union hypothesis and hence does not contradict this statement.
\end{proof}

\begin{claim}{E.19}{Known}{A coherent replacement for arbitrary relative path composition}
Let $\mathcal M$ be $\U$-small and let $W$ be a subcategory containing every object. The hammock localization, followed on each mapping simplicial set by $\operatorname{Ex}^{\infty}$, gives $\U$-small Kan complexes
\[
 \mathcal H_W(a,b)=\operatorname{Ex}^{\infty}
          L^H(\mathcal M,W)(a,b)
 \qquad(a,b\in\mathcal M)
\]
with strictly associative, strictly unital simplicial composition. Their component category is canonically $\mathcal M[W^{-1}]$. Precomposition or postcomposition by any specified arrow in $W$ is a weak equivalence of these mapping complexes. These are coherent mapping objects for the relative localization. No comparison with a specified ordinary colimit of the manuscript's path categories is part of this assertion.
\end{claim}
\begin{proof}
The hammock simplices, their faces, and their composition are those of \citet[Definition~2.1]{DwyerKan1980Calculating}: a simplex is a finite commuting hammock, its leftward and vertical arrows lie in $W$, and composition concatenates horizontally and reduces consecutive same-direction columns and identity columns. The resulting simplicial category has strictly associative composition and identity hammocks. Every set of finite hammocks lies in $\U$, and taking the countable union of lengths and simplicial degrees preserves that size. Proposition~3.1 of the same paper identifies its component category with the ordinary localization. Proposition~3.3 proves the weak-equivalence assertion for composition with an arrow of $W$.

The functor $\operatorname{Ex}$ preserves products because it is a right adjoint; filtered colimits of sets commute with finite products, so its sequential iteration $\operatorname{Ex}^{\infty}$ also preserves finite products. Applying it to the composition maps therefore gives
\[
 \mathcal H_W(b,c)\times\mathcal H_W(a,b)
       \longrightarrow\mathcal H_W(a,c)
\]
with the same strict associativity and unit equations. The natural maps from a simplicial set to its $\operatorname{Ex}^{\infty}$ replacement are weak equivalences and the replacements are Kan. Thus the component category and the specified transport equivalences are unchanged. In a model-category presentation these are the derived function complexes by \citet[Proposition~4.4 and Corollary~4.7]{DwyerKan1980Function}. This supplies an unconditional relative-category construction in the stated small universe, while comparison with any other raw path model remains a separate mathematical assertion, as recorded in E.26--E.27.
\end{proof}

\subsection{An exact comparison for the unrestricted alternating-path model}

In E.20--E.24 we take $\Sigma=\Mor(X)$. We give the comparison for the precise paths and maps under discussion, rather than replace them silently by another categorical cocylinder. Let $J$ be the integer fence: its objects are $j\in\Z$, its nonidentity arrows are
\[
                 2k\longrightarrow 2k-1,
                 \qquad 2k\longrightarrow 2k+1,
\]
and there are no composable pairs of nonidentity arrows. Put $\mathcal P(X)\subseteq\Fun(J,X)$ for the full subcategory of functors which are constant, with identity arrows, outside some finite interval. A natural transformation between two such functors is necessarily constant on each sufficiently distant tail. Its two eventual components define the endpoint functor
\[
 p=(\sigma,\beta):\mathcal P(X)\longrightarrow X\times X.
\]
Thus $\mathcal P(X)$ is the manuscript's two-sided $\Ch_\infty(X)$ with this explicit parity convention. Its strict fibre $\mathcal P(X;a,b)$ has transformations whose two eventual components are $1_a,1_b$. We write $B$ for classifying space and $K=BX$. All homotopy categories of spaces over a base below mean the localization of the \emph{slice category} at underlying weak equivalences, not the slice of the ordinary homotopy category.

\begin{claim}{E.20}{Proved}{Finite support and the constant-path equivalence}
The constant-path section $j:X\to\mathcal P(X)$ is a weak equivalence. For each $a\in X$, both one-endpoint categories $\mathcal P(X;a,?)$ and $\mathcal P(X;?,a)$ are weakly contractible.
\end{claim}
\begin{proof}
For integers $n\leq m$, let $\mathcal P_{n,m}(X)$ be the full subcategory of paths constant outside $[n,m]$. Restriction identifies it with $\Fun(J|_{[n,m]},X)$. Extension by identities gives fully faithful transition functors when the interval increases. These categories form a filtered union on objects and arrows; every finite diagram is contained in one stage. The same statements hold with either or both endpoints fixed.

A finite oriented interval $I$ with vertices $0,\ldots,r$ has functors $r_k:I\to I$, $r_k(i)=\min(i,k)$, for $0\leq k\leq r$. They preserve each directed generating edge, sending an edge beyond $k$ to an identity. The successive functors $r_k,r_{k+1}$ are joined by a natural transformation in the direction of the edge between $k$ and $k+1$: its components on the tail are that edge, and its other components are identities. On every generating edge the naturality square commutes. Precomposition gives a finite zigzag from the constant-at-initial-value endofunctor of $\Fun(I,X)$ to its identity. This zigzag fixes the initial endpoint. Consequently the constant-diagram functor $X\to\Fun(I,X)$ is a weak equivalence; with initial endpoint fixed at $a$ the functor category is contractible. Reversing the enumeration proves the version with the terminal endpoint fixed.

Nerves preserve filtered colimits of categories, since the data and equations of a composable finite string occur at a stage. Weak equivalences of simplicial sets are preserved by filtered colimits. For completeness, $\operatorname{Ex}^{\infty}$ commutes with filtered colimits: its individual simplex formulas use finite subdivided simplices, and its sequential iteration also commutes with such colimits. In the resulting Kan complexes, homotopy groups and components are detected by finitely many simplices and finite homotopies. Applying this observation to the stagewise constant-diagram equivalences proves the first assertion and the two assertions with one endpoint fixed. No contraction of a strict two-endpoint fibre has been claimed.
\end{proof}

\begin{claim}{E.21}{Proved}{The exact endpoint fibres are homotopy fibres}
For each $a,b\in X$, the natural square
\[
\begin{array}{ccc}
 B\mathcal P(X;a,b)&\longrightarrow&B\mathcal P(X)\\
 \downarrow&&\downarrow Bp\\
 \{(a,b)\}&\longrightarrow&K\times K
\end{array}
\]
is homotopy Cartesian. In particular $B\mathcal P(X;a,b)$ has the weak homotopy type of the ordinary path space $P_{a,b}K$. The comparison with the homotopy fibre is natural for functors over $X\times X$.
\end{claim}
\begin{proof}
We first establish the two lifting properties for this exact category. If $n<m$ are even, both endpoints of $J|_{[n,m]}$ are sources. The endpoint functor on $\Fun(J|_{[n,m]},X)$ is a Grothendieck fibration: to lift $(u,v):(a',b')\to(a,b)$ into a diagram with endpoints $(a,b)$, replace the endpoint objects by $a',b'$ and precompose its two outer arrows by $u,v$. Leave every interior object and arrow unchanged. The transformation to the original diagram has components $u,v$ at the endpoints and identities inside. Any transformation over $(u,v)$ factors uniquely through it, so it is Cartesian. Positive length is essential here. If $n<m$ are odd, both endpoints are sinks, and postcomposition supplies the dual coCartesian construction.

For a general functor $q:E\to B_0$, an object $e\in E$, and an arrow $h:b\to q(e)$, write $L_q(e,h)$ for the category of pairs $(d,\gamma)$ with $q(d)=b$, $q(\gamma)=h$, and $\gamma:d\to e$; its arrows are vertical maps commuting with $\gamma$. For the functor $p$ above, this is the filtered union of the lift categories at sufficiently large positive even-even support intervals. Each has a terminal Cartesian lift, and hence is contractible. Every object, arrow, and equation in a finite diagram of lifts occurs at one such stage. The filtered-colimit argument of E.20 shows that $L_p(e,h)$ is weakly contractible. Odd-odd intervals prove, dually, that every outgoing-lift category is weakly contractible. Both parity families are cofinal among all finite supports. This proves the usual smoothness and properness lift tests directly; it does not assert that the union itself is a Grothendieck fibration or opfibration.

We spell out why these lift tests imply the required Quillen--B condition. Let $q:E\to B_0$ have both properties just proved, and let $E_b$ be its strict fibre. The inclusion
\[
                 i_b:E_b\longrightarrow b\backslash q
\]
is a weak equivalence by Quillen A: its over-comma category at $(e,h:b\to q(e))$ is exactly $L_q(e,h)$. For an arrow $h:b\to b'$, let $M_h$ have as objects arrows $\gamma:d\to d'$ of $E$ over $h$, and as arrows commuting squares whose two vertical sides lie over identities. Let $s:M_h\to E_b$ and $t:M_h\to E_{b'}$ be the source and target projections.

For $d_0\in E_b$, an object of $d_0\backslash s$ is an arrow $\gamma:d\to d'$ over $h$, together with a vertical $\alpha:d_0\to d$. Composition sends it to the outgoing lift $\gamma\alpha:d_0\to d'$. The inclusion of that outgoing-lift category takes $\alpha=1_{d_0}$. The composite in one direction is the identity; in the other, the square $(\alpha,1_{d'})$ gives a natural transformation from the composite to the identity. Thus $d_0\backslash s$ is homotopy equivalent to the contractible outgoing-lift category. Quillen A makes $s$ a weak equivalence. An object of $t/d'_0$ instead consists of $\gamma:d\to d'$ and a vertical $\alpha:d'\to d'_0$. The functor $(\gamma,\alpha)\mapsto\alpha\gamma$ retracts this category onto the incoming lifts into $d'_0$. Its inclusion takes $\alpha=1_{d'_0}$, and the squares $(1_d,\alpha)$ give the natural transformation from the identity to the other composite. Thus $t/d'_0$ is contractible and Quillen A makes $t$ a weak equivalence.

The arrows $\gamma$ themselves give a natural transformation
\[
             i_b s\ \Longrightarrow\ h^*i_{b'}t
             \quad\hbox{into }b\backslash q,
\]
where $h^*:b'\backslash q\to b\backslash q$ is precomposition. As $i_b,s,i_{b'},t$ are weak equivalences, this natural transformation and two-out-of-three show that $h^*$ is a weak equivalence. Quillen B therefore identifies $B(b\backslash q)$ with the homotopy fibre of $Bq$ over $b$, and the natural map from $BE_b$ to that homotopy fibre is a weak equivalence \citep[\S1, Theorems A and B]{Quillen1973}. This formulation is also recorded in \citet[Theorem 1.8]{Shoikhet2016}. The comma constructions and the maps to a standard homotopy-fibre model commute with functors over the base, which gives the asserted naturality.

Apply this argument to $p$. E.20 identifies the constant-path section as a weak equivalence over $K^2$, from the diagonal $K\to K^2$ to $Bp$. Consequently their derived fibres agree. The homotopy fibre of the diagonal over $(a,b)$ is the path space from $a$ to $b$. Equivalently, let $P(K)=K^{[0,1]}$ with its endpoint fibration and constant-path section $s:K\to P(K)$. The comparison is the isomorphism
\[
 u_{\mathcal P}=s\,(Bj)^{-1}:B\mathcal P(X)\longrightarrow P(K)
       \quad\hbox{in }\operatorname{Ho}(\mathrm{Spc}/K^2),
\]
followed by the derived-fibre functor. The preceding proof identifies its source derived fibre with the stated strict categorical fibre.
\end{proof}

\begin{claim}{E.22}{Proved}{The actual even translation is the identity in the homotopy category}
Let $T^2:\mathcal P(X)\to\mathcal P(X)$ be given by $(T^2c)(j)=c(j-2)$, with the corresponding shift on arrows and on natural transformations. It induces the identity in $\Hot$ on the total category and on every fixed-endpoint category.
\end{claim}
\begin{proof}
Translation by two is an automorphism of the integer fence, preserves eventual endpoint values, and fixes every constant diagram and constant natural transformation. Thus $T^2$ is an exact endofunctor over $X^2$ satisfying $T^2j=j$. The morphism $Bj$ is invertible in the localized slice $\operatorname{Ho}(\mathrm{Spc}/K^2)$ by E.20. Cancellation gives $B(T^2)=1$ there. Apply the derived-fibre functor at $(a,b)$. Naturality in E.21 identifies the resulting map with the classifying-space map of the \emph{actual restriction} of $T^2$ to $\mathcal P(X;a,b)$. Hence that map is the identity in $\Hot$. The argument neither assumes faithfulness of finite tests nor constructs a bounded zigzag of natural transformations from the two functors.
\end{proof}

\begin{claim}{E.23}{Proved}{The one-sided models and their actual comparison maps}
Put $D=(+,-)$. Let $\mathcal L(X)$ be the colimit of the categories of $D^n$-paths under padding by a constant $D$-block on the left, and let $\mathcal R(X)$ use padding on the right. They have constant-path sections and endpoint maps. Their constant-path sections are weak equivalences, their one-endpoint fibres are weakly contractible, and their strict two-endpoint fibres compute their derived endpoint fibres. Constant-tail extension gives functors
\[
 q_L:\mathcal L(X)\longrightarrow\mathcal P(X),
 \qquad q_R:\mathcal R(X)\longrightarrow\mathcal P(X)
\]
which are weak equivalences on total categories and on every fixed-endpoint fibre. These comparisons commute with the canonical path-space comparisons of E.21.
\end{claim}
\begin{proof}
Realize $\mathcal L(X)$ on the half-fence $j\leq0$, eventually constant at $-\infty$, using stages $[-2n,0]$; realize $\mathcal R(X)$ on $j\geq0$, using $[0,2n]$. Thus $q_L,q_R$ literally extend the other half by the value at zero. Padding is fully faithful, and all the constant-section and one-endpoint assertions follow from the interval contractions and filtered-colimit proof of E.20.

We prove the comparison on the actual half-fence by sliding an identity collar. Let $S:\mathcal L(X)\to\mathcal L(X)$ append two identity arrows at the right end, shifting the old path two places left. Thus $(Sc)(i)=c(\min(i+2,0))$. We give a homotopy $BS\simeq1$, relative to both endpoints. For $j\geq0$, define an endofunctor $\phi_j$ of the negative half-fence by
\[
 \phi_j(i)=
 \begin{cases}
 i+2,&i\leq-j-2,\\
 -j,&-j-2\leq i\leq-j,\\
 i,&i\geq-j.
 \end{cases}
\]
The values agree at the boundaries. Each generating arrow is either translated by two, unchanged, or sent to an identity; thus $\phi_j$ is a functor and fixes zero. Precomposition gives $F_j(c)=c\phi_j$. We have $F_0=S$. For even $j$ there is a natural transformation $F_j\to F_{j+1}$; for odd $j$ its direction is reversed. Its only possibly nonidentity components are at $-j-2,-j-1$, where both are the edge between $c(-j)$ and $c(-j-1)$. The naturality squares on the three adjacent generating edges commute; all other squares have identity components. The two endpoint components are identities.

For every finite diagram in $\mathcal L(X)$, the functors $F_j$ and these transformations become literally the identity on that diagram for sufficiently large $j$, because its paths and transformations are constant sufficiently far to the left. Concatenate the nerve homotopies on successive time intervals with endpoints $1-2^{-j}$, traversing a natural-transformation homotopy backwards when necessary, and put the identity at time $1$. This is continuous: on each nerve simplex it is stationary after finitely many intervals. In compactly generated spaces, realization has the weak topology with respect to its cells, and product with the compact interval preserves this cellwise test. We obtain $BS\simeq1$, on the total category and on every fixed-endpoint category. A related finite-stage identity-pair sliding calculation appears in \citet[\S3.2.1, Proposition 3.3]{Shoikhet2016}; the explicit formulas above establish the required homotopy for our colimit model independently of a categorical cocylinder identification.

Now $q_LS=T^{-2}q_L$ strictly. Every finite diagram in $\mathcal P(X)$ is carried into the image of $q_L$ by a sufficiently large negative even translation. The same holds for any finite diagram representing a homotopy. Homotopy groups and components of these nerves are computed at finite support by the filtered simplicial argument of E.20. These facts prove that $q_L$ is a weak equivalence, including on every fixed-endpoint fibre. Here is the basepoint detail. For a vertex $c\in\mathcal L(X;a,b)$, the collar homotopy supplies paths from $c$ to $S^kc$ inside that category. Use those paths to identify the relevant based homotopy groups. Translation $T^{-2k}$ gives an automorphism of the homotopy group of $B\mathcal P(X;a,b)$, and the relation $q_LS^k=T^{-2k}q_L$ shows that it preserves the image of the homotopy group of $B\mathcal L(X;a,b)$, in both directions, because $S^k$ is a homotopy equivalence. Any class translates into that image at some finite support, so it was already in the image. If a class from the half-fence becomes zero, translate a finite nullhomotopy into the half-fence; it kills its image under $S^k$, and invertibility of $S^k$ kills the original class. The same argument with paths proves bijectivity on components. Reversal of the fence proves the result for $q_R$.

The total and fibrewise weak equivalences $q_L,q_R$, together with E.21 and naturality of the maps to homotopy fibres, transfer the strict-fibre theorem to $\mathcal L$ and $\mathcal R$. Explicitly, the comparison square has the strict-fibre map on top, the map of homotopy fibres induced by $Bq_L$ or $Bq_R$ below, and the natural fibre comparisons vertically. Its top and bottom maps and its target vertical map are weak equivalences, so its source vertical map is one as well. This argument does not assert that a one-sided endpoint map has both strict categorical lifting properties.

If $j_L,j_R$ denote the constant sections, $q_Lj_L=j$ and $q_Rj_R=j$. Cancellation in the localized slice therefore gives $u_{\mathcal P}Bq_L=s(Bj_L)^{-1}$ and likewise for $q_R$, proving compatibility for these actual maps.

We will also use the elementary iterated-fibre identity. For $(f_1,f_2):A\to B_1\times B_2$, if $A_b=f_1^{-1}(b)\to\operatorname{hofib}_b(f_1)$ is a weak equivalence, regard both spaces over $B_2$ by $f_2$. Their derived fibres over $c$ agree; the latter is $\operatorname{hofib}_{(b,c)}(f_1,f_2)$. This follows from the model consisting of a point of $A$ and a pair of paths in the two base spaces. The maps from strict fibres in this identity are natural.
\end{proof}

\begin{claim}{E.24}{Proved}{The manuscript's left-right concatenation is associative}
For the block $D=(+,-)$ and $\Sigma=\Mor(X)$, glue a left-half path and a right-half path at their common value at zero. This is an actual functor
\[
 \kappa:\mathcal L(X)\times_X\mathcal R(X)\longrightarrow\mathcal P(X),
\]
where the pullback uses the terminal and initial endpoints. On fixed fibres define the manuscript's composition by
\[
 \mu_{a,b,c}=B\kappa\,\bigl((Bq_L)^{-1}\times(Bq_R)^{-1}\bigr):
 B\mathcal P(X;a,b)\times B\mathcal P(X;b,c)
       \longrightarrow B\mathcal P(X;a,c)
 \quad\hbox{in }\Hot.
\]
Here the inverses refer to the fixed-fibre weak equivalences of E.23. Under the canonical comparisons to ordinary path spaces, this \emph{particular} operation is ordinary path concatenation. It is therefore associative and unital in $\Hot$; reversal supplies inverses in the usual path-groupoid sense, and the loop object is a group object of $\Hot$.
\end{claim}
\begin{proof}
The gluing definition agrees on the common vertex and on the common component of a natural transformation, so it defines a functor on the indicated strict pullback. At stages $[-2n,0]$ and $[0,2n]$ it is exactly concatenation of two $D^n$-paths into a $D^{2n}$-path. Increasing $n$ pads at the two outer ends, exactly as in the manuscript's formula. This identifies the specified map, including its transition maps.

Write $E=\mathcal L(X)\times_X\mathcal R(X)$ and let $r:E\to X^3$ record the initial, middle, and terminal values. Its constant-pair section $j_E:X\to E$ is a weak equivalence. One proof uses the filtered system of finite supports: a stage with left length $2n$ and right length $2m$ is the functor category on their glued oriented interval, and its constant section is a weak equivalence by the same explicit interval contraction as in E.20. The constants commute with all enlargements.

We need a relative assertion as well. The middle-coordinate map $BE\to K$ is a weak equivalence since its composite with $Bj_E$ is $1_K$. Its strict fibre at $b$ is
\[
 B\mathcal L(X;?,b)\times B\mathcal R(X;b,?),
\]
which is contractible by E.23. Hence this strict fibre computes the middle-coordinate homotopy fibre. On it the remaining two coordinates are the product of the initial-endpoint map on $\mathcal L(X;?,b)$ and the terminal-endpoint map on $\mathcal R(X;b,?)$. Their strict fibres compute homotopy fibres: apply the iterated-fibre identity in E.23 to each full endpoint map, whose strict-fibre assertion is already proved, using its contractible one-endpoint fibre. Taking the remaining two derived fibres, and using that homotopy fibres commute with products, proves
\[
 B\bigl(\mathcal L(X;a,b)\times\mathcal R(X;b,c)\bigr)
       \simeq \operatorname{hofib}_{(a,b,c)}(BE\xrightarrow{Br}K^3).
\]
Every comparison in this paragraph is the natural comparison from the strict fibre, so it also controls the maps induced by projections and gluing.

Let $P=K^{[0,1]}$ and $Q=P\times_K P$, with their endpoint maps to $K^2$ and $K^3$. Both endpoint maps are fibrations; $Q$ is the space of maps from $[0,2]$ with values specified at $0,1,2$. Write $s:K\to P$ and $t:K\to Q$ for constants. These are weak equivalences. Define in the localized slices
\[
 u_L=s(Bj_L)^{-1},\qquad u_R=s(Bj_R)^{-1},\qquad
 u_{\mathcal P}=s(Bj)^{-1},\qquad v=t(Bj_E)^{-1}.
\]
For each projection from $E$ and for $\kappa$, regard the target over $K^3$ by adjoining the unused endpoint coordinate. This amounts to a product with $K$, which preserves weak equivalences. The two maps comparing categorical gluing with topological concatenation agree after composition with $Bj_E$: each sends a constant pair to a constant path, retaining the same middle coordinate. Since $Bj_E$ is invertible in this slice, cancellation proves
\[
                  u_{\mathcal P}B\kappa
                       =\operatorname{concat}\,v.
\]
The same cancellation proves that $v$ commutes with both projections under $u_L,u_R$.

Apply the derived-fibre functor at $(a,b,c)$. The relative assertion for $E$, E.21, E.23, and the ordinary path fibrations identify the derived fibres with the stated strict fibres. The two projection identities identify the induced map of the product fibre as $u_{L,a,b}\times u_{R,b,c}$, because maps into a product in $\Hot$ are determined by their two projections. Thus
\[
 u_{\mathcal P,a,c}\,B\kappa
 =\operatorname{concat}\,(u_{L,a,b}\times u_{R,b,c}).
\]
Compatibility with $q_L,q_R$ from E.23 gives the claimed formula for $\mu$. Ordinary concatenation is associative up to the endpoint-fixed homotopy which changes the two subdivision points in a threefold concatenation. Constants are units up to the analogous reparameterization homotopy. The reversal $c(j)\mapsto c(-j)$ of the integer fence is an actual functor, swaps endpoints, and preserves the constant section; cancellation in the same slices identifies it with reversal of topological paths. The standard contraction of a path followed by its reverse supplies the inverse identities. These prove all the asserted identities in $\Hot$ for the specified operation.
\end{proof}

\subsection{Relative collars and the remaining quantified comparisons}

We now make explicit the greater generality of VII ms.~108--136. An \emph{admissible} class $\Sigma\subseteq\Mor(X)$ contains every identity and is closed under composition. For a finite sign word $\tau$, let $I_\tau$ be the free category on its oriented interval. A $\Sigma$-path of type $\tau$ is a functor $I_\tau\to X$ taking every geometrically backward arrow to $\Sigma$. The category of such paths is the \emph{full} subcategory of $\Fun(I_\tau,X)$ on those objects. In particular, an interior component of a natural transformation is not required to belong to $\Sigma$; fixed-endpoint transformations have identity components at the two ends. These are the manuscript's definitions at VII ms.~65--67.

\begin{claim}{E.25}{Proved}{Canonical relative collars and literal even translation}
For any admissible $\Sigma$, let $\mathcal L_\Sigma(X;a,b)$, $\mathcal R_\Sigma(X;a,b)$, and $\mathcal P_\Sigma(X;a,b)$ be the full relative subcategories of the canonical $D=(+,-)$ models of E.23. The actual constant-tail extension maps $q_L,q_R$ are weak equivalences on these fibres and on their total categories. The literal translation $(T^2c)(j)=c(j-2)$ induces the identity in $\Hot$ on every such fibre. Gluing with the inverses of $q_L,q_R$ consequently defines the canonical relative operation $\mu_D$ in $\Hot$; constant paths are its left and right units.
\end{claim}
\begin{proof}
Each $\phi_j$ in E.23 is nondecreasing on the geometric vertex order and sends a generating edge to an edge of the same geometric direction or to an identity. Hence $F_j$ preserves relative paths. The homotopies between successive $F_j$ are also permitted, since the relative path category is full and their endpoint components are identities. The eventual-stationarity argument therefore proves $BS\simeq1$ on each relative left-half-fence fibre.

There is an equally explicit right-hand version, without invoking path reversal, which need not preserve a relative path class. For $i,j\geq0$ put
\[
 \psi_j(i)=
 \begin{cases}
 i,&i\leq j,\\
 j,&j\leq i\leq j+2,\\
 i-2,&i\geq j+2.
 \end{cases}
\]
Precomposition gives $G_j$. The functor $G_0=A$ inserts two identities at the left end of a right-half path. Consecutive $G_j$ have their two possibly nonidentity comparison components at $j+1,j+2$, given by the edge between $c(j)$ and $c(j+1)$; the transformation points $G_j\to G_{j+1}$ for even $j$ and backwards for odd $j$. These functors preserve relative paths for the same direction-preservation reason. They become the identity on each finite diagram sufficiently far to the right. Thus $BA\simeq1$, relative to the endpoints.

Even translation preserves the relative class and is an automorphism of the two-sided model. The finite-support and basepoint argument of E.23 now applies word for word, with $q_LS=T^{-2}q_L$ and $q_RA=T^2q_R$, to prove both comparison functors weak equivalences. This use of E.23 requires no endpoint-fibration theorem. Cancelling $Bq_L$ in $Bq_L\,BS=B(T^{-2})Bq_L$ gives $B(T^{-2})=1$, hence $B(T^2)=1$, on the actual relative fibre.

The gluing functor preserves the relative condition, so $\mu_D$ is well defined using the two inverses in $\Hot$. Gluing a left-half path to the constant right-half path is exactly $q_L$; gluing a constant left-half path to a right-half path is exactly $q_R$. Both comparisons preserve constants. Cancellation proves the two unit identities. We have not used the derived-fibre argument of E.24 to claim relative associativity: its strict endpoint-fibre hypothesis was established there only for the all-arrows case.
\end{proof}

To specify the remaining questions, we use the actual type category $\mathrm{Ty}$. Its objects are finite sign words, including the empty word $e$. A morphism is an endpoint-preserving nondecreasing functor between the associated oriented interval categories. Equivalently, if the source word is $\varepsilon_1\cdots\varepsilon_r$, a morphism is a choice of integers $l_i\geq0$ whose target word is $\varepsilon_1^{l_1}\cdots\varepsilon_r^{l_r}$; the $i$th source edge maps to that same-sign string, or to an identity when $l_i=0$. This is the description at VII ms.~50'--53. The \emph{useful length} $\ell_u(\tau)$ is the number of maximal constant-sign runs in the word, with $\ell_u(e)=0$ (VII ms.~75--76).

A projective sequence $t:\N^{\op}\to\mathrm{Ty}$ is \emph{left-cofinal in the manuscript's sense} when
\[
       \forall\tau\in\mathrm{Ty}\quad
       \exists n\quad\Hom_{\mathrm{Ty}}(t_n,\tau)\ne\varnothing.
\]
This is cofinality for the reachability preorder of VII ms.~99--102, not an assumption that comma categories are contractible. Equivalently $\ell_u(t_n)$ is unbounded, as stated at VII ms.~138: a source with sufficiently many alternating sign runs maps to any prescribed word by collapsing unused runs and lengthening the retained same-sign runs; a morphism cannot increase useful length. For such $t$ define the actual category
\[
 C_{\infty,\Sigma}(t;a,b)
   =\colim_n\Fun_\Sigma(I_{t_n},X)_{a,b},
\]
with transition functors given by precomposition with the specified transitions of $t$.

\begin{claim}{E.26}{Open}{The remaining cofinal-type and all-type comparisons}
For arbitrary admissible $\Sigma$ and arbitrary left-cofinal projective type sequences as just defined, this audit leaves undecided the following \emph{specified geometric} assertions of VII ms.~108--123:
\begin{enumerate}
\item A natural transformation $u:t\to t'$ of projective sequences, consisting of strictly commuting maps $u_n:t_n\to t'_n$, always induces a weak equivalence
\[
 u^*:C_{\infty,\Sigma}(t';a,b)\longrightarrow
                  C_{\infty,\Sigma}(t;a,b).
\]
\item The isomorphisms obtained from the canonical type projections $t*t'\to t,t'$ are independent of all comparison choices and form the transitive canonical system asserted at ms.~119--123.
\item The manuscript's comparison from a sequential colimit to its \emph{all-type} path category is an isomorphism in $\Hot$ for every such relative datum. Precisely, let $E_t$ be the fibred Grothendieck construction of the displayed sequence over $\N^{\op}$ and let $\mathrm{Ch}_\Sigma(X;a,b)\to\mathrm{Ty}$ have fibre $\Fun_\Sigma(I_\tau,X)_{a,b}$. The assertion is that
\[
 E_t\longrightarrow\mathrm{Ch}_\Sigma(X;a,b)
\]
is a weak equivalence, so that the span with the ordinary-colimit comparison $E_t\to C_{\infty,\Sigma}(t;a,b)$ yields the claimed isomorphism.
\end{enumerate}
\end{claim}
\begin{proof}[Status and missing hypotheses]
The ordinary-colimit leg of the last span is covered by E.8; its fibred convention is obtained from the covariant one by taking opposite categories. That does not prove its all-type projection is a weak equivalence. The comparison maps in the first two items exist and are the stated precomposition functors, but an isomorphism of their ind-objects does not establish their underlying weak equivalence or equality in $\Hot$; E.4 and E.8 explain the general failure of that inference. E.23 and E.25 establish the concrete collar comparisons, not this entire quantified family of transition systems. No counterexample to these particular geometric assertions, and no complete independent proof of their asserted generality, is supplied here. Thus ``Open'' records the precise remaining audit question, not a claim that the literature has designated it as an open problem.
\end{proof}

\begin{claim}{E.27}{Open}{The remaining relative and arbitrary-block composition assertions}
Fix an admissible $\Sigma$ and a word $\tau_0$ with $\ell_u(\tau_0)\geq2$. Form the three actual systems of VII ms.~125--126: the stages $\tau_0^n$ with transitions which collapse only the extra left blocks, the stages $\tau_0^n$ with transitions which collapse only the extra right blocks, and the stages $\tau_0^{2n}$ with symmetric outer collapse. Denote their colimits by $L_{\tau_0,\Sigma}$, $R_{\tau_0,\Sigma}$, and $P_{\tau_0,\Sigma}$, with endpoints understood. Constant padding defines actual maps $q_L:L_{\tau_0,\Sigma}\to P_{\tau_0,\Sigma}$ and $q_R:R_{\tau_0,\Sigma}\to P_{\tau_0,\Sigma}$, and concatenation defines the actual gluing functor $\kappa_{\tau_0}$. Beyond E.24--E.25, the audit does not settle all of the following claims of ms.~125--127: both padding maps are weak equivalences for every block; the resulting operations $\kappa_{\tau_0}(q_L^{-1}\times q_R^{-1})$ are independent of the block under the specified comparison maps; and those operations are associative in $\Hot$ for every admissible $\Sigma$.
\end{claim}
\begin{proof}[Status and smallest uncovered extensions]
For $\tau_0=D=(+,-)$, E.25 proves that the relative operation exists and is unital. The additional assertion of its associativity for an arbitrary proper admissible class $\Sigma\subsetneq\Mor(X)$ is already an extension not proved here. The proof of E.24 cannot simply be reused: it requires the actual strict relative endpoint fibres to be derived fibres of the diagonal on $BX$, a hypothesis absent for general $\Sigma$. E.9 supplies a different derived mapping-space construction only under its explicit rectification hypothesis and does not identify this raw operation with it.

Even with $\Sigma=\Mor(X)$, the written canonical-block argument does not verify every specified comparison for a different block. An explicit short uncovered convention is the alternating word $(+,-,+)$, whose useful length is three, with the left, right, and symmetric transitions specified in the claim. This is a smallest extension of block length beyond the two-letter canonical case, not a claim of a counterexample. A proof of its comparison and composition compatibility, or a proof for the full quantified family, would settle that extension. The general independence and associativity assertions cannot be inferred from their images under $\psi$. We retain them here with their original quantifiers so that the established all-arrows canonical result does not silently replace them.
\end{proof}


\partsection{F}{Finite tests and the information retained by higher morphisms}\label{sec:F}

Let $\mathcal A_{\mathrm{cat}}\subset\Spc$ be a full, essentially small
$\infty$-subcategory containing one representative of every weak homotopy
type $BA$ with $A$ a category having finitely many objects and arrows.
Let $\mathcal A_{\mathrm{fin}}\subset\Spc$ be the full subcategory of finite
CW homotopy types. We choose compatible small skeleta when forming
presheaf categories. The ordinary category called $\Hot_0$ in Chapter~V
is $h\mathcal A_{\mathrm{cat}}$. Neither the word ``finite'' nor a
statement about $\Ind$ will be allowed to identify these two subcategories.
All the mapping spaces in this section are derived mapping spaces of
unpointed spaces; brackets denote their components.

\begin{claim}{F.1}{Proved}{The two relevant finiteness conditions}
Every finite CW homotopy type is represented by a finite poset. A category
with finitely many arrows has a nerve with a finite $2$-skeleton, so each
component of its classifying space has finitely presented fundamental
group. Its nerve need not be homotopy equivalent to a finite CW complex.
In particular the one-object category $BC_2$, with arrows $1,g$ and
$g^2=1$, is finite, whereas $B(BC_2)$ is not a compact object of $\Spc$.
\end{claim}
\begin{proof}
A finite CW complex is homotopy equivalent to a finite simplicial complex
$L$ by simplicial approximation and the mapping-cylinder construction;
we use precisely the finite clause of \citet[Theorem~2C.5]{Hatcher2002}.
For completeness, the inductive construction replaces each attaching map
from a triangulated sphere by a simplicial approximation, and attaches a
finite simplicial mapping cone. The homotopy between the attaching maps
induces the comparison equivalence of their mapping cones. Finitely many
cells require finitely many such replacements. The poset $P(L)$ of
nonempty faces has nerve the barycentric subdivision of $L$, whose
realization is homeomorphic to $|L|$. The empty space is represented by
the empty poset.

For a finite category $A$, its nondegenerate $1$-simplices are its
nonidentity arrows and its nondegenerate $2$-simplices are composable
pairs of nonidentity arrows. Both sets are finite. In any connected
component, choose a maximal tree in the finite $1$-skeleton. The remaining
edges are a finite set of fundamental-group generators and the finitely
many $2$-simplices give a finite set of relations; cells of dimension
at least three do not change the fundamental group. This also proves
finiteness of the number of components.

In the normalized mod-$2$ bar complex of $BC_2$ there is exactly one
nondegenerate simplex $(g,\ldots,g)$ in every degree. Every internal face
is degenerate, because it replaces $g,g$ by $1$, and the two end faces
cancel modulo $2$. Consequently $H^q(B(BC_2);\mathbf F_2)\cong\mathbf F_2$
for every $q\geq0$. This rules out finite CW homotopy type. More precisely,
write $B(BC_2)\simeq\mathbf {RP}^{\infty}$ and use the filtered system
$\mathbf {RP}^{r}\hookrightarrow\mathbf {RP}^{r+1}$. If its identity
factored in $\Hot$ through any $\mathbf {RP}^{r}$, its pullback on
$H^q(-;\mathbf F_2)$ for $q>r$ would be zero, contrary to the identity.
Thus the canonical map
\[
 \colim_r[\mathbf {RP}^{\infty},\mathbf {RP}^{r}]
 \longrightarrow
 [\mathbf {RP}^{\infty},\mathbf {RP}^{\infty}]
\]
is not surjective. This proves noncompactness without replacing a
specified comparison by an abstract bijection. The multiplication table,
all associativity equations, and normalized boundaries through degree
nine are checked in \texttt{code/f\_checks.py}; the preceding symbolic
calculation proves the assertion in every degree. The cohomology-ring
calculation is also \citet[Theorem~3.19]{Hatcher2002}.
\end{proof}

\subsection{Set-valued tests}

Put
\[
 R_1^{\mathrm{cat}}:\Hot\longrightarrow
 \Fun((h\mathcal A_{\mathrm{cat}})^{\op},\Set),
 \qquad R_1^{\mathrm{cat}}(Y)(A)=[A,Y].
\]
The same notation with superscript $\mathrm{fin}$ uses
$h\mathcal A_{\mathrm{fin}}$.

\begin{claim}{F.2}{Refuted}{Faithfulness and conservativity of finite set-valued tests}
Both $R_1^{\mathrm{cat}}$ and $R_1^{\mathrm{fin}}$ fail to be faithful
and fail to be conservative. There is a single endomorphism
$s_B:BG\to BG$ such that $s_B\ne\id_{BG}$, $s_B$ is not an isomorphism
in $\Hot$, and both restricted functors send it to the identity.
\end{claim}
\begin{proof}
Let $G$ be the group of permutations of $\N=\{0,1,\ldots\}$ with finite
support. Define an injective group homomorphism $s:G\to G$ by
\[
 s(\sigma)(0)=0,\qquad s(\sigma)(i+1)=\sigma(i)+1.
\]
It is a homomorphism by substitution in the displayed formula. It is not
surjective because its image fixes $0$, whereas $(0\ 1)\in G$ does not.
Let $s_B=Bs$ be the map of classifying spaces, hence a morphism of
$\Hot$. It cannot be an isomorphism because its map on the fundamental
group is $s$.

If $P$ is any finitely generated group and $\rho:P\to G$ a homomorphism,
the supports of the images of finitely many generators have a finite
union. Choose $m$ such that this union lies in $\{0,\ldots,m\}$. Every
$\rho(p)$ fixes the complement of this set. For
$c=(0\ 1\ \cdots\ m+1)$, considered as an element of $G$, direct
conjugation gives
\[
       c\rho(p)c^{-1}=s(\rho(p))\qquad(p\in P).
\]
The program \texttt{f\_checks.py} verifies this identity for all
$5913$ permutations supported on initial sets of sizes one through
seven; the formula itself proves the identity for arbitrary finite
support.

For a connected CW complex $A$, the canonical classification of maps
to a $1$-type gives
\[
 [A,BG]\cong\Hom(\pi_1(A,a),G)/G,
\]
where $G$ acts by conjugation. One obtains this classification by
constructing the map on the $1$-skeleton from a homomorphism, extending
across $2$-cells using its relations, and then extending uniquely up to
homotopy over higher cells because $\pi_i(BG)=0$ for $i>1$.
Changing the image-basepoint path is exactly conjugation. Thus the
preceding calculation says $(s_B)_*=\id$ on $[A,BG]$ whenever
$\pi_1(A,a)$ is finitely generated. Apply F.1 to every connected component
of a finite-category nerve, or of a finite CW complex. Maps from a
disjoint union are the product of maps from its components, and the
empty test gives the singleton set. Hence the induced natural
endomorphism of either restricted presheaf is the identity. It agrees
with the image of $\id_{BG}$ although the two morphisms of $\Hot$ are
distinct. This proves both failures.

This is the finite-support instance attributed to
\citet[Proposition~4.1]{MatumotoMinamiSugawara1984} by
\citet[Remark~2.3]{ArlinChristensen2023}. We have checked the latter source;
the former paper was not accessible in this run, so its bibliography
entry is explicitly unverified. The proof above does not depend on
that attribution.
\end{proof}

\begin{claim}{F.3}{Refuted}{Fullness of finite set-valued tests}
The functors $R_1^{\mathrm{cat}}$ and $R_1^{\mathrm{fin}}$ are not full.
This is a separate failure from F.2: there are natural transformations
between their values which are not induced by any morphism in $\Hot$.
\end{claim}
\begin{proof}
Choose an uncountable $\U$-small set $I$ and let $F_I$ be the free group
with basis $\{x_i:i\in I\}$. Put $X=BF_I$ and keep the countable group
$G$ from F.2. For each finite subset $J\subset I$, choose pairwise
disjoint transpositions $t_j\in G$, $j\in J$, and let
$\lambda_J:F_J\to G$ be the homomorphism with
$\lambda_J(x_j)=t_j$. Two choices for the same finite labelled set are
simultaneously conjugate in $G$: match the finitely many pairs in their
supports and extend the resulting bijection by a permutation with finite
support. The same observation says that, if $J\subset J'$, the
restriction of $\lambda_{J'}$ to $F_J$ is conjugate to $\lambda_J$.

For a connected test $A$ in either test category, a homomorphism
$\rho:\pi_1(A,a)\to F_I$ has finitely generated image, and therefore
its image lies in $F_J$ for a finite set of basis letters $J$. Define
\[
 \eta_A:[A,X]\longrightarrow[A,BG],
 \qquad [\rho]\longmapsto[\lambda_J\rho].
\]
Passing to $J\cup J'$ proves independence of the chosen finite support.
If $\rho$ is replaced by a conjugate $h\rho h^{-1}$ in $F_I$, enlarge
$J$ to contain the finitely many letters of the word $h$; then applying
$\lambda_J$ changes the image by conjugation by $\lambda_J(h)$.
Thus $\eta_A$ is defined on free homotopy classes. Precomposing a
group homomorphism with the homomorphism induced by a map of test
spaces commutes with this construction; changes of basepoint are the
same conjugacies already checked. On disconnected test spaces use the
construction on each component. It respects maps which send different
components to the same component, because each component is treated
by precisely the same precomposition rule. For the empty test there
is the unique function between singleton sets. We have therefore
constructed an actual natural transformation, not merely a family of
functions.

Suppose $\eta$ were represented by $f:X\to BG$. Choose a group
homomorphism $\lambda:F_I\to G$ representing $f$. Testing the loop
$x_i$ by $S^1$ says that $\lambda(x_i)$ is a transposition. Testing
the pair $x_i,x_j$ by $S^1\vee S^1$ says that the pair
$(\lambda(x_i),\lambda(x_j))$ is simultaneously conjugate to two
disjoint transpositions. In particular $\lambda(x_i)\ne\lambda(x_j)$
when $i\ne j$. Both test spaces have finite poset presentations by
F.1. Thus $i\mapsto\lambda(x_i)$ injects the uncountable set $I$ into
the countable set $G$, a contradiction. All sets and classifying
spaces used here belong to the fixed universe. We searched the primary
sources listed in the report for this particular fullness argument;
we do not attach a priority claim to it.
\end{proof}

\begin{claim}{F.4}{Known}{No small set of unpointed tests repairs conservativity}
For every $\U$-small set $\mathcal G$ of $\U$-small spaces, there is
an endomorphism $f:X\to X$ of a $\U$-small connected $1$-type which
is not an isomorphism in $\Hot$ and satisfies
$[S,f]=\id_{[S,X]}$ for every $S\in\mathcal G$. In particular no such
set jointly reflects isomorphisms or separates all morphisms.
\end{claim}
\begin{proof}
This is the construction proving
\citet[Theorem~2.1]{ArlinChristensen2023}; we give its size argument.
Choose an uncountable regular cardinal $\alpha\in\U$ strictly larger
than the cardinality of $\pi_1(S,s)$ for every $S\in\mathcal G$ and
every component basepoint $s$. The collection of these groups is a
$\U$-small set, so a successor of a sufficiently large cardinal in
$\U$ has this property. Regard $\alpha$ as its initial ordinal and put
\[
 H=\bigcup_{\beta<\alpha}\operatorname{Sym}(\beta)
       \subset\operatorname{Sym}(\alpha),
\]
where each permutation fixes all ordinals outside $\beta$. This is a
$\U$-small group. Let $j:\alpha\hookrightarrow\alpha$ send $\gamma$
to $\gamma+1$. Define $s_H:H\to H$ by transporting a permutation along
$j$ on the successor ordinals and fixing every limit ordinal, including
$0$. Products are preserved, the map is injective, and it is not
surjective because every image fixes $0$. A bounded support remains
bounded after this operation, so it does map $H$ into $H$.

If $|P|<\alpha$ and $\rho:P\to H$, regularity bounds the supports of
all $\rho(p)$ in a single infinite limit ordinal $\beta<\alpha$.
There is a permutation $\tau\in H$ with
$\tau(j(\gamma))=\gamma$ for all $\gamma<\beta$: extend this
bijection from $j(\beta)$ to $\beta$ over a sufficiently large
bounded interval of cardinality $|\beta|$, using an extra interval
of that cardinality to make the two complements equinumerous, and
fix its complement in $\alpha$. The bound remains below $\alpha$
because its cardinality is below $\alpha$. Since $s_H\rho(p)$ is
supported on $j(\beta)$, we obtain
$\tau s_H\rho(p)\tau^{-1}=\rho(p)$ on that support and on its
complement. Thus $Bs_H$ induces the identity on maps from every
connected component of every $S\in\mathcal G$. Products over the
components finish the proof, including arbitrarily many components.
\end{proof}

\begin{claim}{F.5}{Known}{The stronger nonconcreteness obstruction}
The category $\Hot$ admits no faithful functor to $\Set_{\U}$ with
the usual fixed-universe interpretation. The obstruction already
occurs on simply connected, finite-dimensional CW complexes. Therefore
faithfulness of a restricted set-valued Yoneda functor cannot be
obtained by substituting any other $\U$-small full test subcategory.
\end{claim}
\begin{proof}
We use the main theorem of \citet[p.~1 and the conclusion on p.~5 of
the 2004 reprint]{Freyd1970}. The theorem states the failure of every
homotopy-invariant set-valued functor to separate a nonnull map from
the null map. The construction on p.~5 works on $(r+3)$-dimensional,
$r$-connected CW complexes for every $r\geq1$, so based versus
unbased homotopy classes cause no exception. This is an invocation
of the proved theorem, not a new proof of nonconcreteness. For a
$\U$-small category $\mathcal A$, the functor
\[
 \Fun(\mathcal A^{\op},\Set_{\U})\longrightarrow\Set_{\U},
 \qquad F\longmapsto\coprod_{A\in\Ob\mathcal A}F(A)
\]
is faithful: the tag $A$ retains the component of each natural
transformation. A faithful restricted Yoneda functor followed by
this functor would contradict Freyd's theorem.
\end{proof}

\begin{claim}{F.6}{Proved}{The telescope phantom and the failure of epimorphy}
There is an essential unpointed map
$c:T\to K(\Z,2)$, where
$T=\operatorname{Tel}(S^1\xrightarrow{2}S^1\xrightarrow{2}\cdots)$,
such that $c\circ a$ is nullhomotopic for every map $a:BA\to T$
from a finite category. If $T_r$ denotes the $r$th finite
subtelescope, the canonical map $(T_r)_r\to T$ in $\Ind(\Hot)$ is
not an epimorphism.
\end{claim}
\begin{proof}
The telescope has fundamental group $\Z[1/2]$ by van Kampen. Any
map from a sphere of dimension at least two, and any specified
nullhomotopy, lies in a finite subtelescope after cellular
approximation; the latter retracts onto a circle. Hence
$T\simeq K(\Z[1/2],1)$. Contracting the ray of vertices gives a
homotopy equivalent CW complex with one vertex, loops $e_r$, and
one $2$-cell for each relation $e_r=2e_{r+1}$ in its cellular
boundary. Its cochain differential is
\[
 d:\prod_{r\geq0}\Z\longrightarrow\prod_{r\geq0}\Z,
 \qquad d((b_r))=(b_r-2b_{r+1})_r.
\]
Consequently $H^2(T;\Z)=\operatorname{coker}d$. Set $a_{2j}=1$
and $a_{2j+1}=0$. If $a=d(b)$, then for every $N\geq0$
\[
 b_0=\sum_{j=0}^{N}2^ja_j+2^{N+1}b_{N+1}.
\]
In the $2$-adic integers the right-hand partial sum tends to
$1+4+16+\cdots=-1/3$. Since $\Z\hookrightarrow\Z_2$ is injective,
this would give $3b_0=-1$ in $\Z$, which is impossible. Thus $[a]$
is a specified nonzero element of $H^2(T;\Z)$. Its representing
map $c$ exists by \citet[Theorem~4.57]{Hatcher2002}; that theorem
also gives the unpointed identification because the target is
simply connected. The entire cohomology group is
$\Z_2/\Z$, as calculated in
\citet[Examples~3F.7 and~3F.9]{Hatcher2002}, although nonvanishing
alone suffices here.

Every homomorphism from a finitely generated group into
$\Z[1/2]$ has cyclic image, since finitely many denominators have
a common power of two. F.1 and the classification of maps into a
$1$-type therefore imply that a map $BA\to T$, on each connected
component, factors in $\Hot$ through a circle or a point.
Because $[S^1,K(\Z,2)]=0$, its composite with $c$ is nullhomotopic.
No finite-dimensionality assertion about $BA$ is used. A finite
subtelescope also retracts onto a circle, whence
\[
 \Hom_{\Ind(\Hot)}((T_r)_r,K(\Z,2))
       =\lim_r[T_r,K(\Z,2)]=\{0\}.
\]
Precomposing the two distinct arrows $c,0:T\rightrightarrows K(\Z,2)$
with $(T_r)_r\to T$ therefore makes them equal. This is precisely
the failure of that arrow to be an epimorphism.
\end{proof}

\subsection{Space-valued tests}

\begin{claim}{F.7}{Proved}{The exact Ind-completion assertion}
In the fixed universe there is an equivalence
\[
 \Ind_{\omega}(\mathcal A_{\mathrm{fin}})\simeq\Spc
\]
whose restriction to $\mathcal A_{\mathrm{fin}}$ is its given
inclusion. Thus
\[
 R_{\infty}^{\mathrm{fin}}:\Spc\longrightarrow
 \Fun(\mathcal A_{\mathrm{fin}}^{\op},\Spc),
 \qquad X\longmapsto\Map(-,X)|_{\mathcal A_{\mathrm{fin}}}
\]
is fully faithful. The compact objects of $\Spc$ are exactly the
retracts in $\Spc$ of finite CW homotopy types; no assertion that
all such retracts have finite CW type is needed.
\end{claim}
\begin{proof}
The point is compact: its derived mapping-space functor is the
identity. Mapping out of a finite homotopy colimit gives a finite
homotopy limit, and filtered colimits of spaces commute with finite
homotopy limits \citep[Proposition~5.3.3.3]{Lurie2009}. The latter statement can be checked on Kan models:
a map from a finite simplicial set and any of its finitely many
specified horn fillings occurs at one stage; equality and homotopies
of such finite data occur at a later common stage. A finite CW
complex is obtained from the empty space and points by finitely
many homotopy pushouts attaching cells. Induction proves its
compactness for filtered colimits of spaces, not just compactness
of its set of homotopy classes.

A simplicial model for any $\U$-small space is the union of its
finite simplicial subsets. This union is filtered and is its
homotopy colimit, since inclusions are cofibrations; its finite
pieces have finite CW realizations. Hence finite CW homotopy types
generate $\Spc$ under filtered colimits. These two facts verify
all three hypotheses of
\citet[Proposition~5.3.5.11(2)]{Lurie2009}, giving the claimed
equivalence. Alternatively, if
$X=\colim_iK_i$ and $Y=\colim_jL_j$ are such filtered presentations,
compactness and the universal property of the source colimit give
\[
 \Map(X,Y)\simeq
 \lim_i\Map(K_i,Y)\simeq
 \lim_i\colim_j\Map(K_i,L_j),
\]
which is exactly the mapping-space formula in the Ind-category.
These are canonical comparison equivalences, natural in both
variables. The restricted Yoneda functor is the defining embedding
of that Ind-category into presheaves.

If $X$ is compact, compactness applied to its finite-piece
presentation makes $\id_X$ factor in $\Hot$ through some $K_i$.
The corresponding homotopy $r i\simeq\id_X$ gives a homotopy
retract; replacing $i$ by its mapping cylinder realizes the retract
datum in $\Spc$. Conversely a retract of a compact object is compact,
because its mapping-space functor is a retract of the corresponding
compact mapping-space functor. This is also the idempotent-completion
statement \citet[Lemma~5.4.2.4]{Lurie2009}.
\end{proof}

\begin{claim}{F.8}{Proved}{Finite-category space-valued tests are dense, but are not all compact}
The functor
\[
 R_{\infty}^{\mathrm{cat}}:\Spc\longrightarrow
 \Fun(\mathcal A_{\mathrm{cat}}^{\op},\Spc)
\]
is fully faithful. Nevertheless the filtered-colimit extension
$\Ind_{\omega}(\mathcal A_{\mathrm{cat}})\to\Spc$ of the inclusion
is not an equivalence. Thus ``spaces are the Ind-completion of
finite spaces'' proves the intended positive statement only after
specifying finite CW homotopy types, or their retract completion,
as the compact generators.
\end{claim}
\begin{proof}
F.1 gives an inclusion
$i:\mathcal A_{\mathrm{fin}}\hookrightarrow\mathcal A_{\mathrm{cat}}$.
For $X\in\Spc$, the presheaf $R_{\infty}^{\mathrm{cat}}X$ is the
right Kan extension of $R_{\infty}^{\mathrm{fin}}X$ along
$i^{\op}$. Indeed at $A\in\mathcal A_{\mathrm{cat}}$ the
right-Kan-extension formula is
\[
 \lim_{(K\to A)\in(\mathcal A_{\mathrm{fin}}/A)^{\op}}
       \Map(K,X),
\]
where $\mathcal A_{\mathrm{fin}}/A=\mathcal A_{\mathrm{fin}}
\times_{\Spc}\Spc_{/A}$. The density statement in F.7
identifies this canonical limit with $\Map(A,X)$. Restriction
followed by this right Kan extension is the identity on these
presheaves; right Kan extension along a fully faithful inclusion
is fully faithful. Therefore, using F.7, maps between
$R_{\infty}^{\mathrm{cat}}X$ and $R_{\infty}^{\mathrm{cat}}Y$
are exactly maps $X\to Y$, with their entire mapping spaces.

Every representable object of
$\Ind_{\omega}(\mathcal A_{\mathrm{cat}})$ is compact. An
equivalence extending the inclusion would consequently make
$B(BC_2)$ compact in $\Spc$, contrary to F.1. The negative
assertion concerns the specified extension of the inclusion;
we do not infer it merely from the fact that a nerve is
infinite-dimensional.
\end{proof}

\subsection{Finite categorical height}

For each integer $n\geq1$, let $h_n\Spc$ have the same objects as
$\Spc$ and mapping spaces $\tau_{\leq n-1}\Map(X,Y)$. We use
``$n$-category'' in this $(n,1)$ sense. This convention retains
homotopy classes of homotopies at $n=2$; it does not mean the
category of $n$-truncated spaces. For either test subcategory
$\mathcal A$, put
\[
 R_n^{\mathcal A}(X)=
 \tau_{\leq n-1}\Map(-,X)|_{\mathcal A}
 \ \in\ \mathcal P_{n-1}(\mathcal A)
 :=\Fun(\mathcal A^{\op},\Spc_{\leq n-1}).
\]
Because the target is an $n$-category this functor factors through
$h_n\Spc$. The associated functor on the homotopy category has
maps given by coherent natural transformations modulo coherent
homotopy. At $n=2$ these are pseudonatural transformations and
modifications in an equivalent groupoid-enriched model. In
particular, a collection of unrelated homotopies at individual
test spaces is not a modification.

\begin{claim}{F.9}{Known}{Conservativity returns at height two}
For every integer $n\geq2$, both $R_n^{\mathrm{fin}}$ and
$R_n^{\mathrm{cat}}$ reflect equivalences. In fact a map
$f:X\to Y$ is a weak equivalence if
\[
 \Pi_1\Map(S^r,f):\Pi_1\Map(S^r,X)
             \longrightarrow\Pi_1\Map(S^r,Y)
\]
is an equivalence of groupoids for every $r\geq0$.
\end{claim}
\begin{proof}
The second assertion is
\citet[Theorem~4.3]{ArlinChristensen2023}. We give a direct proof
which also checks disconnected and empty spaces. The point is a
retract of $S^0$, so the hypothesis implies that
$\Pi_1 f:\Pi_1X\to\Pi_1Y$ is an equivalence. Thus $f$ is a
bijection on components and an isomorphism on every based
fundamental group. If one space is empty, the other is empty and
there is nothing further to prove. Fix a component and a basepoint
$x\in X$, and use $f$ to identify the two fundamental groups with
a group $\Gamma$.

For $r\geq2$, free homotopy classes of maps from $S^r$ into this
component are the orbit set of $\pi_r(X,x)$ under the
$\Gamma$-action. This follows by recording the path traced by
the basepoint during a free homotopy. The map
$u_r:\pi_r(X,x)\to\pi_r(Y,f(x))$ is a
$\Gamma$-equivariant group homomorphism, and the hypothesis gives
a bijection of its orbit sets. Surjectivity of the orbit map
implies surjectivity of $u_r$: write a target element as
$\gamma u_r(a)=u_r(\gamma a)$. If $u_r(a)=0$, injectivity of the
orbit map forces the orbit of $a$ to equal the orbit of $0$,
which is the singleton $\{0\}$. Thus the kernel is zero and
$u_r$ is injective. This proves the isomorphism on every
homotopy group, hence the weak equivalence. For $n>2$ apply
$\tau_{\leq1}$ to the assumed equivalences of mapping
$(n-1)$-types, and use the case just proved. Spheres and the
point have finite poset presentations, so both test categories
contain all the tests used.
\end{proof}

\begin{claim}{F.10}{Refuted}{Faithfulness at the highest retained dimension}
For every finite $n\geq1$, neither $R_n^{\mathrm{fin}}$ nor
$R_n^{\mathrm{cat}}$ is faithful on the highest retained
morphisms. More precisely, for $T$ from F.6, $Y_n=K(\Z,n+1)$,
and the constant map $0:T\to Y_n$, the comparison on
$\pi_{n-1}$ of mapping spaces has a nonzero kernel. For $n=1$
this means the map on components. Consequently these functors
are not fully faithful as $n$-functors for any finite $n$.
\end{claim}
\begin{proof}
Let $\alpha\ne0$ be the class in $H^2(T;\Z)$ constructed in F.6.
For $n\geq2$, the exponential adjunction and the Eilenberg--Mac~Lane
classification give
\[
 \pi_{n-1}\Map(T,K(\Z,n+1),0)
   \cong H^{n+1}(S^{n-1}\times T,\{*\}\times T;\Z)
   \cong H^2(T;\Z).
\]
The second isomorphism is the relative suspension isomorphism;
its integral sphere factor is free, so no Tor term occurs. The
first classifies families based at the constant map. At $n=1$
use $[T,K(\Z,2)]=H^2(T;\Z)$ instead. Pulling back this class
along any $a:A\to T$ from either finite-test category gives
$a^*\alpha=0$ by F.6. Hence the associated top-dimensional
morphism becomes zero at every object and every mapping-space
component of the test presheaf.

A top-dimensional morphism between maps of $(n-1)$-truncated
presheaves is determined by these components. One can see this
without a strict-model convention: its $(n-1)$-fold loop space
is a limit of discrete spaces, so evaluation at all objects and
all mapping components is injective on its set of points.
Therefore the image of $\alpha$ is zero as a coherent
presheaf morphism, not just at a selected basepoint. Since
$\alpha$ was nonzero before applying $R_n$, the asserted kernel
is nonzero. For $n=2$ this is an explicit phantom homotopy,
namely a nonidentity automorphism of the constant map in the
fundamental groupoid of $\Map(T,K(\Z,3))$.
\end{proof}

\begin{claim}{F.11}{Proved}{A precise gain for countable sources}
Let $X$ have countable CW homotopy type, let $Y$ be any $\U$-small
space, let $n\geq1$, and use finite CW tests. The canonical map
\[
 \theta_n:\tau_{\leq n-1}\Map(X,Y)
   \longrightarrow
 \Map_{\mathcal P_{n-1}(\mathcal A_{\mathrm{fin}})}
       (R_n^{\mathrm{fin}}X,R_n^{\mathrm{fin}}Y)
\]
is surjective on components. It induces isomorphisms on
$\pi_k$ for $0\leq k<n-1$ and surjections on $\pi_{n-1}$,
with basepoints wherever required. In particular coherent
height-two finite CW tests determine equality of homotopy
classes of maps with countable CW source. The same conclusion
about equality holds if one is given coherent tests on all
finite-category homotopy types.
\end{claim}
\begin{proof}
Choose finite subcomplexes
$K_0\subset K_1\subset\cdots$ exhausting a countable CW model
for $X$. Such an exhaustion exists by closure-finiteness and a
countable enumeration of its cells. These inclusions are
cofibrations, so the union is the homotopy colimit and
\[
 \Map(X,Y)\simeq\lim_r\Map(K_r,Y).
\]
Compactness of each test $A\in\mathcal A_{\mathrm{fin}}$ gives
$R_{\infty}^{\mathrm{fin}}X\simeq
\colim_r R_{\infty}^{\mathrm{fin}}K_r$ in the presheaf
$\infty$-category. Objectwise truncation is left adjoint to
inclusion of truncated presheaves. Combining that adjunction
with the preceding colimit and Yoneda gives a canonical
equivalence
\[
 \Map_{\mathcal P_{n-1}}(R_nX,R_nY)
      \simeq\lim_r\tau_{\leq n-1}\Map(K_r,Y).
\]
Under these identifications $\theta_n$ is precisely the
comparison between truncating a homotopy inverse limit and
taking the homotopy inverse limit of the truncations.

Write $m=n-1$. Each map
$\Map(K_r,Y)\to\tau_{\leq m}\Map(K_r,Y)$ has
$m$-connected homotopy fibres. Over any point of the
homotopy limit of the targets, the homotopy fibre of the
comparison is a homotopy inverse limit of a tower of
$m$-connected nonempty spaces. A tower of nonempty connected
spaces has nonempty homotopy inverse limit: after replacing
bonding maps by fibrations, choose points and their path lifts
recursively. More generally, the homotopy inverse limit of a
tower of $m$-connected spaces is $(m-1)$-connected. This
follows either by the same recursive extension over spheres,
or from the homotopy-limit fibre description as the homotopy
equalizer of the identity and shift on the product, whose long
exact sequence has zero homotopy groups in degrees below $m$.
For $m=0$ the assertion is only nonemptiness. Thus the
comparison is $m$-connected, which yields exactly the
surjectivity and isomorphisms stated. Truncating its source
at $m$ does not change these groups.

For $n=2$ the map on $\pi_0$ is a bijection. Thus a coherent
natural equivalence between the two images of maps
$f,g:X\to Y$ implies $f=g$ in $\Hot$. If the coherent
equivalence was given on $\mathcal A_{\mathrm{cat}}$, restrict
it to $\mathcal A_{\mathrm{fin}}$ before applying this
argument. We do not infer a fullness assertion on
$\mathcal A_{\mathrm{cat}}$, since its noncompact tests do not
satisfy the colimit comparison used above. F.10 shows that the
possible failure of injectivity on the highest group is real,
even for this countable source.
\end{proof}

\begin{claim}{F.12}{Refuted}{Ordinary fullness and faithfulness at every finite height}
For every finite $n\geq2$, neither $hR_n^{\mathrm{fin}}$ nor
$hR_n^{\mathrm{cat}}$ is full or faithful on the ordinary sets
of homotopy classes of maps. Nonfaithfulness has examples whose
source has CW cardinality $\aleph_{n-1}$; nonfullness has examples
whose source has CW cardinality $\aleph_n$. Moreover, for each
$0\leq k\leq n-1$, the map on the retained $\pi_k$ of a mapping
space can fail to be injective. All the groups, modules, and spaces
in these examples are $\U$-small.
\end{claim}
\begin{proof}
We first record the group-cohomological input and its cardinal control.
For $d\geq1$, put
\[
 G_d=\bigoplus_{\aleph_{d-1}}\Z/2.
\]
This is a locally finite abelian group of cardinality $\aleph_{d-1}$.
The established theorem
\citet[Corollary~6.10(2)]{BC_DKLT2002}, attributed there to Chen,
states that a torsion abelian group $G$ of cardinality $\aleph_r$
has cohomological dimension $r+1$ over any nonzero ring in which
the orders of its finite subgroups are invertible. In particular
\[
       \operatorname{cd}_{\Q}G_d=d.
\]
Thus there is a $\Q G_d$-module $M$ with
$H^d(G_d;M)\ne0$. We can choose such an $M$ of cardinality at
most $\aleph_{d-1}$: in the rational free bar resolution of the
trivial module $\Q$, let $\Omega^i\Q$ be its $i$th syzygy and
put $M=\Omega^d\Q$. The sequence
\[
 0\longrightarrow\Omega^d\Q\longrightarrow P_{d-1}
   \longrightarrow\Omega^{d-1}\Q\longrightarrow0
\]
does not split, since $\Omega^{d-1}\Q$ is not projective when
$\operatorname{pd}_{\Q G_d}\Q=d$. Its nonzero extension class
corresponds, by dimension shifting, to a nonzero element of
$\operatorname{Ext}^d_{\Q G_d}(\Q,M)=H^d(G_d;M)$.
Each bar module has cardinality $\aleph_{d-1}$, so the syzygies
have the required bound. For $d=2$ one may more explicitly take
$M=\Q G_2$, since \citet[Theorem~5.4]{BC_DKLT2002} proves that
$H^2(G_2;\Q G_2)\ne0$. These are invocations of the cited proved
group-cohomology theorems, not of Holt's conjectures for general
locally finite groups.

For a $\Q G$-module $M$ and $r\geq2$, write
\[
 L_G(M,r)=EG\times_G K(M,r)\longrightarrow BG
\]
for the split Eilenberg--Mac~Lane fibration, where the action on
$M$ is realized functorially on the simplicial abelian group
$K(M,r)$. It has the zero section. Its derived space of sections
over $BH\to BG$, for any subgroup $H\leq G$, is
\[
 \Gamma_H(M,r)=\Map_H(EH,K(M,r)),\qquad
 \pi_i\Gamma_H(M,r)=H^{r-i}(H;M)\quad(0\leq i\leq r),
 \tag{F.12a}
\]
and its homotopy groups above degree $r$ vanish. The equation for
$i=0$ is the classification of components as an abelian group.
One explicit verification uses the free simplicial $\Q H$-module
on $EH$: its normalized complex is the rational bar resolution.
Equivariant maps to the Dold--Kan model of $M[r]$ give the good
nonnegative truncation of the cochain complex
$\Hom_{\Q H}(P_\bullet,M)$ shifted by $r$. Its homology in
degree $i$ is precisely $H^{r-i}(H;M)$. This proves (F.12a),
including its naturality under restriction of subgroups.
When $H$ is finite, the invariants functor on $\Q H$-modules is
exact: an arbitrary lift of an invariant vector can be averaged
by $|H|^{-1}\sum_{h\in H}h$. Consequently
\[
 H^j(H;M)=0\ (j>0),\qquad
       \Gamma_H(M,r)\simeq K(M^H,r).
 \tag{F.12b}
\]

We next explain why this calculation controls coherent natural
transformations, rather than merely their values at individual tests.
Fix $n\geq2$, put $m=n-1$, and abbreviate
$\mathcal P_m(\mathcal A_{\mathrm{cat}})$ to $\mathcal P_m$.
For a locally finite group $G$, let $\mathcal F(G)$ be the filtered
poset of its finite subgroups. There is a canonical equivalence
\[
 R_n^{\mathrm{cat}}(BG)
     \simeq\colim_{H\in\mathcal F(G)}R_n^{\mathrm{cat}}(BH).
 \tag{F.12c}
\]
To verify this, evaluate first before truncation at a finite-category
test $A$. Each component of $A$ has finitely generated fundamental
group by F.1. Maps from that component to $BG$ form the groupoid
of group homomorphisms to $G$ and conjugating elements in $G$.
A homomorphism has finite image, since $G$ is locally finite,
and finitely many specified conjugating elements lie in a common
finite subgroup with that image. Thus this mapping groupoid is
the filtered colimit of the corresponding mapping groupoids for
$BH$. There are only finitely many components of $A$, so the
same assertion holds for the product over its components.
Taking nerves verifies the assertion on mapping spaces. Objectwise
truncation preserves colimits into $\mathcal P_m$, giving (F.12c).

Let $p:E\to BG$ be a fibration and consider the space of sections
of its image in $\mathcal P_m$. By (F.12c) and the universal
property of a colimit in the slice category, it is
\[
 \begin{aligned}
 &\Map_{\mathcal P_m/R_nBG}(R_nBG,R_nE)\\
 &\qquad\simeq
 \lim_{H\in\mathcal F(G)^{\op}}
 \operatorname{hofib}_{BH\hookrightarrow BG}
 \left(\tau_{\leq m}\Map(BH,E)
          \longrightarrow\Map(BH,BG)\right).
 \end{aligned}
 \tag{F.12d}
\]
Here $BH$ is itself a finite-category test, so truncated Yoneda
identifies the first mapping space inside each fibre exactly.
The second mapping space is a $1$-type, already $m$-truncated.
For a map with $m$-truncated base, taking the homotopy fibre
commutes with $m$-truncation of its total space. This follows
from the homotopy-group exact sequence, since the base has no
homotopy groups above degree $m$. Hence each term of (F.12d)
is $\tau_{\leq m}$ of the actual space of sections of $p$ over
$BH$. In particular, if those section spaces have homotopy type
$K(M^H,r)$ with $r>m$, every term of (F.12d) is contractible,
and so is their limit. This last conclusion is a fully coherent
statement; no compatible choice of individual nullhomotopies
has been assumed.

For nonfaithfulness, take $G=G_n$ and a nonzero
$\alpha\in H^n(G;M)$ as above, and put $E=L_G(M,n)$.
The zero section $s_0$ and a section $s_\alpha$ representing
$\alpha$ are different components of the actual section space
by (F.12a). They also determine different \emph{unpointed} homotopy
classes $BG\to E$. Indeed, the homotopy fibre sequence
\[
 \Gamma_G(M,n)\longrightarrow\Map(BG,E)
       \xrightarrow{p_*}\Map(BG,BG)
\]
shows that two fibre components can merge in the total space
only if they belong to the same orbit of
$\pi_1(\Map(BG,BG),\id)$. Composition with $s_0$ is a section
of $p_*$, so every loop of the base lifts to a loop at $s_0$.
Its action fixes the zero component. That orbit is a singleton,
and cannot contain $\alpha\ne0$.
On the other hand, (F.12b)--(F.12d) make the space of sections
of $R_np$ contractible. The two induced transformations are
therefore homotopic in $\mathcal P_m$. This proves failure of
ordinary faithfulness for finite-category tests. Restricting
the coherent homotopy proves the same failure for finite CW
tests.

For the additional assertion about every retained homotopy group,
keep $G=G_n$ and replace $E$ by $L_G(M,n+k)$ for
$1\leq k\leq n-1$. Formula (F.12a) puts the nonzero group
$H^n(G;M)$ in $\pi_k$ of its section space at zero.
The base $\Map(BG,BG)$ is a $1$-type, so the fibre exact
sequence injects this group into $\pi_k\Map(BG,E)$; for
$k\geq2$ it is an isomorphism onto that homotopy group.
Formula (F.12d) again makes the presheaf section space
contractible, so this nonzero element maps to zero.

For nonfullness, take $G=G_{n+1}$ and choose
$0\ne\beta\in H^{n+1}(G;M)$ with the same cardinal control.
Let $z,s_\beta$ be the zero and $\beta$ sections of
$L_G(M,n+1)\to BG$. Define $E_\beta\to BG$ to be the
fibrewise path fibration whose fibre at $b$ consists of paths
in the fibre of $L_G(M,n+1)$ from $z(b)$ to $s_\beta(b)$.
Its fibres are $K(M,n)$, and a homotopy section would give a
path from $z$ to $s_\beta$ in $\Gamma_G(M,n+1)$, contradicting
$\beta\ne0$. Thus $E_\beta\to BG$ has no homotopy section.
For finite $H\leq G$, (F.12b) implies that the two restricted
sections are connected by a path, and their path space has
homotopy type $K(M^H,n)$. Formula (F.12d) therefore provides
a section
\[
 \eta:R_nBG\longrightarrow R_nE_\beta
 \quad\text{with}\quad R_np\circ\eta\simeq\id_{R_nBG}
\]
in $\mathcal P_m$.

This section cannot be induced, even up to coherent homotopy,
by a map $f:BG\to E_\beta$. The reason, valid for either test
category $\mathcal A$, is the exact comparison for any $1$-type
$Z$ and $n\geq2$:
\[
\begin{split}
 \Map_{\mathcal P_m(\mathcal A)}(R_nX,R_nZ)
 &\simeq
 \Map_{\Fun(\mathcal A^{\op},\Spc)}
                 (R_\infty X,R_\infty Z)\\
 &\simeq\Map(X,Z).
\end{split}
 \tag{F.12e}
\]
The first equivalence is the truncation adjunction, because
$R_\infty Z$ already takes values in $1$-types; the second
is F.7 or F.8. If $R_nf\simeq\eta$, (F.12e) applied to
$Z=BG$ would give $pf\simeq\id_{BG}$, an impossible homotopy
section. Thus $hR_n^{\mathrm{cat}}$ is not full. Restrict
$\eta$ to finite CW tests; the same argument using (F.12e)
shows that it has no lift there either. Finally, the bar and
Eilenberg--Mac~Lane constructions and the finite homotopy
pullbacks used above preserve the indicated infinite cardinal
bounds. The entire construction remains in $\U$.
\end{proof}

This counterexample identifies why the tower argument in F.11 cannot
be extended to arbitrary filtered indexing categories. The canonical
comparison for finite CW tests remains
\[
 \tau_{\leq n-1}\lim_{i\in I^{\op}}\Map(K_i,Y)
   \longrightarrow
 \lim_{i\in I^{\op}}\tau_{\leq n-1}\Map(K_i,Y),
\]
but F.12 supplies actual failures of both injectivity and surjectivity
on its components. At height two the nonfaithful source has cardinality
$\aleph_1$, so the countable-source restriction in F.11 cannot simply
be removed. The higher Brown representability theorem
\citep[Theorem~4.2.11]{NguyenRaptisSchrade2022} concerns functors
on the full $n$-category satisfying its exactness hypotheses; it
does not assert fullness of these restricted finite-test presheaves.
Likewise the conservative comparison theorems in
\citet[Theorems~35 and~37]{Arlin2020} do not assert that these
ordinary morphism sets are faithfully represented.

\begin{center}
\begin{tabular}{L{0.16\linewidth}L{0.18\linewidth}L{0.22\linewidth}L{0.30\linewidth}}
\hline
Height & Equivalence detection & Highest-dimensional faithfulness & Ordinary map-class fullness and faithfulness\\
\hline
$n=1$ & False (F.2) & False (F.2) & Both false (F.2--F.3)\\
$2\leq n<\infty$ & True (F.9) & False (F.10, F.12) & Both false in general (F.12); finite CW tests give both for countable sources (F.11)\\
$n=\infty$ & True & True & Fully faithful, including every coherence (F.7--F.8)\\
\hline
\end{tabular}
\end{center}

\subsection{The identities sought in Chapters VII and IX}

\begin{claim}{F.13}{Proved}{What testing can justify for composition and reversal}
Let $P$ be a space with a specified morphism
$m:P\times P\to P$ in $\Hot$, and let $T:P\to P$ be a
specified endomorphism. Equality of the two associativity
composites, or of $T^2$ and $\id_P$, does not follow in general
from equality after every finite set-valued test. For each fixed
finite height there are actual operations and endomorphisms for
which the corresponding coherent finite-test identities hold
but these identities in $\Hot$ fail. A coherent
homotopy of their space-valued finite-test natural
transformations does imply the corresponding equality in
$\Hot$. A coherent height-two finite-test homotopy suffices
when the source of the proposed equality has countable CW
homotopy type. For arbitrary sources no fixed finite height
supplies a general equality criterion: F.12 gives distinct
ordinary homotopy classes with coherently identical finite tests
at every prescribed finite height.
\end{claim}
\begin{proof}
The associativity composites are the two typed maps
\[
 m\circ(m\times\id_P),\quad m\circ(\id_P\times m):
                  P\times P\times P\longrightarrow P
\]
in $\Hot$, using its Cartesian product. The reversal identity
compares $T\circ T,\id_P:P\to P$.

For a literal counterexample to both implications, fix $n\geq2$
and use the split fibration $p:E=L_G(M,n)\to BG$ and
$0\ne\alpha\in H^n(G;M)$ in F.12. Translation by the section
$s_\alpha$ defines a fibrewise automorphism
$T=t_\alpha:E\to E$, with inverse $t_{-\alpha}$.
F.12 gives a coherent path $R_ns_\alpha\simeq R_ns_0$ over
$R_nBG$. The functor $R_n$ preserves the fibre products over
$BG$ used in fibrewise addition: at each test the base mapping
space is a $1$-type, and truncation at $n-1\geq1$ preserves
such pullbacks. Therefore $R_nT\simeq\id_{R_nE}$.
But $T^2s_0=s_{2\alpha}$ is not homotopic to $s_0$ by the
unpointed argument in F.12, since the rational vector space
$H^n(G;M)$ has $2\alpha\ne0$. Thus $T^2\ne\id_E$.
Now put $m=T\circ\operatorname{pr}_1:E\times E\to E$.
Its finite-test image is coherently homotopic to
$\operatorname{pr}_1$, an associative operation. Its actual
associators are $T^2\operatorname{pr}_1$ and
$T\operatorname{pr}_1$ on $E^3$. If they were homotopic,
fixing the other two coordinates and composing with $T^{-1}$
would give $T\simeq\id_E$, contrary to
$Ts_0=s_\alpha\not\simeq s_0$. These examples also give
height-one false positives by applying $\pi_0$ to the
height-two coherent comparison. They refute the proposed
general testing implications, without asserting a failure of
any particular geometrically defined path operation.

A coherent space-valued homotopy
is a path in the mapping space of presheaves between the two
induced natural transformations. F.7 or F.8 identifies that
path with a homotopy of the original maps, proving their
equality in $\Hot$. If the domain is countable CW, F.11
already gives the required injectivity on $\pi_0$ at height
two. An equivalence-detection theorem has a premise about
invertibility of a single morphism and a conclusion about its
invertibility; it cannot be substituted for this premise
about a path between two parallel morphisms.

For the applications at VII ms.~127, VII ms.~150, and IX
ms.~13, one can therefore use a geometrically constructed
associator or reversal homotopy directly, or retain the
coherent data required by the preceding tests. Merely checking
that each finite family of paths gives equal elements of a
set supplies no such coherent datum. In a strict Moore-path
model, concatenation is an actual associative operation
because lengths add and functions are joined on successive
intervals; reversal of a path twice is the original path.
Transferring those identities to a manuscript model still
requires the specified comparison equivalence and
compatibility of its composition and reversal maps. The
testing theorem neither constructs that comparison nor
makes its compatibility automatic.
\end{proof}

\partsection{G}{Finite categorical presentations and minimum sizes}
\label{sec:G}

We count identity arrows. A finite category has finitely many arrows, and hence finitely many objects. In this section ``a model of $Y$'' means a category $C$ together with an isomorphism $BC\simeq Y$ in $\Hot$; it does not mean a representative of the functor $X\mapsto[BX,Y]$ on $\Cathot$. A finite monoid is a one-object finite category. We distinguish minimization over finite categories from minimization over all small categories. These two object minima need not agree. We use $M(\Z/2,1)=S^1\cup_2D^2\simeq\mathbb{RP}^2$; an arbitrary space having only the same reduced homology would not specify this homotopy type.

\begin{claim}{G.1}{Proved}{Finite posets and necessary finiteness conditions}
A homotopy type is the classifying-space type of a finite poset if and only if it has a finite CW model. Every homotopy type with a finite CW model consequently has a finite categorical model. For any finite category $C$, $BC$ has finitely many cells in each dimension, finitely many components, finitely presented fundamental groups, and finitely generated integral homology in every degree. These latter conditions do not imply that $BC$ has a finite CW model.
\end{claim}
\begin{proof}
The nerve of a finite poset $P$ has a nondegenerate simplex for each strictly increasing chain. There are finitely many such chains, so $BP$ is a finite simplicial complex. Conversely, a finite CW complex is homotopy equivalent to a finite simplicial complex $L$; this is the finite version of simplicial approximation and mapping-cylinder replacement, recorded in \citet[Theorem 2C.5]{Hatcher2002}. The face poset $P(L)$ has one object for each nonempty simplex and a unique arrow $\sigma\to\tau$ exactly when $\sigma\subseteq\tau$. The simplices of its nerve are flags of faces, so $BP(L)$ is the barycentric subdivision of $|L|$. This is a specified homeomorphism followed by the chosen homotopy equivalence $|L|\to Y$. The empty poset represents the empty type, and disjoint unions treat finitely many components.

If $C$ has $m$ arrows, its nerve has at most $m^r$ simplices in degree $r\geq1$. The normalized cellular chains in every degree are therefore finitely generated free abelian groups. The one-skeleton has finitely many edges and the two-skeleton has finitely many triangles; choosing a maximal tree in each component gives a finite presentation of its fundamental group. These arguments also prove the homology assertion. They do not bound the dimension. For example the one-object category $C_2$ with multiplication modulo two has $BC_2=\mathbb{RP}^{\infty}$, whose mod-two homology is nonzero in every degree. Its complete two-element multiplication table occurs in the exhaustive computation below. No finite CW complex has this homology.
\end{proof}

\begin{claim}{G.2}{Refuted}{Finite type is not sufficient for finite categorical realization}
The assertion that every simply connected homotopy type with a simplicial model finite in every degree has a finite categorical model is false.
\end{claim}
\begin{proof}
List all finite categories up to isomorphism as $C_0,C_1,\ldots$; finite composition tables give such a countable list. Choose a base object in each nonempty $C_i$. Put $A_i=\Z/2$ if $\pi_{i+2}(BC_i)$ is zero and put $A_i=0$ otherwise; for the empty category use $A_i=\Z/2$. Set
\[
                 Y=\prod_{i\geq0}K(A_i,i+2).
\]
The standard simplicial abelian-group model for $K(A_i,i+2)$ is trivial below degree $i+2$ and finite in every degree. Consequently this product is finite in every degree. It is connected and simply connected, and its group in degree $i+2$ is exactly $A_i$. If $BC_i$ is nonempty, that group differs from the group of its chosen component by construction; a disconnected $BC_i$ already differs on components. The empty case is immediate. Thus no $BC_i$ is equivalent to $Y$. This diagonal construction is an existence proof, not an algorithm for recognizing the relevant homotopy groups.
\end{proof}

\begin{claim}{G.3}{Known}{Unrestricted monoids and countable models}
Every connected CW homotopy type is the classifying-space type of a small discrete monoid. For a connected countable CW homotopy type the monoid can be chosen countable.
\end{claim}
\begin{proof}
The unrestricted realization theorem is McDuff's theorem, with a primary-source alternative proof in \citet[Theorems 3.5 and 4.7]{Fiedorowicz1984}; we use those verified theorem statements. We give the countability reduction, since finiteness does not follow from that theorem.

Let $BM\simeq Y$ be such a realization and suppose $Y$ has a countable CW model. All homotopy groups of $Y$ are countable. Choose, for every degree, representatives of a generating set of $\pi_r(BM)$ and include all the finitely many arrow labels needed for each representative in a countable submonoid $M_0\subseteq M$. At stage $j$, the CW complex $BM_j$ is countable, so the kernels of $\pi_r(BM_j)\to\pi_r(BM)$ are countable for all $r$. For each element of each kernel choose a nullhomotopy in $BM$ and include its finitely many arrow labels in $M_{j+1}$. Each resulting submonoid is countable. The assertion about finitely many labels follows equally by simplicial approximation: a map or a homotopy with finite simplicial domain into a filtered union of nerves factors, after subdivision, through one member. The union $M_\omega$ is countable. Every class of a homotopy group of $BM_\omega$ is represented at a finite stage, and each kernel element is killed at the next stage at which it appears. Thus $BM_\omega\to BM$ is a weak equivalence. Both spaces are CW complexes, so it represents an isomorphism in $\Hot$. For a noncontractible type this monoid is nonempty; in particular the object minimum among all small categories is one.
\end{proof}

\begin{claim}{G.4}{Proved}{A finite monoid cannot have a nonsimply connected finite-dimensional type}
Let $M$ be a finite monoid. Its universal group $U(M)=\pi_1(BM)$ is finite. If $BM$ is homotopy equivalent to a finite-dimensional CW complex, then $U(M)=1$. In particular neither $S^1$ nor $\mathbb{RP}^2$ is the classifying-space type of a finite monoid.
\end{claim}
\begin{proof}
Let $\eta:M\to U(M)$ be the universal homomorphism. For $a\in M$, finiteness gives $a^r=a^s$ for some $r<s$, so $\eta(a)^{s-r}=1$. Every inverse of a generator is therefore a nonnegative power of that generator. It follows that every word in the generators and their inverses equals $\eta(m)$ for some $m\in M$. Thus $\eta$ is surjective, and $U(M)$ is finite. The identification with $\pi_1(BM)$ follows directly from the edge-and-triangle presentation: the generators are the nonidentity elements and the relations are $[a][b]=[ab]$.

Choose an idempotent $e$ in the minimal nonempty two-sided ideal $I$ of $M$. Such an ideal exists by finiteness, and a sufficiently large power of any of its elements is idempotent. We verify that $H=eMe$, with identity $e$, is a group. For $x\in eMe$, minimality gives $MxM=I$, so $e=axb$ for some $a,b\in M$. Multiplying by $e$ at both ends yields
\[
                    e=(eae)x(ebe)
\]
in the finite monoid $eMe$. In any finite monoid a product is invertible only if each factor is invertible: a right inverse makes right multiplication injective, hence surjective, and gives a two-sided inverse; the dual argument treats a left inverse. Thus $x$ is invertible. The group $H$ is finite, and the homomorphism $H\to U(M)$ is surjective since $\eta(eme)=\eta(m)$.

There is a canonical morphism $BH\to BM$ in $\Hot$ even though the inclusion $H\subset M$ sends $1_H$ to $e$, not necessarily to $1_M$. To type this map, let $\operatorname{Kar}(M)$ be the category of idempotents of $M$ with
\[
       \Hom_{\operatorname{Kar}(M)}(f,g)=gMf.
\]
The full one-object subcategory at $1_M$ is $M$ and the full one-object subcategory at $e$ is $H$. Every object is a retract of $1_M$. The inclusion $j:M\to\operatorname{Kar}(M)$ is a weak equivalence: for each object $f$, the comma category $j/f$ is a retract of $j/1_M$, which has a terminal object. A retract of a contractible space is contractible, and Quillen's Theorem A applies. Hence the typed morphism is
\[
          BH\longrightarrow B\operatorname{Kar}(M)
                     \xleftarrow[\simeq]{Bj} BM.
\]
It induces the surjection $H\to U(M)$ on fundamental groups, after the basepoint path given by the retract arrows is chosen.

Suppose $h:BM\to Z$ is a homotopy equivalence with $Z$ a connected finite-dimensional CW complex. Miller's theorem states that every map from the classifying space of a torsion group to $Z$ induces the zero map on fundamental groups \citep[Theorem B, p.~40]{Miller1984}. Apply it to the composite $BH\to BM\xrightarrow{h}Z$. Its map on fundamental groups is both zero and surjective. Consequently $\pi_1(Z)=U(M)=1$. Notice that this step uses the precise fundamental-group theorem, with no assumption that $Z$ was simply connected.
\end{proof}

\begin{claim}{G.5}{Known}{The homotopy-finite part of finite-monoid realization}
Among connected homotopy types having a finite CW model, the types realized by finite monoids are exactly the simply connected ones. The realizing monoid in the simply connected case may be chosen aperiodic.
\end{claim}
\begin{proof}
Necessity is G.4. Sufficiency is \citet[Theorem C, proved in Section 5]{BenBassat2026}, whose complete primary-source preprint was opened during this audit. Its stronger assertion is realization by a finite aperiodic monoid. The stated construction realizes the homotopy cofiber of a map from a finite connected CW complex to the classifying space of a finite aperiodic monoid (Theorem B), then applies this to a finite CW model having one vertex and no one-cells. This is a cited result, not a priority claim or a claimed independent verification of every step of that recent preprint. The necessity argument in G.4 does not depend on it.
\end{proof}

\begin{claim}{G.6}{Refuted}{Finite fundamental group is not sufficient for finite monoids}
The assertion that every connected finite CW complex with finite fundamental group is the classifying space of a finite monoid is false.
\end{claim}
\begin{proof}
Take $\mathbb{RP}^2$, whose fundamental group is $\Z/2$. G.4 excludes any finite monoid realization. In particular this refutes Conjecture 7.1 in the verified version of \citet{BenBassat2026}; it does not refute that paper's simply connected Theorem C. The explicit finite categorical and poset models below verify that the obstruction concerns the one-object restriction.
\end{proof}

\begin{claim}{G.7}{Proved}{Exact sphere minima for posets}
For every integer $n\geq1$, a poset whose classifying space is $S^n$ has at least $2n+2$ objects and at least $2(n+1)^2$ arrows. Both bounds are attained simultaneously. Thus the object/arrow minima for $S^1,S^2,S^3$ are respectively $(4,8),(6,18),(8,32)$.
\end{claim}
\begin{proof}
We prove a slightly stronger combinatorial bound. A finite flag simplicial complex with nonzero reduced $n$th homology over a field has a full subcomplex $L$ with this property which is minimal under vertex deletion. For each vertex $v$ of $L$, the Mayer--Vietoris sequence for
\[
        L=(L-v)\cup\operatorname{st}(v),\qquad
        (L-v)\cap\operatorname{st}(v)=\operatorname{lk}(v)
\]
shows that $\widetilde H_{n-1}(\operatorname{lk}(v))\ne0$: the nonzero group $\widetilde H_n(L)$ injects into it because $\widetilde H_n(L-v)=0$ and the star is contractible. For $n=0$ a complex with nonzero reduced homology has at least two vertices. Induction therefore implies that every vertex of $L$ has at least $2n$ neighbors. No vertex is adjacent to all other vertices: in a flag complex this would make $L$ a cone. Consequently $L$ has at least $2n+2$ vertices. Counting vertex-edge incidences gives at least $n|\Ob L|$ edges. Hence
\[
                   f_0(L)+f_1(L)\geq(n+1)f_0(L)
                                         \geq2(n+1)^2.
\]
The order complex of a poset is flag, and its edges are exactly the strict comparable pairs. Its objects plus its edges therefore count all categorical arrows. Passing to a full subcomplex only decreases these numbers. The same lower bound holds if the original poset is infinite, since any candidate below a finite upper bound is itself finite.

For the upper bound take $n+1$ ordered levels, each consisting of two incomparable elements, with every element of a lower level smaller than every element of a higher level. The order complex is the join of $n+1$ copies of $S^0$, hence is homeomorphic to $S^n$. It has $2(n+1)$ objects and
\[
             2(n+1)+4\binom{n+1}{2}=2(n+1)^2
\]
arrows. The supplied program verifies the complete categories and normalized homology for $n=1,2,3$. The object bound and uniqueness of the minimum model agree with \citet[Theorem 2.13 and Corollary 2.14]{BarmakMinian2006}; our argument also gives the global arrow bound, rather than assuming that a model with minimum objects must minimize arrows.
\end{proof}

\begin{claim}{G.8}{Known}{The projective-plane poset and its size bounds}
The minimum number of objects in a poset of type $\mathbb{RP}^2$ is $13$. Every such poset with $13$ objects has $49$ arrows. The global arrow minimum $a_{\mathrm{pos}}(\mathbb{RP}^2)$ satisfies
\[
                        28\leq a_{\mathrm{pos}}(\mathbb{RP}^2)\leq49.
\]
We do not identify that global minimum.
\end{claim}
\begin{proof}
The lower bound of thirteen objects and classification of the two thirteen-object models, opposite to one another, are \citet[Corollary 5.8 and Theorem 5.10, arXiv version~1]{CianciOttina2015}. Here the type being minimized is that of the order-complex realization, equivalently the weak homotopy type of the associated finite $T_0$ space.

For an explicit model, give $S^2$ the boundary triangulation of the octahedron, with vertices $(i,\epsilon)$ for $i\in\{0,1,2\}$ and $\epsilon\in\{0,1\}$. A face contains at most one vertex with a given first coordinate. The antipodal map switches every second coordinate. Its quotient is a regular CW decomposition of $\mathbb{RP}^2$ with three vertices, six edges and four triangular faces: no closed face contains an antipodal pair, so its quotient characteristic map remains injective. The face poset has thirteen objects. There are twelve vertex-edge incidences, twelve edge-face incidences and twelve vertex-face incidences; including identities gives $13+36=49$ arrows. Its order complex is the subdivision of this regular CW decomposition. This construction gives both the type and the complete incidence relation. The code forms the face orbits, closes incidence under transitivity, checks composition, and computes $H_1=\Z/2$ and $H_2=0$ by exact Smith forms. The cited classification shows that the opposite poset is the only other object-minimal case and has the same arrow count.

For the lower arrow bound choose an arrow-minimal model; the finite upper bound makes it finite. A vertex with exactly one comparable neighbor is a beat point, and deleting it retracts the order complex without changing its homotopy type while decreasing the arrow count. There are no isolated vertices in a connected nontrivial model. Thus the comparability graph of an arrow-minimal model has minimum degree at least two, giving $f_1\geq f_0$. If $f_0=13$ its arrow count is $49$. Otherwise $f_0\geq14$ and $|\Mor P|=f_0+f_1\geq28$. This proves precisely the displayed interval; it does not promote the thirteen-object classification to an arrow-minimum theorem.
\end{proof}

\begin{claim}{G.9}{Proved}{The five-element sphere monoid}
Let $P_2=\{1\}\amalg\{x_{ij}:i,j\in\{0,1\}\}$, with $1$ an identity and $x_{ij}x_{kl}=x_{il}$. Then $BP_2\simeq S^2$.
\end{claim}
\begin{proof}
The multiplication is associative since either bracketing of three nonunits keeps the first row and the last column. All elements are idempotent, so the universal group is trivial. Write $R=\Z[P_2]$ and use left modules. There is an augmented projective resolution
\[
0\longrightarrow Rx_{00}\oplus Rx_{01}
 \xrightarrow{d_2} R\xrightarrow{d_1}Rx_{00}
 \xrightarrow{\varepsilon}\Z\longrightarrow0,
\quad
 d_2(a,b)=a+b,\quad d_1(c)=c(x_{10}-x_{00}).
\]
Each $Rx_{ij}$ is projective, being the image of an idempotent right-multiplication endomorphism of $R$. The kernel of $\varepsilon$ is spanned by $x_{10}-x_{00}$. The map $d_1$ kills all four nonunits and sends $1$ to that difference, so its kernel is their free span. The two left ideals in the first term have the disjoint bases $\{x_{00},x_{10}\}$ and $\{x_{01},x_{11}\}$; hence $d_2$ maps isomorphically onto that kernel. Tensoring with the trivial right $R$-module gives
\[
                  0\longrightarrow\Z^2\xrightarrow{(1\ 1)}\Z
                     \xrightarrow{0}\Z\longrightarrow0.
\]
Thus $H_r(BP_2;\Z)$ is $\Z$ for $r=0,2$ and zero otherwise. Since $BP_2$ is simply connected, Hurewicz supplies a map $S^2\to BP_2$ inducing a homology generator, and the homological Whitehead theorem makes it a homotopy equivalence. This is the explicit monoid in \citet[unnumbered Lemma, pp.~33--34]{Fiedorowicz2002}. The program checks all products and the exact integral resolution; the argument is not based on a finite range of bar homology alone.
\end{proof}

\begin{claim}{G.10}{Proved}{The seven-element three-sphere monoid}
Let $P_3=P_2\amalg\{q,r\}$ with, in addition to the preceding rectangular-band products,
\[
 x_{ij}q=x_{ij}r=x_{i0},\qquad qx_{ij}=rx_{ij}=x_{ij},
 \qquad q^2=qr=q,\qquad r^2=rq=r.
\]
The identity remains $1$. Then $BP_3\simeq S^3$.
\end{claim}
\begin{proof}
Checking whether a factor belongs to the rectangular band or to $\{q,r\}$ proves associativity; all $7^3$ triples are also checked by the supplied program. All elements are idempotent, so $\pi_1(BP_3)=1$. Put $R=\Z[P_3]$. An augmented projective resolution of the trivial left module is
\[
0\to Rx_{01}\oplus Rq
 \xrightarrow{d_3}Rx_{00}\oplus R
 \xrightarrow{d_2}Rq
 \xrightarrow{d_1}Rx_{00}
 \xrightarrow{\varepsilon}\Z\to0,
\]
where
\[
 d_3(c,d)=(0,c+d),\qquad d_2(a,b)=a+b(r-q),\qquad
 d_1(c)=c(x_{10}-x_{00}).
\]
These formulas are $R$-linear maps with the indicated targets. The abelian bases of $Rx_{00}$, $Rq$ and $Rx_{01}$ are respectively
\[
 \{x_{00},x_{10}\},\qquad \{x_{00},x_{10},q,r\},\qquad
 \{x_{01},x_{11}\}.
\]
The kernel of $\varepsilon$ is generated by $x_{10}-x_{00}$. The map $d_1$ kills $x_{00},x_{10}$ and sends both $q,r$ to that difference; its kernel is the free span of $x_{00},x_{10},r-q$. Right multiplication by $r-q$ kills every nonunit and sends $1$ to $r-q$. Hence $d_2$ has that kernel as its image, and its own kernel is precisely $(0,I)$, where $I$ is the free span of the six nonunits. The two bases in $Rx_{01}\oplus Rq$ are disjoint and exhaust those six nonunits, so $d_3$ is injective with image $(0,I)$. Exactness in every degree follows.

After tensoring with the trivial right module the unaugmented complex is
\[
0\to\Z^2\xrightarrow{\left(\begin{smallmatrix}0&0\\1&1\end{smallmatrix}\right)}
\Z^2\xrightarrow{(1\ 0)}\Z\xrightarrow{0}\Z\to0.
\]
It has the homology of $S^3$. Simple connectivity, Hurewicz and the homological Whitehead theorem prove the claim. The construction is \citet[Example 8.2, version~1]{Sweeney2025v1}, whose six-element \emph{semigroup} becomes the seven-element monoid here after adjoining the identity. Omitting that extra arrow would give an incorrect monoid minimum. The exact integral calculation, including primitive images, is reproduced by the code.
\end{proof}

\begin{claim}{G.11}{Proved}{A three-object, nine-arrow projective plane}
There is a finite category $R_2$ with three objects and nine arrows such that $BR_2$ is homeomorphic to $\mathbb{RP}^2$.
\end{claim}
\begin{proof}
The objects are $0,1,2$. Each endomorphism set consists only of its identity. For $i<j$ put $\Hom(i,j)=\Z/2$, put $\Hom(j,i)=\varnothing$, and compose labels by addition modulo two. The only composition of two nonidentity arrows has type $0\to1\to2$, so associativity is immediate; the program nevertheless checks all typed triples.

Let $Q$ be the six-object poset with three levels, two elements $(i,\epsilon)$ at each level, and all comparisons between different levels. Define a functor $p:Q\to R_2$ by $p(i,\epsilon)=i$ and by labeling $(i,\epsilon)\to(j,\delta)$ with $\delta-\epsilon$. It preserves composition. The involution $(i,\epsilon)\mapsto(i,1-\epsilon)$ acts freely, and every composable string in $R_2$ has exactly two lifts to $Q$, once its first vertex is chosen. Therefore $NQ/(\Z/2)\cong NR_2$ degree by degree, with the displayed lift rule verifying all face and degeneracy maps. Realization gives the quotient of the octahedral $S^2$ by its antipodal action. Thus $BR_2\cong\mathbb{RP}^2$. There are three identities and two arrows for each of the three ordered object pairs, giving nine in total. This model is also the $n=2$ instance of \citet[Theorem 6.2]{CarcaciaCampos2026}; the preceding proof is independent of that preprint's minimality arguments.
\end{proof}

\begin{claim}{G.12}{Proved}{Removing strongly connected components only for fundamental groups}
Let $C$ be a connected finite category such that $BC$ has a finite-dimensional CW homotopy type. There exists a finite acyclic category $D$ with no more objects or arrows than $C$, and a functor $p:C\to D$ inducing an isomorphism on fundamental groups. We assert no equivalence on higher homotopy groups.
\end{claim}
\begin{proof}
Partition the objects by mutual reachability. A part $S$ is a full strongly connected subcategory. Fix $a\in S$ and arrows $u_x:a\to x$, $v_x:x\to a$ for every $x\in S$, with identities at $a$. In the groupoid completion, for $f:x\to y$ the corresponding loop at $a$ satisfies
\[
 [u_y]^{-1}[f][u_x]=[v_yu_y]^{-1}[v_yfu_x].
\]
Both terms on the right come from endomorphisms of $a$. Thus the universal group of the finite monoid $\operatorname{End}_C(a)$ surjects onto $\pi_1(BS,a)$. Its finite subgroup $e\operatorname{End}_C(a)e$, constructed in G.4, also surjects. The resulting map from that finite group's classifying space to $BC$ factors through $BS$ in $\Hot$. Miller's Theorem B implies that its image on $\pi_1(BC)$ is zero. Consequently every $\pi_1(BS)$ has trivial image in $\pi_1(BC)$.

Form a categorical quotient by identifying all objects of each $S$ and making all arrows internal to $S$ identities. One explicit definition uses the original arrows as generators and imposes these relations together with the original composition relations. Every new composable word has a representative which is one original arrow: insert an internal arrow between the target of one factor and the source of the next factor, then compose in $C$; that inserted arrow is the identity in the quotient. Thus the quotient has at most as many arrows as $C$. Its only endomorphisms are identities, since its object graph is the acyclic condensation of the mutual-reachability relation. Call it $D$, and write $p:C\to D$ for the quotient functor.

Taking groupoid completions in this presentation identifies $\pi_1(BD)$ with the quotient of $\pi_1(BC)$ by the normal closure of the images of all $\pi_1(BS)$. To see that no additional loop relation is hidden here, choose a tree connecting the objects inside each $S$ before choosing a global maximal tree; the internal tree edges account for object identification, while the remaining internal arrows account exactly for those subgroup images. All these images were just proved zero. Hence $p$ induces an isomorphism on fundamental groups. The finite-dimensional hypothesis was used for this last assertion through Miller's theorem, not for constructing the quotient.
\end{proof}

\begin{claim}{G.13}{Proved}{Exact finite-category minima for the projective plane}
Among finite categories of type $\mathbb{RP}^2$, the object minimum is three and the arrow minimum is nine.
\end{claim}
\begin{proof}
Apply G.12. An acyclic category with at most two objects has a nerve which is a graph, hence free fundamental group. Thus a finite category of projective-plane type has at least three objects.

Suppose a connected acyclic category has at most eight arrows. If it has at least four objects, connectivity forces at least one fewer nonidentity arrows than objects. It therefore has exactly four objects and at most four nonidentity arrows. Its one-skeleton is a connected multigraph of cycle rank at most one. Every nondegenerate nerve triangle has three distinct vertices and three distinct edges. Its boundary, if there is a cycle, is an embedded triangular cycle and represents a primitive generator of that rank-one graph group. Attaching the triangles therefore leaves a free group of rank zero or one. Higher simplices do not alter the fundamental group.

For three objects with no directed chain of length two the nerve is again a graph. Otherwise linearly order the objects as $0<1<2$ and write $a,b,c$ for the sizes of $\Hom(0,1)$, $\Hom(1,2)$ and $\Hom(0,2)$. All three numbers are positive. If $a=1$, the triangle relations express every $1\to2$ arrow as one $0\to2$ arrow after the unique $0\to1$ arrow is put in a maximal tree; the relation belonging to a chosen $1\to2$ tree arrow kills one $0\to2$ generator. The remaining presentation is free. The case $b=1$ is dual. If $c=1$, use one arrow of each adjacent Hom set for a tree. The triangle relations with one of these tree arrows first kill every remaining adjacent arrow, so the fundamental group is trivial. A nonfree fundamental group thus requires $a,b,c\geq2$, giving at least $3+6=9$ arrows. G.12 transfers this lower bound to all finite categories with finite-dimensional classifying-space type. The model of G.11 attains both bounds.
\end{proof}

\begin{claim}{G.14}{Proved}{The computed small-category lower bounds}
No finite category with at most four arrows has classifying-space type $S^2$. No connected finite category with at least two objects and at most six arrows has classifying-space type $S^3$.
\end{claim}
\begin{proof}
The attached program \texttt{code/g\_finite\_models.py} is an exhaustive certificate for these bounded statements. For $n$ objects and $m$ arrows it fixes identities $0,\ldots,n-1$, distributes the $m-n$ remaining arrows among the $n^2$ source-target pairs, and ranges over every possible value of each composable product. Identities determine all products involving an identity; the remaining products are accepted only when every associativity equation holds. Every finite category can be labeled in this form. Disconnected endpoint graphs are discarded, since the target types are connected. No quotient by isomorphism is used, so enumeration need not recognize isomorphism to be exhaustive.

For every accepted table the program constructs the normalized integral nerve boundary. It first computes the Betti vector through degree three over $\F_2$. This alone excludes all four-arrow candidates for $S^2$. For the second assertion only ten tables survive the $\F_2$ vector $(1,0,0,1)$: two with two objects and five arrows, and eight with two objects and six arrows. Exact rational ranks give $(1,0,0,0)$ in every surviving case. They consequently do not have the homology of $S^3$. In particular one must not identify the mod-two signature of $BC_2*BC_2$ in these low degrees with that of a three-sphere.

The monoid counts with a fixed labeled identity in orders one through four are $1,2,11,156$. The two-object counts for three through six arrows are $2,11,76,977$; the three-object counts for five and six arrows are $6,54$. The output contains the full surviving tables and the rational checks. The vanishing assertions depend only on homology in degrees at most three, so finite-range computation is sufficient for these exclusions.
\end{proof}

\begin{claim}{G.15}{Known}{Exact finite-category and finite-monoid sphere minima}
Among finite categories the object/arrow minima are $(2,4)$ for $S^1$, $(1,5)$ for $S^2$, and $(1,7)$ for $S^3$. Among finite monoids the arrow minima for $S^2,S^3$ are five and seven, respectively; $S^1$ has no finite monoid model.
\end{claim}
\begin{proof}
Two parallel arrows $0\rightrightarrows1$, together with identities, have a circle as their nerve realization. G.4 excludes a one-object finite model, and a connected category with fewer than four arrows either has one object or has a tree as its underlying graph; neither has type $S^1$. G.9 and G.14 establish the two-sphere claims.

For the three-sphere upper bound use G.10. The remaining lower-bound case after G.14 is a monoid of order at most six. We use the complete monoid enumeration in \citet[Table 5 and Theorem 2.7, version~2]{Sweeney2026v2}: none of the reduced monoids of orders at most six has $(H_1,H_2,H_3)=(0,0,\Z)$. The reduction is valid because, when an idempotent $e\ne1$ satisfies $eMe=eM$ or $eMe=Me$, its classifying space is homotopy equivalent to $B(eMe)$, a strictly smaller monoid. Iterating terminates. For clarity, the nontrivial reduced homology triples occurring up to order six are
\[
\begin{gathered}
(C_k,0,C_k)\ (k=2,3,4,5,6),\quad
(C_2^2,C_2,C_2^3),\quad(C_2,0,C_6),\\
(0,\Z,0),\quad(0,\Z,\Z),\quad(0,\Z,C_2),\quad(0,C_2,0),
\end{gathered}
\]
plus the trivial triple. All exclude $S^3$. This is an explicitly cited external exhaustive computation. We have reproduced the multi-object exclusion and the monoids through order four, but have not claimed to rerun all monoids of orders five and six. The distinction is retained in the results ledger.
\end{proof}

\begin{claim}{G.16}{Proved}{Size table and the unrestricted convention}
The following table records separate minima of objects and arrows among finite models. A dash means no finite model in that column's class; an interval records a lower and upper bound, not a known minimum.
\[
\begin{array}{c|cc|cc|cc}
Y&\multicolumn{2}{c|}{\text{finite categories}}&\multicolumn{2}{c|}{\text{posets}}&\multicolumn{2}{c}{\text{finite monoids}}\\
 &\Ob&\Mor&\Ob&\Mor&\Ob&\Mor\\ \hline
S^1&2&4&4&8&-&-\\
S^2&1&5&6&18&1&5\\
S^3&1&7&8&32&1&7\\
M(\Z/2,1)&3&9&13&[28,49]&-&-
\end{array}
\]
If arbitrary small categories and arbitrary small monoids are allowed, the object minimum becomes one for every displayed type. The least arrow cardinalities for arbitrary monoids are respectively $\aleph_0,5,7,\aleph_0$. The arrow minima for arbitrary categories and posets are the same as in the corresponding finite columns, with the same unresolved poset entry.
\end{claim}
\begin{proof}
All finite entries follow from G.4 and G.7--G.15. For the unrestricted object minimum, G.3 supplies a one-object category. The empty category represents none of these nonempty types. The countable reduction in G.3 provides a countable monoid of type $\mathbb{RP}^2$; $\Z$ itself supplies a countable monoid of type $S^1$. G.4 excludes finite monoids in these two cases, so the minimum arrow cardinal is exactly $\aleph_0$. A category with fewer arrows than a displayed finite upper bound has finitely many objects and arrows. Thus allowing infinite categories cannot lower any arrow minimum that already has a finite candidate. This also explains why the minimum objects and minimum arrows for unrestricted categories need not be attained by the same model.
\end{proof}

\begin{claim}{G.17}{Open}{The remaining realization and size questions}
We have not determined the exact integer $a_{\mathrm{pos}}(\mathbb{RP}^2)$ in $[28,49]$. We have not provided an intrinsic classification of all possibly infinite-dimensional homotopy types realized by finite categories or by finite monoids. In particular this audit does not settle whether $K(\Z,2)$ has a finite categorical model or a finite monoid model.
\end{claim}

The thirteen-point classification controls only the case of thirteen objects; a larger poset can in principle use fewer arrows, which is the precise missing step for the first question. The finite CW part of the second question is settled by G.1 and G.5, and G.2 shows why finite generation of homotopy or homology data cannot replace a classification in general. We searched the exact titles and citations of McDuff, Fiedorowicz, Barmak--Minian, Cianci--Ottina, Sweeney, Martinian--Steinberg and the two November/October 2026 preprints, together with the queries ``finite category K(Z,2)'', ``finite monoid CP infinity'', and ``finite monoid real projective plane''. We found no verified answer to these last two finite-model questions. This is an audit limitation, not an assertion that those questions are new or absent from the literature.

\clearpage
\bibliographystyle{apalike}
\bibliography{CATHOT_references}
\end{document}
